\documentclass[UTF8, a4paper, 12pt, fontset=fandol, scheme=chinese]{ctexart}
% ============================================================
%  SY02-17 《仿射 Grassmann 簇与几何 Satake 等价》译稿导言区
%  底本：Xinwen Zhu, An introduction to affine Grassmannians and the
%        geometric Satake equivalence, arXiv:1603.05593（PCMI 2015 讲义，95 页）
%  体例：定理类与公式共用一个计数器、按节编号（复现原文 Theorem 1.1、(1.9)）
% ============================================================

% ---------- 页面与字体 ----------
\usepackage{geometry}
\geometry{a4paper, left=2.6cm, right=2.6cm, top=2.8cm, bottom=2.6cm}
\linespread{1.22}
\setlength{\parskip}{0.25em}
\setlength{\emergencystretch}{1.5em}
\sloppy

% ---------- 数学 ----------
\usepackage{amsmath}
\usepackage{amsthm}
\usepackage{amssymb}
\usepackage{amsfonts}
\usepackage{mathtools}
\usepackage{mathrsfs}
\usepackage{stmaryrd}
\usepackage{bm}
\allowdisplaybreaks[2]

% ---------- 图表 ----------
\usepackage{graphicx}
\usepackage{xcolor}
\usepackage{tikz}
\usetikzlibrary{cd, arrows.meta, positioning, calc}
\usepackage[most]{tcolorbox}
\usepackage{enumitem}
\setlist{itemsep=2pt, topsep=3pt, parsep=0pt}
\setlist[enumerate,1]{label=(\arabic*)}

% ---------- 超链接 ----------
\usepackage[hyperfootnotes=false, colorlinks=true, linkcolor=blue!45!black, citecolor=green!35!black,
            urlcolor=teal!60!black, bookmarksnumbered=true]{hyperref}
\usepackage{xurl}% 允许 URL 在任意位置断行（译注中的 nLab 永久链接较长）
\usepackage{bookmark}

% ---------- 中文微调 ----------
\xeCJKsetup{CJKecglue={\hskip 0.24em plus 0.05em minus 0.05em}}
\ctexset{
  section    = {name={,}, aftername={.\enspace}, format=\Large\bfseries\centering},
  subsection = {aftername={.\enspace}, format=\large\bfseries},
  subsubsection = {aftername={.\enspace}, format=\normalsize\bfseries}
}

% ---------- 定理环境：全部挂在 equation 计数器上（原文即如此） ----------
% 原文 pcmi.cls 设 \numberwithin{equation}{subsection}，故公式为 (1.2.1)、
% 定理同计数器，也为「定理 1.1.3」；引言用 \section*，故编号形如 0.3.6。
\numberwithin{equation}{subsection}
\theoremstyle{plain}
\newtheorem{thm}[equation]{定理}
\newtheorem{prop}[equation]{命题}
\newtheorem{lem}[equation]{引理}
\newtheorem{sublem}[equation]{子引理}
\newtheorem{lem-def}[equation]{引理-定义}
\newtheorem{cor}[equation]{推论}
\newtheorem{conj}[equation]{猜想}
\theoremstyle{definition}
\newtheorem{dfn}[equation]{定义}
\newtheorem{construction}[equation]{构造}
\newtheorem{hypothesis}[equation]{假设}
\theoremstyle{remark}
\newtheorem{rmk}[equation]{注}
\newtheorem{ex}[equation]{例}
\newtheorem{exs}[equation]{练习}
\newtheorem{ques}[equation]{问题}

\renewcommand{\proofname}{证明}
\renewcommand{\qedsymbol}{$\square$}
\renewcommand{\refname}{参考文献}
\renewcommand{\contentsname}{目\quad 录}

% ---------- 译注标记（独立脚注序列） ----------
\usepackage{amscd}
\usepackage{manyfoot}
\DeclareNewFootnote[plain]{trans}
\renewcommand{\thefootnotetrans}{译\arabic{footnotetrans}}
\newcommand{\zy}[1]{\footnotetrans{#1}}
\newtcolorbox{yizhu}[1][]{
  enhanced, breakable, colback=orange!4!white, colframe=orange!55!black,
  boxrule=0.5pt, left=6pt, right=6pt, top=4pt, bottom=4pt,
  fonttitle=\bfseries\small, title=#1,
  before skip=6pt, after skip=6pt
}

\providecommand{\cprime}{'}
\providecommand{\vprime}{''}
\providecommand{\sprime}{'}

% ---------- 原文宏 ----------
% 由 工程文件/scripts/port_macros2.py 从原文导言区自动移植，勿手改。
\newcommand{\fraka}{{\mathfrak a}}
\newcommand{\frakb}{{\mathfrak b}}
\newcommand{\frakc}{{\mathfrak c}}
\newcommand{\frakd}{{\mathfrak d}}
\newcommand{\frake}{{\mathfrak e}}
\newcommand{\frakf}{{\mathfrak f}}
\newcommand{\frakg}{{\mathfrak g}}
\newcommand{\frakh}{{\mathfrak h}}
\newcommand{\fraki}{{\mathfrak i}}
\newcommand{\frakj}{{\mathfrak j}}
\newcommand{\frakk}{{\mathfrak k}}
\newcommand{\frakl}{{\mathfrak l}}
\newcommand{\frakm}{{\mathfrak m}}
\newcommand{\frakn}{{\mathfrak n}}
\newcommand{\frako}{{\mathfrak o}}
\newcommand{\frakp}{{\mathfrak p}}
\newcommand{\frakq}{{\mathfrak q}}
\newcommand{\frakr}{{\mathfrak r}}
\newcommand{\fraks}{{\mathfrak s}}
\newcommand{\frakt}{{\mathfrak t}}
\newcommand{\fraku}{{\mathfrak u}}
\newcommand{\frakv}{{\mathfrak v}}
\newcommand{\frakw}{{\mathfrak w}}
\newcommand{\frakx}{{\mathfrak x}}
\newcommand{\fraky}{{\mathfrak y}}
\newcommand{\frakz}{{\mathfrak z}}
\newcommand{\frakA}{{\mathfrak A}}
\newcommand{\frakB}{{\mathfrak B}}
\newcommand{\frakC}{{\mathfrak C}}
\newcommand{\frakD}{{\mathfrak D}}
\newcommand{\frakE}{{\mathfrak E}}
\newcommand{\frakF}{{\mathfrak F}}
\newcommand{\frakG}{{\mathfrak G}}
\newcommand{\frakH}{{\mathfrak H}}
\newcommand{\frakI}{{\mathfrak I}}
\newcommand{\frakJ}{{\mathfrak J}}
\newcommand{\frakK}{{\mathfrak K}}
\newcommand{\frakL}{{\mathfrak L}}
\newcommand{\frakM}{{\mathfrak M}}
\newcommand{\frakN}{{\mathfrak N}}
\newcommand{\frakO}{{\mathfrak O}}
\newcommand{\frakP}{{\mathfrak P}}
\newcommand{\frakQ}{{\mathfrak Q}}
\newcommand{\frakR}{{\mathfrak R}}
\newcommand{\frakS}{{\mathfrak S}}
\newcommand{\frakT}{{\mathfrak T}}
\newcommand{\frakU}{{\mathfrak U}}
\newcommand{\frakV}{{\mathfrak V}}
\newcommand{\frakW}{{\mathfrak W}}
\newcommand{\frakX}{{\mathfrak X}}
\newcommand{\frakY}{{\mathfrak Y}}
\newcommand{\frakZ}{{\mathfrak Z}}
\newcommand{\bA}{{\mathbb A}}
\newcommand{\bB}{{\mathbb B}}
\newcommand{\bC}{{\mathbb C}}
\newcommand{\bD}{{\mathbb D}}
\newcommand{\bE}{{\mathbb E}}
\newcommand{\bF}{{\mathbb F}}
\newcommand{\bG}{{\mathbb G}}
\newcommand{\bH}{{\mathbb H}}
\newcommand{\bI}{{\mathbb I}}
\newcommand{\bJ}{{\mathbb J}}
\newcommand{\bK}{{\mathbb K}}
\newcommand{\bL}{{\mathbb L}}
\newcommand{\bM}{{\mathbb M}}
\newcommand{\bN}{{\mathbb N}}
\newcommand{\bO}{{\mathbb O}}
\newcommand{\bP}{{\mathbb P}}
\newcommand{\bQ}{{\mathbb Q}}
\newcommand{\bR}{{\mathbb R}}
\newcommand{\bS}{{\mathbb S}}
\newcommand{\bT}{{\mathbb T}}
\newcommand{\bU}{{\mathbb U}}
\newcommand{\bV}{{\mathbb V}}
\newcommand{\bW}{{\mathbb W}}
\newcommand{\bX}{{\mathbb X}}
\newcommand{\bY}{{\mathbb Y}}
\newcommand{\bZ}{{\mathbb Z}}
\newcommand{\mA}{{\mathcal A}}
\newcommand{\mB}{{\mathcal B}}
\newcommand{\mC}{{\mathcal C}}
\newcommand{\mD}{{\mathcal D}}
\newcommand{\mE}{{\mathcal E}}
\newcommand{\mF}{{\mathcal F}}
\newcommand{\mG}{{\mathcal G}}
\newcommand{\mH}{{\mathcal H}}
\newcommand{\mI}{{\mathcal I}}
\newcommand{\mJ}{{\mathcal J}}
\newcommand{\mK}{{\mathcal K}}
\newcommand{\mL}{{\mathcal L}}
\newcommand{\mM}{{\mathcal M}}
\newcommand{\mN}{{\mathcal N}}
\newcommand{\mO}{{\mathcal O}}
\newcommand{\mP}{{\mathcal P}}
\newcommand{\mQ}{{\mathcal Q}}
\newcommand{\mR}{{\mathcal R}}
\newcommand{\mS}{{\mathcal S}}
\newcommand{\mT}{{\mathcal T}}
\newcommand{\mU}{{\mathcal U}}
\newcommand{\mV}{{\mathcal V}}
\newcommand{\mW}{{\mathcal W}}
\newcommand{\mX}{{\mathcal X}}
\newcommand{\mY}{{\mathcal Y}}
\newcommand{\mZ}{{\mathcal Z}}
\newcommand{\al}{{\alpha}}
\newcommand{\be}{{\beta}}
\newcommand{\ga}{{\gamma}}
\newcommand{\la}{{\lambda}}
\newcommand{\ve}{{\varepsilon}}
\newcommand{\Ga}{{\Gamma}}
\newcommand{\La}{{\Lambda}}
\newcommand{\ad}{{\on{ad}}}
\newcommand{\Ad}{{\on{Ad}}}
\newcommand{\aff}{{\on{aff}}}
\newcommand{\Aff}{{\on{Aff}}}
\newcommand{\Aut}{{\on{Aut}}}
\newcommand{\Bun}{{\on{Bun}}}
\newcommand{\cha}{{\on{char}}}
\newcommand{\der}{{\on{der}}}
\newcommand{\Der}{{\on{Der}}}
\newcommand{\diag}{{\on{diag}}}
\newcommand{\End}{{\on{End}}}
\newcommand{\Fl}{{\mF\ell}}
\newcommand{\Gal}{{\on{Gal}}}
\newcommand{\Gr}{{\on{Gr}}}
\newcommand{\Hom}{{\on{Hom}}}
\newcommand{\IC}{{\on{IC}}}
\newcommand{\id}{{\on{id}}}
\newcommand{\Id}{{\on{Id}}}
\newcommand{\ind}{{\on{ind}}}
\newcommand{\Ind}{{\on{Ind}}}
\newcommand{\inv}{{\on{Inv}}}
\newcommand{\Iso}{{\on{Isom}}}
\newcommand{\Lie}{{\on{Lie}}}
\newcommand{\Pic}{{\on{Pic}}}
\newcommand{\pr}{{\on{pr}}}
\newcommand{\Ql}{{\overline{\bQ}_\ell}}
\newcommand{\res}{{\on{res}}}
\newcommand{\Res}{{\on{Res}}}
\newcommand{\s}{{\on{sc}}}
\newcommand{\sgn}{{\on{sgn}}}
\newcommand{\Spec}{\on{Spec}}
\newcommand{\St}{{\on{St}}}
\newcommand{\tr}{{\on{tr}}}
\newcommand{\Tr}{{\on{Tr}}}
\newcommand{\Mod}{{\mathrm{-Mod}}}
\newcommand{\mmu}{\mu_\bullet}
\newcommand{\boxstar}{\makebox[0pt][l]{$\square$}\raisebox{.115ex}{\hspace{0.09em}\Large$\star$}\ }
\newcommand{\fstar}{\makebox[0pt][l]{$\square$}\raisebox{.115ex}{\hspace{0.09em}\Large$\star$}^f }
\newcommand{\ostar}{\text{$\star$}}  % 原文用 tikz 画圈，译文改用普通星号
\newcommand{\GL}{{\on{GL}}}
\newcommand{\GSp}{{\on{GSp}}}
\newcommand{\GU}{{\on{GU}}}
\newcommand{\SL}{{\on{SL}}}
\newcommand{\SU}{{\on{SU}}}
\newcommand{\SO}{{\on{SO}}}
\newcommand{\gl}{{\frakg\frakl}}
\newcommand{\Hk}{{\on{Hk}}}
\newcommand{\lhk}{\Hk^{\on{loc}}}
\newcommand{\pf}{{p^{-\infty}}}
\newcommand{\bb}{{\mathbf{b}}}
\newcommand{\bGr}{{\overline{\Gr}}}
\newcommand{\OGr}{\on{OGr}}
\newcommand{\wGr}{{\widetilde{\Gr}}}
\newcommand{\Sat}{{\on{Sat}}}
\newcommand{\Sp}{\on{Sp}}
\newcommand{\Ran}{\on{Ran}}
\newcommand{\Rat}{\on{Rat}}
\newcommand{\Rep}{\on{Rep}}
\newcommand{\Loc}{\on{Loc}}
\newcommand{\ppart}{(\!(t)\!)}
\newcommand{\pparu}{(\!(u)\!)}
\newcommand{\xcoch}{\mathbb{X}_\bullet}
\newcommand{\xch}{\mathbb{X}^\bullet}
\newcommand{\on}{\operatorname}


% ---------- 原文类文件 pcmi.cls 提供、但正文另需的宏 ----------
\providecommand{\scr}[1]{\mathscr{#1}}


\begin{document}
\pagestyle{plain}

% ============ 扉页 ============
\begin{titlepage}
\centering
\vspace*{2.2cm}

{\Huge\bfseries 仿射 Grassmann 簇\\[6pt] 与几何 Satake 等价}

\vspace{1.0cm}
{\LARGE\itshape An Introduction to Affine Grassmannians\\[4pt] and the Geometric Satake Equivalence}

\vspace{0.9cm}
{\large 中文译注本}

\vspace{1.4cm}
{\large 朱歆文（Xinwen Zhu）}

\vspace{0.4cm}
{\large PCMI 2015 暑期学校讲义（2017 年扩充版）}

\vspace{1.6cm}
{\normalsize 原文：arXiv:1603.05593，95 页，2016 年}

\vspace{0.3cm}
{\normalsize 译者：数译（DSH agent）\quad 版本：v1.0\quad 日期：2026 年 11 月}

\vspace{1.2cm}
\rule{0.72\textwidth}{0.4pt}

\vfill
{\small 本译本仅供学术研读；原文版权归作者所有。}
\end{titlepage}

% ============ 译者说明 ============
\section*{译者说明}
\addcontentsline{toc}{section}{译者说明}

\textbf{原文性质。}这是朱歆文（Xinwen Zhu）在 2015 年 PCMI 暑期学校（Park City，专题「几何朗兰兹纲领与自守形式」）的系列讲义，
后经扩充收入 IAS/Park City Mathematics Series 第 24 卷（2017）。全书五讲加一个附录：仿射 Grassmann 簇的基本性质、
Schubert 簇与横截片、Beilinson–Drinfeld Grassmann 簇与分解结构、对 $G$-丛模空间的应用、几何 Satake 等价，
附录是层论与形式代数几何的补充。讲义以定义与证明为主线，是几何朗兰兹纲领中几何 Satake 这一技术核心的入门标准文献。

\textbf{编号约定。}原文把定理与公式挂在同一个计数器上（PCMI 体例），按节编号，例如「定理 1.1」与「(1.9)」并列出现。
译文照此：全部定理类环境与公式共用计数器、按节重置，编号与原文逐一对应。引理、命题、注、例等的编号也按出现次序顺排，
与原文一致。

\textbf{译注标记。}文中「译\emph{n}」是译者所加，用独立脚注序列，不占用作者的脚注编号；作者脚注保留原编号。
译注只写三类内容：易错预警、含糊之处的点明、背景直觉。每条译注给出依据或标明「译者按（上下文推断，无外部依据）」。

\textbf{术语约定。}核心术语在全书首次出现时保留英文原词，其后只用中文。affine Grassmannian 译「仿射 Grassmann 簇」，
perverse sheaf 译「反常层」，convolution 译「卷积」，fusion product 译「融合积」，Beilinson–Drinfeld Grassmannian 译
「Beilinson–Drinfeld Grassmann 簇」，uniformization 译「单值化」。完整对照见 \verb|资料卡/04-术语文献勘误.md|。

\newpage
\tableofcontents
\newpage

% ============ 正文 ============

% ch00  引言（原文 555–726 行）

\section*{引言}
\addcontentsline{toc}{section}{引言}

\subsection{一些动机}

设 $F$ 是完备离散赋值域，$\mO\subset F$ 是其整数环，$k$ 是剩余域。又设 $V=F^n$ 是 $n$ 维 $F$-向量空间。
$V$ 的一个\emph{格}（lattice）是指 $V$ 的有限生成 $\mO$-子模 $\La$，且满足 $\La\otimes F=V$。例如 $\La_0=\mO^n$ 就是 $V$ 中的一个格。
本讲义的主要目的，是给 $V$ 的格构成的集合（及其变形与推广）赋予无穷维代数簇的结构，即所谓的\emph{仿射 Grassmann 簇} $\Gr$；随后研究它的几何性质，并给出若干应用。

不过，我们先给出研究这一对象的一份不完整的动机清单。

第一个动机来自环路代数、环路群与 Kac--Moody 代数、Kac--Moody 群之间的关系。由于每个格都可以经 $V$ 的线性自同构搬到 $\La_0$，
作为集合，$\Gr$ 自然等同于 $\GL_n(F)/\GL_n(\mO)$。当 $F=k((t))$ 时，群 $\GL_n(F)$ 与 $k$ 上的仿射 Kac--Moody 群密切相关。
把 $\Gr$ 实现为该 Kac--Moody 群的（部分）旗簇，对仿射 Kac--Moody 代数的表示论很重要。两者的确切关系将在 \S\ref{Kac-Moody} 中简述。

第二个动机来自代数曲线上向量丛的模空间。设 $X$ 是 $k$ 上光滑投影代数曲线，$x\in X(k)$ 是一点。取 $x$ 附近的局部坐标 $t$，
把 $X$ 在 $x$ 处的完备局部环与 $k[[t]]$ 等同。于是可以把 $\Gr$ 等同于下述集合
\[
\Gr=\left\{\, (\mE,\beta) \ \middle|\ \mE \text{ 是 } X \text{ 上的向量丛},\ \beta:\mE|_{X-x}\simeq \mO_{X-x}^n \text{ 是平凡化} \,\right\}.
\]
于是有从 $\Gr$ 到 $X$ 上秩 $n$ 向量丛的同构类集合的映射。后者实际上是一个精细得多的代数几何对象的点集，即 $X$ 上秩 $n$ 向量丛的模叠 $\Bun_n$。
因此很自然希望把上述集合论层面的映射升级为代数几何对象之间的态射
\[
u_x:\Gr\to \Bun_n .
\]
这样的态射称为\emph{单值化映射}（uniformization map），它是研究 $\Bun_n$ 的某些几何性质的重要工具。关于这一想法的详细解释，见 \S\ref{one point uniformization}。

第三个动机来自 $\Gr$ 中一些（局部）闭子簇的研究，它们出现于朗兰兹纲领的各个方面。这类子簇包括仿射 Springer 纤维（与非 Archimedes 局部域上约化群的调和分析有关）、
仿射 Deligne--Lusztig 簇（与志村簇模 $p$ 的约化有关）以及 Kisin 簇（与 Galois 表示的形变空间有关）。这些簇的简介见 \S\ref{digression}。

第四个动机同样来自朗兰兹纲领。描述 $\GL_n(F)$ 上 $\GL_n(\mO)$-双不变函数环的 Satake 同构\zy{更确切地说，是 Langlands 对 Satake 同构的重新解释。}，
是朗兰兹对偶的出发点。事实证明，Satake 同构有一个极大的推广，称为\emph{几何 Satake 等价}。它是几何朗兰兹纲领的出发点，
近年来还在经典算术朗兰兹纲领中找到了一些出人意料的用途（例如见 \cite{Laf}）。但即便只是为了陈述几何 Satake 等价，我们也需要知道 $\Gr$ 上的代数几何结构。

\subsection{范围与内容}

\subsubsection{} 本讲义以 2015 年夏在 Park City 数学研究所的五次讲座为基础。它的定位是对本课题的一次严肃导引。
话虽如此，我们只处理（我们认为）本课题最基本也最重要的方面。本讲义当然没有触及许多其他重要课题，一部分是因为篇幅有限，一部分是因为作者能力不足。

另一方面，我们希望本讲义是对既有文献的补充，而不是取代它们。因此，对每个结果我们并不都给出详细证明，尤其是当文献中已有这样的证明时；
这种情况下我们只指出相关文献。反过来，有些众所周知且被广泛使用的结果，文献中似乎并没有（完整的）证明。遇到这种情况，我们会尽量补出细节。

本讲义包含少数此前未见于文献的结果，尽管它们很可能已为专家所熟知。例如 Beilinson--Drinfeld Grassmann 簇的相对 Picard 层的计算（\S\ref{line bundle on BD Grass}）。

\subsubsection{} 下面简述各节内容。

第 \ref{Lecture I} 节解释基本构造，包括仿射 Grassmann 簇的各种（模）解释。我们还将引入重要的行列式线丛，
并讨论复数域上的仿射 Grassmann 簇与基于环路群之间的关系。虽然我们主要在等特征的局部域上考虑群，
但在本节末尾会简述为 $p$-进群定义仿射 Grassmann 簇的近期进展。

第 \ref{Lecture II} 节研究仿射 Grassmann 簇的一些基本几何。我们将引入仿射 Grassmann 簇中的各种子簇，包括 Schubert 簇、对偶 Schubert「簇」、横截片，
以及一些源于朗兰兹纲领的重要子簇。对单连通单群 $G$，我们研究其仿射 Grassmann 簇的 Picard 群。
我们还简述仿射 Grassmann 簇与仿射 Kac--Moody 群的部分旗簇之间的关系。

把仿射 Grassmann 簇解释为曲线上丛的某种模空间，这一想法有一个极大的推广，即 Beilinson--Drinfeld Grassmann 簇。
第 \ref{Lecture III} 节专门讨论这些对象。Beilinson--Drinfeld Grassmann 簇的一个显著特征是\emph{分解性质}（factorisation property）。
我们较详细地解释这一性质，特别是用曲线的 Ran 空间来表述它。我们还讨论 Beilinson--Drinfeld Grassmann 簇上的线丛。

第 \ref{Lecture IV} 节讨论仿射 Grassmann 簇对曲线上丛模空间研究的一些应用。基本工具是单值化定理。这样的定理实际上有两类：
当 $G$ 半单时成立的单点单值化，以及 ad\'elique 单值化——它是 Weil 对曲线上 $G$-丛同构类的 ad\`ele 描述的代数几何版本。
作为应用，当 $G$ 单且单连通时，我们给出光滑投影曲线上 $G$-丛模空间的 Picard 群的描述，并把模空间上线丛的整体截面解释为共形块（conformal block）。

第 \ref{Lecture V} 节专门讨论几何 Satake 等价。这一理论自然分成两部分。第一部分是从仿射 Grassmann 簇上某类反常层构造一个 Tannaka 范畴。
关键要素是该范畴上的\emph{融合积}结构，它的构造依赖 Beilinson--Drinfeld Grassmann 簇的分解性质。
第二部分是研究 Tannaka 群。这里的主要工具是 Mirkovi\'c--Vilonen 关于半无限轨道几何的理论。
我们还解释如何把几何 Satake 等价推广到适合（几何）朗兰兹纲领应用的形式。本节末尾，我们说明如何从几何版本恢复经典 Satake 同构。

另外有一个附录，收集第 \ref{Lecture V} 节用到的 $\ell$-进层的若干背景知识。

\subsection{约定与记号}

下面固定全讲义通用的记号与约定。读者\emph{务必}先跳过本小节，只在需要时回来查阅。

\subsubsection{$k$-空间}\label{Intro:sp}

设 $k$ 是交换域。以 $\on{Aff}_k$ 记仿射 $k$-概形范畴，即交换 $k$-代数范畴 $k\on{-alg}$ 的反范畴。
\emph{预层}（prestack 叫\emph{预叠}）是指从 $\Aff_k$ 到集合范畴（相应地，群胚 $2$-范畴）的共变函子（相应地，$2$-函子）。
我们把集合看成一个离散群胚，因此有时把预层看成离散的预叠。

给 $\Aff_k$ 赋以 fpqc 拓扑，即 $R$ 的一个覆盖是忠实平坦映射 $R\to R'$。$k$-\emph{空间}是指关于该拓扑在 $\Aff_k$ 上的层。
具体地说，空间 $\mF$ 是共变函子 $k\on{-alg}\to \on{Set}$，满足
\begin{enumerate}
\item 保持有限乘积，即
\[
\mF\Big(\prod_i R_i\Big)\to \prod_i \mF(R_i)
\]
是同构；
\item 且当 $R\to R'$ 忠实平坦时，
\begin{equation}\label{fpqc sheaf}
\mF(R)\to \mF(R')\rightrightarrows\mF(R'\otimes_RR')
\end{equation}
是等化子。
\end{enumerate}
两个空间之间的态射即函子的自然变换。$k$-空间范畴记为 $\Sp_k$。众所周知，$\Sp_k$ 以 $k$-概形范畴 $\on{Sch}_k$ 为全子范畴。
回忆：预层映射 $f:\mF\to\mG$ 称为\emph{概形化的}（schematic），若对每个概形 $T$，纤维积 $\mF\times_{\mG}T$ 都是概形。叠的概念可以类似地定义。

设 $X$ 与 $Y$ 是两个预叠。纤维积 $X\times Y$ 总是指 $X\times_kY$。若 $Y=\Spec R$，有时也把它写成 $X_R$、$X\otimes R$ 或 $X\otimes_kR$。
当 $Y=\Spec A$、$Z=\Spec B$ 时，我们也把 $X\times_YZ$ 写成 $X\otimes_AB$。

设 $X$ 是 $k$ 上的预叠。定义 $X$ 上的正则函数环为
\begin{equation}\label{reg fun}
\Gamma(X,\mO_X)=\Hom(X, \bG_a),
\end{equation}
其中 $\Hom$ 在预叠范畴中计算。也就是说，$X$ 上的正则函数 $f$ 是这样一种相容的方式：给每个 $x\in X(R)$ 指定元素 $f_x\in R$，并满足自然的相容性条件。
类似地，定义 $X$ 上线丛的 Picard 群胚为
\begin{equation}\label{Pic line}
\Pic(X)=\Hom(X,\bB\bG_m),
\end{equation}
其中 $\bB\bG_m$ 是 $\bG_m$ 的分类叠（见下面 \S\ref{Intro:torsor}）。也就是说，$\Pic(X)$ 的对象（称为 $X$ 上的线丛）是这样一种相容的方式：
给每个 $x\in X(R)$ 指定 $\Spec R$ 上的一个线丛，并满足自然的相容性条件。

\subsubsection{归纳概形}\label{Intro:Indsch}

本讲义中有一类特殊的 $k$-空间扮演重要角色。

\begin{dfn}\label{ind-sch}
$k$ 上的\emph{归纳概形}（ind-scheme）是指可写成子空间的\emph{滤过}正向极限
\begin{equation}\label{presentation}
X=``\varinjlim''_{i\in I}X_i
\end{equation}
的空间，其中每个 $X_i$ 都可由概形表示，且对 $I$ 中每个箭头 $i\to j$，相应映射 $X_i\to X_j$ 都是闭浸入。
\end{dfn}

\begin{rmk}
文献中，上面定义的归纳概形有时称为\emph{严格}归纳概形，以区别于概形范畴中一般的归纳对象。后一概念在本讲义中不会出现。
\end{rmk}

\begin{dfn}
设 $X$ 是 $k$ 上的归纳概形。称 $X$ 为：(i) \emph{归纳仿射的}；(ii) \emph{在 $k$ 上归纳有限型的}；(iii) \emph{在 $k$ 上归纳 proper 的}；
(iv) \emph{在 $k$ 上归纳投影的}；(v) \emph{既约的}，若存在形如 \eqref{presentation} 的表示，使得每个 $X_i$ 都可取为
(i) 仿射的；(ii) 在 $k$ 上有限型的；(iii) 在 $k$ 上 proper 的；(iv) 在 $k$ 上投影的；(v) 既约的。
定义 $X_{\on{red}}:=``\varinjlim''_{i\in I}(X_i)_{\on{red}}$。
\end{dfn}

\begin{ex}
形式概形是这样一种归纳概形：$X_{\on{red}}$ 是概形。
\end{ex}

\subsubsection{主丛}\label{Intro:torsor}

设 $G$ 是 $k$-群，即 $k$-空间范畴中的群对象。预叠映射 $E\to X$ 称为 $G$-\emph{主丛}（$G$-torsor），若 $G$ 作用在 $E$ 上，并且：
\begin{enumerate}
\item 映射 $E\to X$ 是 $G$-等变的（其中 $G$ 平凡地作用在 $X$ 上）；
\item 对每个 $\Spec R\to X$，存在忠实平坦映射 $\Spec R'\to\Spec R$，使得 $\Spec R'\times_XE$ 与 $\Spec R'\times G$ 是 $G$-等变同构的。
\end{enumerate}
若 $G$ 是 $k$ 上仿射群概形，这等价于要求 $E\to X$ 概形化且忠实平坦，且自然作用映射 $G\times E\cong E\times_X E$ 是同构。
若 (2) 中的 $\Spec R'$ 可取成使 $\Spec R'\to\Spec R$ 为 fppf（相应地，\'etale、Zariski）覆盖，则称该 $G$-主丛是 fppf（相应地，\'etale、Zariski）局部平凡的。

设 $X$ 是带 $k$-群 $G$ 作用的 $k$-空间。以 $[X/G]$ 记 fpqc 叠，其 $R$-点由 $\Spec R$ 上的 $G$-主丛 $E$ 连同 $G$-等变映射 $E\to X$ 组成。
自然投影 $X\to[X/G]$ 便是 $G$-主丛。若 $X=\Spec k$，通常把 $[X/G]$ 记为 $\bB G$，称为 $G$ 的分类叠。若 $G$ 在 $X$ 上的作用自由，则 $[X/G]$ 是 $k$-空间。
设 $E\to X$ 是 $G$-主丛，$Y$ 是 $G$-空间（即带 $G$ 作用的 $k$-空间），定义\emph{扭曲积}为
\[
X\tilde\times Y:= E\times^G Y:=[E\times Y/G],
\]
其中 $G$ 对角地作用在 $E\times Y$ 上。有自然投影 $X\tilde\times Y\to X$，其纤维同构于 $Y$。若 $Y=V$ 是 $G$ 的线性表示，
有时也把 $E\times^GV$ 写成 $V_E$，并把它看成 $X$ 上的向量丛。

若 $G=\GL_n$ 是一般线性群，$V=k^n$ 是其标准表示，则 $(\mE\to X)\mapsto(V_\mE\to X)$ 给出 $X$ 上 $\GL_n$-主丛与 $X$ 上秩 $n$ 向量丛的群胚等价。
本讲义中将自由使用这一等价。

上述记号显然可以推广到某个基 $k$-概形 $S$ 上的（非常值）群空间 $\underline G$。以 $\mE^0$ 记\emph{那个}平凡 $\underline G$-主丛。

\subsubsection{圆盘}\label{Intro:disc}

记
\[
D=\Spec k[[t]],\quad D^*=\Spec k((t)),
\]
称之为\emph{那个}\emph{圆盘}与\emph{去心圆盘}，以强调坐标 $t$ 的选定。对 $k$-代数 $R$，记
\[
D_R:= D\hat\times\Spec R:=\Spec R[[t]], \quad D_R^*:= D^*\hat\times \Spec R:= \Spec R((t)),
\]
并把它们看成由 $\Spec R$ 参数化的圆盘族（相应地，去心圆盘族）。

除另有说明，所谓局部域 $F$，指 $k$ 上的域，它关于一个在 $k$ 上平凡的离散赋值完备，且剩余域为 $k$。以 $\mO\subset F$ 记其整数环。
于是选定素元 $t\in\mO$ 后，有 $F\cong k((t))$、$\mO\cong k[[t]]$。设
\[
\mO\hat\otimes R:=\varprojlim (\mO/t^n\otimes R), \quad F\hat\otimes R:=\mO\hat\otimes R[1/t],
\]
从而 $\mO\hat\otimes R\cong R[[t]]$、$F\hat\otimes R\cong R((t))$。记 $D_{F,R}=\Spec (\mO\hat\otimes R)$，$D^*_{F,R}=\Spec (F\hat\otimes R)$。

设 $X$ 是 $k$ 上光滑几何连通曲线。以 $|X|$ 记 $X$ 的闭点集。对 $x\in|X|$，记 $X^*=X-\{x\}$。
以 $\mO_x$ 记 $X$ 在 $x$ 处的完备局部环，$F_x$ 记 $\mO_x$ 的分式域。对 $k$-代数 $R$，为简单起见把 $D_{F_x,R}$、$D^*_{F_x,R}$ 记成 $D_{x,R}$、$D^*_{x,R}$。
若 $x:\Spec R\to X$ 是 $X$ 的 $R$-点，记 $\Ga_x\subset X_R$ 为 $x$ 的图像，$\hat{\Ga}_x$ 为 $\Ga_x$ 在 $X_R$ 中的形式完备化。

\subsubsection{群论数据}\label{Intro:group}

设 $G$ 是 $k$ 上连通分裂约化群。以 $G_\der$ 记 $G$ 的导群，$G_\ad$ 记 $G$ 的伴随形式，$G_\s$ 记 $G_\der$ 的单连通覆盖。

以 $T$ 记 $G$ 的抽象 Cartan。回忆它定义为 Borel 子群 $B\subset G$ 对其幂幺根群的商。可以证明 $T$ 在不计典范同构的意义下与 $B$ 的选取无关。
当需要把 $T$ 作为极大环面嵌入 $G$ 时，我们会明确说明。

以 $\xch=\xch(T)$ 记权格，即自由 Abel 群 $\Hom(T,\bG_m)$；以 $\xcoch=\xcoch(T)$ 记余权格，即 $\xch(T)$ 的对偶。
以 $\Phi\subset\xch$（相应地，$\Phi^\vee\subset\xcoch$）记根系（相应地，余根系）。四元组 $(\xch,\xcoch,\Phi,\Phi^\vee)$ 称为 $G$ 的根数据。
以 $W$ 记 Weyl 群。Borel 子群 $B\subset G$ 决定了支配余权构成的半群 $\xcoch(T)^+\subset\xcoch(T)$ 与正根子集 $\Phi^+\subset\Phi$，它们其实与 $B$ 的选取无关。于是有典范双射
\[
\xcoch(T)/W\cong \xcoch(T)^+.
\]
回忆 $\xcoch(T)$ 上有偏序 $\leq$：若 $\mu-\la$ 是单余根的非负整系数线性组合，则记 $\la\leq\mu$；若 $\la\leq\mu$ 但 $\la\neq\mu$，则记 $\la<\mu$。
该偏序在 $\xcoch(T)^+$ 上的限制有时称为 $\xcoch(T)^+$ 上的 Bruhat 序。通常以 $2\rho\in\xch(T)$ 记所有正根之和，并令 $\rho=\frac12(2\rho)\in\xch(T)\otimes\bQ$。

\subsection{致谢}

本讲义基于 2015 年夏在 Park City 数学研究所的讲座。我要感谢 PCMI 的组织者、听众与工作人员。
我还要感谢Weizhe Zheng就 \S\ref{ULA sheaf} 的有益讨论，以及 Michael Harris、Jiuzu Hong与 Joel Kamnitzer 的宝贵意见。

% ============================================================
%  ch01  第一讲：仿射 Grassmann 簇及其初步性质
%  原文 Lecture I（原文行 727–902）
% ============================================================

\section{第一讲：仿射 Grassmann 簇及其初步性质}\label{Lecture I}

我们先给出将要研究的对象。

\subsection{$\GL_n$ 的仿射 Grassmann 簇}\label{aff GLn}

先要说明什么叫「一族格」。所谓 $k((t))^n$ 中的一族格，是指 $k((t))^n$ 的一族有限生成投射 $R[[t]]$-子模，张量回 $R((t))$ 之后恰好给出整个 $R((t))^n$。

\begin{dfn}
设 $R$ 是 $k$-代数。所谓 $k((t))^n$ 中的一个 \textbf{$R$-族格}，是指 $R((t))^n$ 的一个有限生成投射 $R[[t]]$-子模 $\La$，使得
\[
\La\otimes_{R[[t]]}R((t))=R((t))^n .
\]
\end{dfn}

\begin{dfn}
（$\GL_n$ 的）\textbf{仿射 Grassmann 簇} $\Gr_{\GL_n}$ 是这样的预层：它给每个 $k$-代数 $R$ 配上 $k((t))^n$ 中全体 $R$-族格构成的集合。
\end{dfn}

为简便起见，本小节把 $\Gr_{\GL_n}$ 简记作 $\Gr$。整个理论都建立在下面这个定理之上。

\begin{thm}\label{rep}
仿射 Grassmann 簇 $\Gr$ 可由一个归纳射影概形表示。
\end{thm}

回忆一下：这尤其意味着 $\Gr$ 可以写成一列有限维射影概形的递增并（见 \S\ref{Intro:Indsch}）。

定理的证明略有技术性，但背后的想法却非常简单，所以我们先做一个非正式的说明（只在 $k$-点层面）。取标准格 $\La_0=k[[t]]^n$。以 $\Gr^{(N)}\subset\Gr$ 记分类那些夹在 $t^N\La_0\subset\La\subset t^{-N}\La_0$ 之间的格 $\La$ 的子空间，则 $\Gr$ 是这些 $\Gr^{(N)}$ 的递增并。于是只需证明 $\Gr^{(N)}$ 可由射影概形表示。

对格 $\La\in\Gr^{(N)}$，商 $\La/t^N\La_0$ 可以看作 $k^{2nN}\cong t^{-N}\La_0/t^N\La_0$ 的子空间，并且它在 $t$ 的作用下稳定。这样一来，$\Gr^{(N)}$ 就实现为通常 Grassmann 簇 $\Gr(2nN)$ 的闭子概形，而后者分类 $k^{2nN}$ 的有限维子空间；因此 $\Gr^{(N)}$ 是射影概形。下面把这个想法写成严格证明。

\begin{proof}
给定一族 $R$-格 $\La\subset R((t))^n$，总存在足够大的整数 $N$，使得
\begin{equation}\label{subfun}
t^NR[[t]]^n\subset \La\subset t^{-N}R[[t]]^n .
\end{equation}
所以 $\Gr$ 是满足 \eqref{subfun} 的那些格所成子函子 $\Gr^{(N)}$ 的并。关键的观察是下面这个引理。

\begin{lem}\label{projectivity}
通过映射
\[
\La\mapsto Q=t^{-N}R[[t]]^n/\La ,
\]
$\Gr^{(N)}$ 与下述预层 $\Gr^{(N),f}$ 等同：
\[
\Gr^{(N),f}:=\left\{Q\ \middle|\ \text{$Q$ 是 $\dfrac{t^{-N}R[[t]]^n}{t^NR[[t]]^n}$ 的商模，且作为 $R$-模投射}\right\}.
\]
\end{lem}

后面会看到，这个论断对 $k$-点来说相当显然，对一般的 $R$ 却颇为微妙。我们还需要一个引理。

以 $\Gr(m)$ 记通常的 Grassmann 簇，它分类 $k^m$ 的有限维子空间，是各个 $\Gr(r,m)$ 的无交并，后者分类 $r$ 维子空间。

\begin{lem}\label{Embedding}
函子 $\Gr^{(N),f}$ 可由 $\Gr(2nN)$ 的一个闭子概形表示。
\end{lem}

先证引理 \ref{Embedding}。把 $t^{-N}k[[t]]^n/t^Nk[[t]]^n$ 与 $k^{2nN}$ 等同，则乘以 $t$ 给出该向量空间上的幂零算子 $\Phi$，而 $\Gr^{(N),f}$ 就是 $\Gr(2nN)$ 中分类那些在 $\Phi$ 下稳定的子空间的闭子概形。注意 $\Gr^{(N),f}$ 实际上是一些 Springer 纤维的并（\cite[\S~1]{Yu}）。

再证引理 \ref{projectivity}。首先要说明引理中的映射
\begin{equation}\label{I:map}
\Gr^{(N)}\to \Gr^{(N),f}
\end{equation}
确实有定义，即 $Q$ 作为 $R$-模投射。由于 $\Lambda/t\Lambda$ 是 $R$-投射的，
\[
R((t))^n/\Lambda\cong\bigoplus_{k\geq 0}t^{-k-1}\Lambda/t^{-k}\Lambda
\]
也 $R$-投射。于是由 $R$-模的短正合列
\[
0\to Q\to R((t))^n/\Lambda\to R((t))^n/t^{-N}R[[t]]^n\to 0
\]
即得 $Q$ 的投射性。

因此 \eqref{I:map} 把 $\Gr^{(N)}$ 实现为 $\Gr^{(N),f}$ 的子预层。还需说明 \eqref{I:map} 在每个 $R$ 的 $R$-点上满射。换言之，对每个 $k$-代数 $R$ 以及 $\Gr^{(N),f}$ 的一个由商映射 $t^{-N}R[[t]]^n/t^NR[[t]]^n\to Q$ 表示的 $R$-点，要说明 $R[[t]]$-模
\[
\Lambda:=\ker\bigl( t^{-N}R[[t]]^n\to  t^{-N}R[[t]]^n/t^NR[[t]]^n\to Q\bigr)
\]
是有限投射的。此时 $\Lambda\otimes R((t))\simeq R((t))^n$ 自动成立。

由引理 \ref{Embedding}，$\Gr^{(N),f}$ 局部有限型，故（见 \cite[Tag01ZC]{St}）
\[
\Gr^{(N),f}(\varinjlim R_i)=\varinjlim \Gr^{(N),f}(R_i).
\]
把 $R$ 写成有限生成子 $k$-代数的并，便可假设 $R$ 在 $k$ 上有限生成。

此时由于 $R$ 是 Noether 环，自然环同态 $R[t]\to R[[t]]$ 平坦（见 \cite[Tag00MB]{St}），并诱导同构 $t^{-N}R[t]^n/t^NR[t]^n\cong t^{-N}R[[t]]^n/t^NR[[t]]^n$。于是可定义 $R[t]$-模
\[
\Lambda_f:= \ker\bigl( t^{-N}R[t]^n\to  t^{-N}R[t]^n/t^NR[t]^n\to Q\bigr).
\]
由 $R[t]\to R[[t]]$ 的平坦性，$\Lambda=\Lambda_f\otimes_{R[t]}R[[t]]$。因此只需证明 $\Lambda_f$ 是有限投射 $R[t]$-模。显然 $\Lambda_f$ 有限生成，故只需说明它在 $R[t]$ 上平坦。因 $\La_f$ 已 $R$-平坦，由 \cite[Tag00MH]{St}，只需说明对每个 $x\in \Spec R$（剩余域记为 $K$），基变换 $\Lambda_f\otimes_{R[t]}K[t]$ 在 $K[t]$ 上平坦。由于 $Q$ 是 $R$-平坦的，
\[
\Lambda_f\otimes_{R[t]}K[t]= \ker \bigl( t^{-N}K[t]^n\to Q\otimes_R K\bigr),
\]
它无 $t$-挠，因而在 $K[t]$ 上平坦。引理 \ref{projectivity} 证毕。
\end{proof}

最后，不难看出含入 $\Gr^{(N)}\to \Gr^{(N+1)}$ 是闭嵌入。定理得证。

\begin{rmk}把 $R[[t]]$ 换成 $R[t]$ 这一想法，是 Beauville--Laszlo 用一条整体曲线实现 $\Gr$ 的特例。见下文定理 \ref{glob aff}。
\end{rmk}

\subsection{一般群的仿射 Grassmann 簇}

设 $G$ 是光滑仿射 $k$-群。回忆主丛的概念（见 \S\ref{Intro:torsor}）以及圆盘的 $R$-族（见 \S\ref{Intro:disc}）。$G$ 的\textbf{仿射 Grassmann 簇} $\Gr_{G}$ 定义为
\begin{equation}\label{loc aff}
\Gr_{G}(R)=\left\{(\mE,\beta)\ \left|\ \begin{array}{l}
\mE \text{ 是 } D_R \text{ 上的 } G\text{-主丛},\\
\beta:\mE|_{D^*_R}\cong  \mE^0|_{D^*_R} \text{ 是一个平凡化}
\end{array}\right.\right\}.
\end{equation}

注意这个定义对非常值群概形也说得通。确切地说，若 $\underline G$ 是 $k[[t]]$ 上光滑的仿射群概形，则可完全仿照 \eqref{loc aff} 定义 $\underline G$ 的仿射 Grassmann 簇 $\Gr_{\underline G}$，只须把 $G$ 换成 $\underline G$。

\begin{thm}\label{I:rep aff}
预层 $\Gr_{\underline G}$ 可由 $k$ 上一个局部有限型的归纳概形表示。若 $\underline G$ 约化，则 $\Gr_{\underline G}$ 是归纳射影的。
\end{thm}

\begin{rmk}
(i) 上述定义与定理中当然可以把 $k[[t]]$ 换成任一局部域的整数环。

(ii) 容易看出，对加法群 $G=\bG_a$，$\Gr_{\bG_a}$ 可由仿射空间的归纳极限表示，因此不是归纳射影的。
\end{rmk}

\begin{rmk}在本讲义中，凡是不需要额外代价就能对一般群概形叙述的结果，我们都按一般情形叙述。不熟悉群概形的读者尽可把 $\underline G$ 换成 $k$-群 $G$。
\end{rmk}

定理的证明基于以下两个结果。

\begin{prop}\label{lin rep}
存在忠实的 $k[[t]]$-线性表示 $\rho: \underline G\to \GL_n$，使得 $\GL_n/\underline G$ 拟仿射。此外，若 $\underline G$ 约化，则可选取 $\rho$ 使 $\GL_n/\underline G$ 仿射。
\end{prop}

\begin{proof}第一句见 \cite[Proposition 1.3]{PR}，第二句见 \cite[Corollary 9.7.7]{Ap}。
\end{proof}

\begin{prop}\label{gen G}
设 $\rho:\underline G\to \GL_n$ 是线性表示，且 $\GL_n/\underline G$ 拟仿射，则诱导映射 $f_\rho:\Gr_{\underline G}\to \Gr_{\GL_n}$ 是局部闭嵌入。此外，若 $\GL_n/\underline G$ 仿射，则它是闭嵌入。
\end{prop}

这个证明与后文无关，愿意接受该命题的读者可以跳过。

\begin{proof}
设 $\Spec R\to \Gr_{\GL_n}$ 是一映射，由 $(\mE,\beta)$ 表示。要说明基变换映射
\[
\mF:=\Spec R\times_{\Gr_{\GL_n}}\Gr_{\underline G}\to \Spec R
\]
是局部闭嵌入，并且在 $\GL_n/\underline G$ 仿射时是闭嵌入。以 $\pi:\mE\to D_R$ 记投影，并把 $\beta$ 看作 $\pi$ 在 $D^*_R$ 上的截面。以
\[
\bar\pi:[\mE/\underline G]\to D_R
\]
记 fpqc 商。由于 $\GL_n/\underline G$（拟）仿射，由（拟）仿射态射的忠实平坦下降（见 \cite[Tag0244, Tag0246]{St}），存在 $D_R$ 上仿射且有限展示的概形 $W$，使得 $[\mE/\underline G]\subset W$ 是拟紧开子集。截面 $\beta$ 诱导 $\bar\pi$ 在 $D^*_R$ 上的截面，记作 $\bar\beta$。把 $\mE$ 约化为 $\underline G$-主丛 $\mE'$，等价于给出 $\bar\pi$ 的一个截面。于是 $\mF$ 就是 $\Spec R$ 上的预层：它给每个 $R\to R'$ 配上 $\bar\pi$ 在 $D_{R'}$ 上的截面 $\beta'$，并要求 $\beta'|_{D^*_{R'}}=\bar\beta|_{D^*_{R'}}$。我们有

\begin{lem}
设 $p:V\to D_R$ 是有限展示的仿射概形，$s$ 是 $p$ 在 $D_R^*$ 上的截面。则 $\Spec R$ 上那个预层——它给每个 $R\to R'$ 配上 $p$ 在 $D_{R'}$ 上的截面 $s'$ 且要求 $s'|_{D^*_{R'}}=s|_{D^*_{R'}}$——可由 $\Spec R$ 的闭子概形表示。
\end{lem}

\begin{proof}
把 $V\subset \bA^N_{D_R}$ 嵌入某个（很大的）仿射空间，成为闭子概形。用坐标把 $s$ 写成
\[
s=(s_1(t),\ldots,s_N(t)),\quad s_i(t)=\sum_j s_{ij}t^{j}\in R((t)).
\]
所论预层便由闭子概形 $\Spec A\subset \Spec R$ 表示，其中 $A$ 由 $\{s_{ij}=0,\ i=1,\ldots,N,\ j<0\}$ 生成的理想定义。
\end{proof}

把该引理用于 $V=W$、$s=\bar\beta$，得到闭子概形 $\Spec A\subset\Spec R$，以及一个恒同映射 $\beta_A: D_A\to W_{D_A}$，它在 $D_A^*$ 上的限制是 $\bar\beta|_{D^*_A}$。沿 $t=0$ 给出的闭嵌入 $\Spec A\to D_A$ 作基变换，得到截面 $s_0:\Spec A\to W\times_{D_R}\Spec A$。于是 $\mF$ 由 $\Spec A$ 的开子概形 $s_0^{-1}([\mE/\underline G]\times_{D_R}\Spec A)$ 表示。
\end{proof}

\begin{rmk}\label{general base}
本讲义在固定的基域 $k$ 上工作。但对某些应用（例如 \cite{PR2,PZ}），需要在基概形 $S=\Spec A$ 上工作。从 $A[[t]]$ 上仿射光滑群概形 $\underline G$ 出发，同样可用 \eqref{loc aff} 定义仿射 Grassmann 簇 $\Gr_{\underline G}$，它是 $S$ 上的预层。若相应的命题 \ref{lin rep} 在这种一般性下成立，则与上面相同的论证说明 $\Gr_{\underline G}$ 可表示为 $S$ 上有限展示的归纳概形。已知命题 \ref{lin rep} 在以下情形成立：(i) $\underline G$ 是 $S$ 上线性约化群概形的基变换；(ii) $S$ 是 Dedekind 整环（见 \cite[Appendix]{PZ}）。
\end{rmk}

有意思（也很重要）的非常值群概形，来自 Bruhat--Tits 理论中的抛孔群概形（综述见 \cite{Ti}）。我们不一般地解释它们是什么，只给出两个基本例子。

\begin{ex}\label{hyperspacial and Iwahori}
设 $G$ 是 $k$ 上连通约化群，$B\subset G$ 是 Borel 子群。
\begin{enumerate}
\item 群概形 $\underline G:=G\otimes_k \mO$ 是抛孔的，通常称为 $G$ 的\textbf{超特殊}抛孔。
\item 仍设 $\underline G$ 如上。有求值映射 $\on{ev}: \underline G(\mO)\to G(k)$，$g\mapsto (g\bmod t)$。则存在 $\mO$ 上光滑、纤维连通且仿射的群概形 $\underline G'$，连同同态 $\underline G'\to \underline G$，它诱导同构 $\underline G'\otimes_{\mO}F\cong \underline G\otimes_\mO F=G\otimes_k F$，并把 $\underline G'(\mO)$ 与 $\on{ev}^{-1}(B(k))$ 等同。通常称 $\underline G'$ 为 $G$ 的 \textbf{Iwahori 群概形}。
\end{enumerate}
更多讨论见例 \ref{hyperspacial and Iwahori2}。
\end{ex}

\begin{rmk}\label{terminology}
本讲义把 $\Gr_{\underline G}$ 称为 $\underline G$ 的仿射 Grassmann 簇。这与文献中的标准术语\emph{并不一致}。例如当 $\underline G$ 是如上 $G$ 的 Iwahori 群概形时，文献中通常把 $\Gr_{\underline G}$ 称为 $G$ 的\textbf{仿射旗簇}，并记作 $\Fl_{G}$。
\end{rmk}

下述命题推广了定理 \ref{I:rep aff} 的第二句。

\begin{prop}\label{indproper}
设 $\underline G$ 是 $\mO$ 上纤维连通光滑的仿射群概形，其一般纤维连通约化。则 $\Gr_{\underline G}$ 归纳射影，当且仅当 $\underline G$ 是抛孔群概形。
\end{prop}

这与经典结果类似：连通约化群 $G$ 的齐性空间射影，当且仅当它是 $G$ 的部分旗簇。该结果在本讲义别处不用，故略去证明。

下面引入卷积 Grassmann 簇 $\Gr_{\underline G}\tilde\times\Gr_{\underline G}$，它定义为
\begin{equation}\label{I:conv}
\Gr_{\underline G}\tilde\times\Gr_{\underline G}(R)=\left\{(\mE_1,\mE_2,\beta_1,\beta_2)\ \left|\ \begin{array}{l}
\mE_1,\mE_2 \text{ 是 } D_R \text{ 上的 } \underline G\text{-主丛},\\
\beta_1:\mE_1|_{D_R^*}\cong \mE^0|_{D_R^*},\ \beta_2:\mE_2|_{D_R^*}\cong  \mE_1|_{D_R^*}
\end{array}\right.\right\}.
\end{equation}
有映射
\begin{equation}\label{conv m0}
m:\Gr_{\underline G}\tilde\times\Gr_{\underline G}\to \Gr_{\underline G},\quad (\mE_1,\mE_2,\beta_1,\beta_2)\mapsto (\mE_2, \beta_1\beta_2),
\end{equation}
以及自然投影
\[
\pr_1:\Gr_{\underline G}\tilde\times\Gr_{\underline G}\to \Gr_{\underline G}, \quad (\mE_1,\mE_2,\beta_1,\beta_2)\mapsto (\mE_1,\beta_1),
\]
二者一起诱导同构
\begin{equation}\label{I: isom}
(\pr_1,m):\Gr_{\underline G}\tilde\times\Gr_{\underline G}\cong \Gr_{\underline G}\times \Gr_{\underline G}.
\end{equation}
特别地，卷积 Grassmann 簇可表示。映射 $m$ 通常称为\textbf{卷积映射}。

类似地有 $n$ 重卷积 Grassmann 簇 $\Gr_{\underline G}\tilde\times\cdots\tilde\times\Gr_{\underline G}$，它分类 $(\mE_i,\beta_i)$，$i=1,\ldots,n$，其中 $\mE_i$ 是 $D_R$ 上的 $\underline G$-主丛，$\beta_i: \mE_{i}|_{D_R^*}\simeq \mE_{i-1}|_{D_R^*}$ 是同构（这里 $\mE_0=\mE^0$ 平凡）。对每个 $j=1,\ldots,n-1$，记下 $(\mE_i,\beta_i)$，$i=0,\ldots,j$，便得到到 $j$ 重卷积 Grassmann 簇的自然投影；另有 $n$ 重卷积映射
\begin{equation}\label{conv m}
m:\Gr_{\underline G}\tilde\times\cdots\tilde\times\Gr_{\underline G}\to \Gr_{\underline G},
\end{equation}
它把 $(\mE_i,\beta_i)$ 送到 $(\mE_n, \beta_1\cdots\beta_n)$。

% ============================================================
%  ch02  第一讲 §1.3：附着于去心圆盘的群
%  原文行 909–1032
% ============================================================

\subsection{附着于去心圆盘的群}

$\Gr_{\underline G}$ 还有一种有用的解释。先引入射流群与环路群。

\begin{dfn}
设 $X$ 是 $\mO=k[[t]]$ 上的预层。$X$ 的 \textbf{$n$-射流空间} $L^nX$ 是这样一个预层：它给每个 $R$ 配上集合
\[
L^nX(R)=X(R[t]/t^n).
\]
$X$ 的\textbf{形式射流空间}（有时也称为\textbf{弧空间}或\textbf{正环路空间}）$L^+X$ 是这样一个预层：它给每个 $R$ 配上集合
\[
L^+X(R)=X(R[[t]]).
\]
设 $X$ 是 $F=k((t))$ 上的预层。$X$ 的\textbf{环路空间} $LX$ 是这样一个预层：它给每个 $R$ 配上集合
\[
LX(R)=X(R((t))).
\]
若 $X$ 是 $k$ 上的预层，为简单起见，在不引起混淆时把 $L^+(X\otimes_k\mO)$ 记作 $L^+X$，把 $L(X\otimes_k F)$ 记作 $LX$。
\end{dfn}

\begin{prop}~
\begin{enumerate}
\item 设 $X$ 是 $\mO=k[[t]]$ 上有限型概形。则 $L^nX$ 可由 $k$ 上有限型概形表示。此外
\[
L^+X\cong \varprojlim\ L^nX
\]
可由概形表示。若 $X$ 仿射，则 $L^+X$ 也仿射。若 $f:X\to Y$ 是开（相应地，闭）嵌入，则 $L^+f: L^+X\to L^+Y$ 也是。

\item 设 $X$ 是 $F=k((t))$ 上有限型仿射概形。则 $LX$ 可由 $k$ 上归纳仿射概形表示。若 $f:X\to Y$ 是闭嵌入，则 $Lf:LX\to LY$ 也是。此外，若 $X=\underline X\otimes_\mO F$，其中 $\underline X$ 是 $\mO$ 上有限型仿射概形，则 $L^+\underline X\subset LX$ 是闭子概形。
\end{enumerate}
\end{prop}

\begin{proof}见 \cite[1.a]{PR}。
\end{proof}

\begin{ex}
若 $X=\bA^1$，则 $L^+X(R)=R[[t]]=\{\sum r_it^i \mid r_i\in R,\ i<0 \text{ 时 } r_i=0\}$，而 $LX(R)=R((t))=\{\sum r_it^i\mid  r_i\in R,\ i\ll 0 \text{ 时 } r_i=0\}$。因此 $L^+X\cong \Spec k[r_0,r_1,\ldots]$，而 $LX=\varinjlim_i \Spec k[r_{-i}, r_{-i+1},\ldots]$。
\end{ex}

\begin{ex}\label{hyperspacial and Iwahori2}
考虑例 \ref{hyperspacial and Iwahori} 中的群概形。自然同态 $\underline G'\to \underline G$ 诱导闭嵌入 $L^+\underline G'\subset L^+\underline G$。事实上，求值映射可以提升为仿射 $k$-群的映射 $\on{ev}:L^+G\to G$，而 $L^+\underline G'=\on{ev}^{-1}(B)$。但要注意，同态 $\underline G'\to \underline G$ 并不是闭嵌入，因为基变换到 $F$ 之后它是同构。
\end{ex}

回到仿射 Grassmann 簇。注意有自然作用
\begin{equation}\label{I:action}
L\underline G\times \Gr_{\underline G}\to \Gr_{\underline G},\quad (A, (\mE,\beta))\mapsto (\mE,A\beta).
\end{equation}

\begin{prop}\label{I:quotient}
仿射 Grassmann 簇 $\Gr_{\underline G}$ 可以等同于 fpqc 商 $[L\underline G/L^+\underline G]$。
\end{prop}

\begin{proof}
先需要一个引理。

\begin{lem}\label{et loc triv}
$D_R$ 上每个 $\underline G$-主丛都能在某个 \'etale 覆盖 $\Spec R'\to \Spec R$ 的 $D_{R'}$ 上平凡化。
\end{lem}

\begin{proof}
首先，映射 $\mE\otimes_{R[[t]]}R\to\Spec R$ 光滑（因为 $\underline G$ 光滑），故在某个 \'etale 覆盖 $\Spec R'\to \Spec R$ 上有截面。又由 $\mE\to D_R$ 的光滑性，该截面可提升为 $\mE$ 在 $D_{R'}$ 上的截面 $\epsilon$。
\end{proof}

现在来证命题。要证的论断等价于说 $L\underline G$ 是 $\Gr_{\underline G}$ 上的 $L^+\underline G$-主丛。确切地说，对 $k$-代数 $R$ 有典范同构
\[
L\underline G(R)\cong\left\{(\mE,\beta, \epsilon)\ \left|\ \begin{array}{l}
(\mE,\beta)\in \Gr_{\underline G}(R),\\
\epsilon: \mE^0\cong \mE \text{ 是平凡化}
\end{array}\right.\right\}.
\]
事实上，所要的同构由 $A\mapsto (\mE^0,A,\on{id})$ 给出，其逆映射为 $(\mE,\beta,\epsilon)\mapsto A:=\beta\epsilon$。
\end{proof}

\begin{rmk}
由命题的证明，$L^+\underline G$-主丛 $L\underline G\to \Gr_{\underline G}$ 在 \'etale 拓扑下局部平凡。若 $\underline{G}=G\otimes \mO$ 是常值约化的（更一般地，是所谓驯顺分歧的），它在 Zariski 拓扑下也局部平凡（见引理 \ref{big cell}）。
\end{rmk}

卷积 Grassmann 簇 \eqref{I:conv} 也有类似的解释。利用 $L^+\underline G$-主丛 $L\underline G\to \Gr_{\underline G}$ 以及 $L^+\underline G$ 在 $\Gr_{\underline G}$ 上的左作用（由 \eqref{I:action} 定义），可作扭曲积 $L\underline G\times^{L^+ \underline G}\Gr_{\underline G}$（见 \S\ref{Intro:torsor}）。同样的论证表明 $L\underline G\times^{L^+ \underline G}\Gr_{\underline G}$ 同构于 $\Gr_{\underline G}\tilde\times\Gr_{\underline G}$（特别地，我们的记号是相容的）。

由此得到一个直接推论。

\begin{cor}
仿射 Grassmann 簇 $\Gr_{\underline G}$ 形式光滑。
\end{cor}

\begin{proof}
事实上，$L\underline G$ 与 $L^+\underline G$ 都形式光滑。
\end{proof}

\begin{rmk}\label{formal smoothness}
对归纳概形而言，形式光滑是个很弱的概念。例如 $L\bG_m$ 形式光滑，但容易看出它\textbf{极不}既约。
\end{rmk}

下面是环路群的一个更深入的应用。设 $k$ 代数闭，$\underline G_F=\underline G\otimes F$ 连通约化。记 $I=\Gal(\overline F/F)$，它作用在 $\underline G_F$ 的基本群 $\pi_1(\underline G_F)$ 上。以 $\pi_1(\underline G_F)_I$ 记其余不变量群。以 $p$ 记 $k$ 的特征指数。

\begin{thm}\label{components}~
\begin{enumerate}
\item 有典范同构 $\pi_0(L\underline G)\cong \pi_1(\underline G_F)_I$。映射 $L\underline G\to \Gr_{\underline G}$ 诱导同构 $\pi_0(L\underline G)\cong \pi_0(\Gr_{\underline G})$。以 $(\Gr_{\underline G})^0$ 记单位连通分支。
\item 设 $\underline G_F$ 半单且单连通。则 $\Gr_{\underline G}$ 既约。
\item 设 $\underline G_F$ 半单且 $p\nmid \pi_1(\underline G_F)$。以 $G_{\on{sc}}$ 记 $\underline G_F$ 的单连通覆盖，并设它可延拓为 \'etale 覆盖 $\underline G_{\on{sc}}\to \underline G$。则它诱导同构 $\Gr_{\underline G_{\on{sc}}}\cong (\Gr_{\underline G})^0$。特别地，$\Gr_{\underline G}$ 既约。
\item 一般地，设 $p\nmid \pi_1(G_\der)$，其中 $G_\der$ 是 $\underline G_F$ 的导群。则映射 $\underline G_{\on{sc}}\to \underline G$ 诱导同构 $\Gr_{\underline G_{\on{sc}}}\cong (\Gr_{\underline G})^0_{\on{red}}$。
\end{enumerate}
\end{thm}

这个定理（在各种一般性下）的证明可从多处提取（\cite[\S 4.5]{BD}、\cite{BL}、\cite{LS}、\cite{BLS}、\cite{Fa}，最终可追溯到 \cite{PR}）。注意 (1) 类似于拓扑中众所周知的事实：环路空间的 $\pi_0$ 等于原空间的 $\pi_1$（当 $k=\bC$ 时这其实不止是类比，见定理 \ref{top triv}）。

由该定理，$\Gr_{\underline G}$ 上的许多问题都可化归到 $G$ 单连通的情形。

\begin{ex}
(i) 取 $G=\bG_m$。由定理，$\pi_0(\Gr_G)\cong \bZ$。事实上 $(\Gr_G)_{\on{red}}=\pi_0(G)=\bZ$。但如注 \ref{formal smoothness} 所述，它不既约。

(ii) 在 \cite[Remark 6.4]{PR} 中证明了：当 $G=\on{PGL}_2$、$\cha\ k=2$ 时，$\Gr_G$ 不既约。
\end{ex}

除 $L\underline G$ 在 $\Gr_{\underline G}$ 上的自然作用外，$\Gr_{\underline G}$ 上还有另一个重要的对称性——至少在 $\underline G=G\otimes_k\mO$ 为常值时如此。以 $\Aut(D)$ 记 $\mO=k[[t]]$ 的自同构群归纳概形。确切地说，它是 $\Aff_k$ 上的群预层，定义为
\[
\Aut(D)(R)=\Aut_R(R[[t]]).
\]
$R[[t]]$ 的任一 $R$-自同构 $\varphi$ 由 $\varphi(t)=a_0+a_1t+a_2t^2+\cdots\in R[[t]]$ 的像决定。容易验证：存在 $\psi$ 使 $\psi(\varphi(t))=t$，当且仅当 $a_0$ 幂零且 $a_1$ 可逆。因此
\[
\Aut(D)\cong \on{Spf} k[[a_0]] \times \Spec k[a_1^{\pm 1}]\times \Spec\ k[a_2,\cdots].
\]
注意 $\Aut^+(D)=\Spec k[a_1^{\pm 1},a_2,\cdots]$ 是闭正规子群概形，它分类那些模 $t$ 后为恒等映射的 $R[[t]]$ 自同构。于是
\begin{equation}\label{decomp1}
\Aut(D)\cong \Aut^+(D)\rtimes \hat{\bG}_a,
\end{equation}
其中 $\hat{\bG}_a=\on{Spf} k[[a_0]]$ 是形式加法群。此外，$\Aut^{++}(D)=\Spec\ k[a_2,\cdots]$ 是 $\Aut^+(D)$ 的超幂幺根。记 $\bG_m^{\on{rot}}=\Spec k[a_1^{\pm 1}]$，通常称为\textbf{旋转环面}。则
\begin{equation}\label{decomp2}
\Aut^+(D)=\Aut^{++}(D)\rtimes \bG_m^{\on{rot}}.
\end{equation}

\begin{rmk}
若 $F$ 是一般的局部域，且没有指定的素元可选，则不存在典范的半直积分解 \eqref{decomp1} 与 \eqref{decomp2}。但相应的群归纳概形短正合列仍然有定义。
\end{rmk}

现在，若 $\underline G$ 是 $k[[t]]$ 上的群概形，带有提升 $k[[t]]$ 上典范作用的 $\Aut(D)$-作用（例如常值的 $\underline G=G\otimes_k k[[t]]$），则有作用
\begin{equation}\label{auto of disc}
\Aut(D)\times \Gr_{\underline G}\to \Gr_{\underline G},\quad (g, (\mE,\beta))\mapsto (g^*\mE,g^*\beta).
\end{equation}
特别地，旋转环面 $\bG_m^{\on{rot}}$ 作用在 $\Gr_{\underline G}$ 上。

% ============================================================
%  ch03  第一讲 §1.4 Beauville--Laszlo 定理、§1.5 行列式线丛
%  原文行 1033–1191
% ============================================================

\subsection{Beauville--Laszlo 定理}

$\Gr_{\underline G}$ 还有一个重要的模空间解释，归功于 Beauville--Laszlo。

设 $X$ 是 $k$ 上既约连通曲线（不必光滑），$x\in |X|$ 是光滑闭点，$\underline G$ 是 $X$ 上光滑仿射群概形。定义预层
\begin{equation}\label{aff glob}
\Gr_{\underline G,x}(R)=\left\{(\mE,\beta)\ \left|\ \begin{array}{l}
\mE \text{ 是 } X_R \text{ 上的 } \underline G\text{-主丛},\\
\beta:\mE|_{X_R^*}\cong \mE^0|_{X_R^*} \text{ 是一个平凡化}
\end{array}\right.\right\}.
\end{equation}
令 $\underline G_x=\underline G\otimes_X \mO_x$。把 $X_R$ 上的 $\underline G$-主丛限制到 $D_{x,R}$，给出自然态射
\[
\on{res}:\Gr_{\underline G,x}\to \Gr_{\underline G_x}.
\]

\begin{thm}\label{glob aff}
上述态射是同构。
\end{thm}

同样，定理的证明有技术性，但想法很简单，先作非正式说明。若 $\mE$ 是 $D_x$ 上带平凡化 $\beta$（在 $D_x^*$ 上）的 $\underline G_x$-主丛，则可用 $\beta$ 把 $\mE$ 与 $X^*$ 上平凡的 $\underline G|_{X^*}$-主丛 $\mE^0$ 粘起来。这一构造给出 $\on{res}$ 的逆映射。

要把这个想法写成严格证明，先要讨论 Beauville--Laszlo 的下降定理（见 \cite{BL2}）。

\begin{thm}\label{descent lemma}
设 $p: A\to \tilde{A}$ 是交换环同态，$f\in A$。设 $f$（相应地，$p(f)$）在 $A$（相应地，$\tilde{A}$）中不是零因子，且诱导映射 $A/f\to \tilde{A}/p(f)$ 是同构。以 $\mM_f(A)$ 记无 $f$-挠的 $A$-模范畴，以 $\mM_{p(f)}(\tilde{A})$ 记 $(\tilde{A},p(f))$ 的类似范畴。以 $\mM$ 记三元组 $(M_1,M_2,\varphi)$ 的范畴，其中 $M_1$ 是 $A_f$-模，$M_2\in \mM_{p(f)}(\tilde{A})$，而 $\varphi$ 是 $M_1\otimes_{A_f}\tilde{A}_{p(f)}$ 与 $M_2\otimes_{\tilde{A}}\tilde{A}_{p(f)}$ 之间的同构。则
\begin{enumerate}
\item 函子 $M\mapsto M\otimes_A\tilde{A}$ 把 $\mM_f(A)$ 映入 $\mM_{p(f)}(\tilde{A})$，故有函子 $F : \mM_f(A)\to \mM$，它把 $M\in\mM_f(A)$ 送到 $(M_f, M\otimes_A\tilde{A}, \varphi)$，其中 $\varphi$ 是 $M_f\otimes_{A_f}\tilde{A}_{p(f)}$ 与 $(M\otimes_A\tilde{A})_{p(f)}$ 之间的自然同构。
\item $F : \mM_f(A)\to \mM$ 是范畴等价。
\item $M\in\mM_f(A)$ 是有限（相应地，平坦、有限投射）$A$-模，当且仅当 $M_f$ 与 $M\otimes_A\tilde{A}$ 都具有该性质。
\end{enumerate}
\end{thm}

在 \cite{BL2} 中取 $\tilde{A}=\hat A$，即 $A$ 关于 $f$ 的完备化；但同样的论证在此也适用。

\begin{rmk}
这个定理的内容是从覆盖 $(A_f, \tilde A)$ 把模下降到 $A$。它不能由通常的忠实平坦下降推出，因为没有假定 $A\to \tilde A$ 平坦。例如 $A$ 非 Noether 时，$A\to \hat A$ 通常不平坦。但即使在 Noether 情形，同构 $\varphi$ 也没有给出完整的下降数据，即不含 $\hat R\otimes_R\hat R$ 上的同构。
\end{rmk}

下面是该定理的推论。

\begin{cor}\label{global interpretation of loop}
环路群 $L\underline G_x$ 表示下述函子：它给每个 $k$-代数 $R$ 配上三元组 $(\mE,\al,\beta)$ 的集合，其中 $\mE$ 是 $X_R$ 上的 $\underline G$-主丛，$\al$ 是 $\mE|_{D_{x,R}}$ 的平凡化，$\beta$ 是 $\mE|_{X^*_R}$ 的平凡化。
\end{cor}

\begin{proof}
详细论证见 \cite[Lemma 6.1]{PZ}，思路如下。

可设 $X=\Spec A$ 仿射，$x$ 由 $t=0$ 定义。设 $R$ 是 $k$-代数，$g\in L\underline G_x(R)$。以 $H_1$（相应地，$H_2$）记 $\underline G\times_X D_{x,R}$（相应地，$\underline G\times_XX^*_R$）的正则函数环，它是平坦 $R[[t]]$-代数（相应地，$(A[t^{-1}]\otimes R)$-代数）。元素 $g$ 诱导代数同构 $\phi:H_1\otimes_{R[[t]]} R((t))\cong H_2 \otimes_{(A[t^{-1}]\otimes R)} R((t))$。

把定理 \ref{descent lemma} 用于三元组 $(H_1,H_2,\phi)$，得到平坦的 $(A\otimes R)$-代数 $H$。于是 $\mE=\Spec H$ 就是 $X_R$ 上的 $\underline G$-主丛，并带有典范平凡化 $\al$ 与 $\beta$。

反之，给定 $(\mE,\al,\beta)$，由 $g=(\beta|_{D_{x,R}^*})(\al|_{D_{x,R}^*})^{-1}$ 得到 $g$。
\end{proof}

现在比较命题 \ref{I:quotient} 与推论 \ref{global interpretation of loop}，定理 \ref{glob aff} 即得。也可以仿照上面的论证，直接构造 $\on{res}$ 的逆映射。细节留给读者。

\begin{rmk}
回忆 $\Aff_k$ 上 fpqc 层范畴与 $\on{Sch}_k$ 上 fpqc 层范畴等价。在此等价下，\eqref{aff glob} 中 $\Gr_{\underline G,x}$ 的模空间描述可自然推广到所有 $k$-概形。反之，把 \eqref{loc aff} 的模空间描述推广到所有 $k$-概形则很别扭：对非仿射概形 $S$，$S\hat{\times} \Spec k((t))$ 并不真的有意义。
\end{rmk}

\subsection{行列式线丛}\label{determinant line}

下一个主题是在 $\Gr_{\underline G}$ 上构造一个极丰沛线丛。回忆定理 \ref{rep} 证明中的交换图
\[\begin{CD}
\Gr^{(N)}@>>>\Gr^{(N+1)}\\
@VVV@VVV\\
\Gr(2nN)@>>>\Gr(2n(N+1))
\end{CD}\]
在通常的 Grassmann 簇 $\Gr(m)$ 上有行列式线丛 $\mL$：它把子空间 $L\subset k^m$ 配上商空间 $k^m/L$ 的最高外幂。这实际上是 $\Gr(m)$（各连通分支的）Picard 群的丰沛生成元。这些线丛在嵌入 $\Gr(m)\subset \Gr(m+1)$ 下相容，因此在 $\Gr_{\GL_n}$ 上定义了丰沛线丛 $\mL_{\det}$。我们更愿意给出一个更典范的构造。

先引入一个重要概念。

\begin{dfn}
若一个拓扑向量空间是某个离散向量空间的拓扑对偶（即通常的对偶，其 $0$ 的开邻域基由有限维子空间的零化子组成），则称它\textbf{线性紧}。若一个拓扑空间有一组由线性紧子空间组成的 $0$ 的邻域基，则称它\textbf{线性局部紧}，或\textbf{Tate 空间}。Tate 向量空间 $V$ 中的\textbf{格}是指 $V$ 的线性紧开子空间。
\end{dfn}

Tate 向量空间中任意两个格 $L_1, L_2$ 是\textbf{可公度的}，即商 $L_1/(L_1 \cap L_2)$ 与 $L_2/(L_1 \cap L_2)$ 都有限维。

\begin{rmk}\label{exact str}
Tate 向量空间构成范畴，其 Hom 为连续线性映射。这是 Quillen 意义下的正合范畴：容许单射为闭嵌入，容许满射为开满射。

Tate 向量空间还有不借助拓扑的范畴化途径。回忆对正合范畴 $\mC$，Beilinson 构造了「局部紧对象」范畴 $\underleftrightarrow{\lim} \mC$，作为 $\mC$ 的 pro-ind-对象范畴的满子范畴（\cite[Appendix]{Be}）。若把 $k$ 上有限维向量空间范畴 $\on{Vect}_k$ 看作正合范畴，则 Tate 向量空间范畴恰是 Beilinson 范畴 $\underleftrightarrow{\lim}\on{Vect}_k$。这一途径的好处是不用拓扑，并可推广到其他情形。
\end{rmk}

设 $V$ 是 $k$ 上 Tate 向量空间。对 $k$-代数 $R$，定义拓扑 $R$-模
\[V\hat\otimes_k R= \varprojlim_L (V/L)\otimes_kR,\]
其中反向极限取遍 $V$ 中的格。于是可定义群预层 $\GL(V)$，其 $R$-点为 $V\hat\otimes_k R$ 的连续 $R$-线性自同构。

\begin{dfn}
$V$ 中的一个 \textbf{$R$-族格}是指开有界子模 $\Lambda\subset V\hat\otimes_k R$，且 $V\hat\otimes_k R/\Lambda$ 是投射 $R$-模。
\end{dfn}

\begin{dfn}
\textbf{Sato Grassmann 簇} $\Gr(V)$ 是这样的预层：它给每个 $k$-代数 $R$ 配上 $V$ 中全体 $R$-族格构成的集合。
\end{dfn}

与定理 \ref{rep} 完全类似，
\begin{thm}
Sato Grassmann 簇 $\Gr(V)$ 可由归纳射影概形表示。
\end{thm}
事实上，用类似（而且更简单）的论证有自然同构
\[\Gr(V)\cong \varinjlim_{L_1\subset L_2}\Gr(L_2/L_1),\]
其中 $L_1\subset L_2$ 是两个格，故 $L_2/L_1$ 是有限维 $k$-向量空间，而 $\Gr(L_2/L_1)$ 表示分类 $L_2/L_1$ 的有限维子空间的通常 Grassmann 簇。

下面在 $\Gr(V)\times \Gr(V)$ 上引入行列式线丛（空间上线丛的定义见 \S\ref{Intro:sp}）。先回忆：有限投射 $R$-模 $M$ 的最高外幂记作 $\det(M)$，它是可逆 $R$-模。对 $\Lambda_1,\La_2\in\Gr(V)(R)$，定义 $(\La_1,\La_2)$ 的\textbf{相对行列式线}为
\[\det(\La_1|\La_2)=\det(L\hat\otimes R/\La_1)\otimes \det(L\hat\otimes R/\La_2)^{-1},\]
其中 $L$ 是 $V$ 中某个格，使得 $(L\hat\otimes R)/\La_1$ 与 $(L\hat\otimes R)/\La_2$ 都是投射 $R$-模。这在典范同构意义下与 $L$ 的选取无关。让 $R$ 变动，得到 $\Gr(V)\times\Gr(V)$ 上的线丛 $\mL_{\det}$。下面这个引理虽然初等，却概括了相对行列式线的重要性质。

\begin{lem}~
\begin{enumerate}
\item 对任意 $g\in  \GL(V)(R)$ 与 $\La_1,\La_2\in\Gr(V)(R)$，有典范同构 $\det(g\La_1|g\La_2)\cong\det(\La_1|\La_2)$，且对 $g,g'$，同构 $\det(gg'\La_1|gg'\La_2)\cong\det(g'\La_1|g'\La_2)\cong\det(\La_1|\La_2)$ 与 $\det(gg'\La_1|gg'\La_2)\cong\det(\La_1|\La_2)$ 一致。换言之，$\GL(V)$ 在 $\Gr(V)\times\Gr(V)$ 上的对角作用提升为 $\mL_{\det}$ 上的作用。
\item 对任意 $\La_1,\La_2,\La_3$，有典范同构 $\ga_{123}:\det(\La_1|\La_2)\otimes\det(\La_2|\La_3)\cong\det(\La_1|\La_3)$，使得对任意 $\La_1,\La_2,\La_3,\La_4$ 有 $\ga_{134}\ga_{123}=\ga_{124}\ga_{234}$。换言之，若以 $\mL_{\det}^\times$ 记 $\mL_{\det}$ 的全空间，则得到群胚
\[\mL_{\det}^\times\rightrightarrows\Gr(V).\]
\end{enumerate}
\end{lem}

\begin{rmk}
该群胚定义了 $\bG_m$-gerbe
\[\mD_V=[\Gr(V)/ \mL_{\det}^\times],\]
通常称为 $V$ 的\textbf{行列式 gerbe}（这个术语由 Kapranov 首先引入）。
\end{rmk}

\begin{rmk}
上述构造可以族化。确切地说，对交换环 $R$，有 Tate $R$-模范畴 $\on{Tate}_R$，它定义为 $\underleftrightarrow{\lim}\on{Proj}_R$ 的幂等完备化，其中 $\on{Proj}_R$ 是有限投射 $R$-模范畴（拓扑途径见 \cite{Dr}）。对 Tate $R$-模 $M$，有相应的 Sato Grassmann 簇 $\Gr(M)\to \Spec R$（见\emph{同上}），它可由归纳代数空间表示。此外，$\Spec R$ 上有行列式 gerbe $\mD_M$。
\end{rmk}

\begin{rmk}\label{graded line}
事实上，上述构造应当升级为分次版本。确切地说，对有限投射 $R$-模，可以把 $\det(M)$ 定义为分次可逆 $R$-模，其中 $\det(M)$ 的次数取为 $M$ 的秩（它是 $\Spec R$ 上局部常值函数）。此时上面定义的 $\det(\La_1\mid\La_2)$ 成为 $\Spec R$ 上的分次线丛。记住分次对行列式线上的分解结构至关重要（见注 \ref{graded factorization} 与注 \ref{det factorization2}），而后者对符号与互反律的应用又很重要（见 \cite{BBE, OZ1, OZ2}）。
\end{rmk}

取 $V=k((t))^n$，配以通常的 $t$-进拓扑。它是 Tate 向量空间，含格 $\Lambda_0=k[[t]]^n$。由引理 \ref{projectivity} 得典范闭嵌入 $i:\Gr_{\GL_{n}}\to\Gr(V)$。于是通过拉回得到线丛 $\mL_{\det}|_{\{\La_0\}\times\Gr_{\GL_n}}$，仍记作 $\mL_{\det}$。注意这是 $\Gr_{\GL_n}$ 上的极丰沛线丛。

一般地，设 $\rho: \underline G\to \GL_n$ 是命题 \ref{lin rep} 中忠实的 $\mO$-线性表示。则由命题 \ref{gen G}，它诱导局部闭嵌入 $f_\rho: \Gr_{\underline G}\to \Gr_{\GL_n}$，而 $\mL_{\det}$ 的拉回定义了 $\Gr_{\underline G}$ 上的极丰沛线丛。

\begin{rmk}\label{Pfaffians}
设 $\on{char} k\neq 2$（去掉这一假设见 \cite[\S\ 4.2.16]{BD}）。设 $V$ 是 $k$ 上 Tate 向量空间。$V$ 上的连续对称双线性型 $B$ 称为\textbf{非退化}的，若诱导映射 $V\to V^*$ 是同构，其中 $V^*$ 是 $V$ 的拓扑对偶。此时对格 $\Lambda$ 可定义其正交补 $\Lambda^\perp$。若 $\Lambda=\Lambda^\perp$，则称 $\Lambda$ 为\textbf{拉格朗日格}。正交 Sato Grassmann 簇 $\OGr(V,B)\subset \Gr(V)$ 是分类拉格朗日格的闭子空间。如 \cite[\S\ 4]{BD} 所构造，行列式线丛 $\mL_{\det}$ 在 $\OGr(V,B)\times \OGr(V,B)$ 上的限制有典范平方根 $\mL_{\on{Pf}}$，称为 \textbf{Pfaffian 线丛}。
\end{rmk}

% ============================================================
%  ch04  第一讲 §1.6 复数域上的仿射 Grassmann 簇、§1.7 $p$-进群
%  原文行 1192–1254
% ============================================================

\subsection{复数域上的仿射 Grassmann 簇}\label{top loop group}

本小节讨论复数域上的仿射 Grassmann 簇，故设 $k=\bC$。设 $\underline G=G\otimes_k \mO$，并把 $\Gr_{\underline G}$ 简记作 $\Gr$。此时 $\Gr$ 可以看作无穷维复解析空间（特别地，是 CW 复形）。

设 $K\subset G$ 是极大紧子群，$S^1$ 是单位圆。以 $\Omega K$ 记多项式映射 $(S^1,1)\to (K,e_K)$ 构成的空间。确切地说，用 $e^{i\phi}$（$\phi\in\bR$）参数化 $S^1$，并取定嵌入 $K\subset \on{SO}(n,\bR)$。则 $\Omega K$ 是这样的映射 $(S^1,1)\to (K,e_K)$ 构成的空间：与 $K\subset \on{SO}(n,\bR)$ 复合之后，映射的矩阵元都是 $e^{i\phi}$ 的 Laurent 多项式。把 $\Omega K$ 看作拓扑群。有下述定理。

\begin{thm}\label{top triv}~
\begin{enumerate}
\item 典范映射 $\Omega K\to LG$ 诱导同胚 $\Omega K\cong \Gr$。等价地说，有分解 $LG=\Omega K\cdot L^+G$。

\item 有同构 $\Omega K\times\Omega K\cong \Gr\tilde\times\Gr$，使下图交换
\[\begin{CD}
\Omega K\times \Omega K@>m>>\Omega K\\
@VVV@VVV\\
\Gr\tilde\times\Gr @>m>>\Gr,
\end{CD}\]
其中上行为环路群 $\Omega K$ 的通常乘法，下行为卷积映射 \eqref{conv m0}。
\end{enumerate}
\end{thm}

(1) 属于 Pressley--Segal（见 \cite[\S\ 8.3]{PS}）。严格说来，\emph{同上}只证明了 $G=\GL_n$ 的情形。但一般情形可由 Nadler 一个远为一般的定理推出（见 \cite[\S\ 4]{Na}）。确切地说，$\Omega K$ 可看作无穷维复群 $G(\bC[t,t^{-1}])$ 的一个「实形式」，Nadler 给出了 $G(\bC[t,t^{-1}])$ 的某些实形式在 $\Gr$ 上的轨道参数化；特殊到 $\Omega K$ 时，他的定理给出 (1)。(2) 由 (1) 直接推出。\zy{这里与拓扑环路群正式接上：nLab 词条 loop group（\url{https://ncatlab.org/nlab/show/loop+group}，2026-03-21 修订）记拓扑环路群的同胚分解 $\mathcal{L}G\simeq G\times\Omega G$；本文所得是更强的陈述 $\Gr\cong\Omega K$（A 级：\cite[\S 8.3]{PS}）。[来源：nLab loop group，2026-03-21；A 级 \cite{PS}。]}

\begin{rmk}
特别地，$\Gr$ 具有 $H$-空间的同伦型。此外，$L^+G$-主丛 $LG\to \Gr$ 平凡，且作为拓扑空间 $\Gr\tilde\times\Gr\cong \Gr\times\Gr$（不要把这个同构与 \eqref{I: isom} 混淆）。\S\ref{Lecture V} 将说明：对任意域 $k$ 上 $G$ 的仿射 Grassmann 簇，这些事实在上同调层面（事实上在 \'etale 同伦型层面）都有对应的体现。
\end{rmk}

\subsection{$p$-进群的仿射 Grassmann 簇}\label{p-adic}

简要讨论 $p$-进域 $F$（即 $\bQ_p$ 的有限扩张）上群的仿射 Grassmann 簇。与 $F=k((t))$ 的情形一样，第一步是让「由 $\bF_p$-代数 $R$ 的 $\Spec R$ 参数化的 $V=\bQ_p^n$ 的一族格」有意义；尤其要让圆盘 $D_R$ 有意义。虽然纤维积 $\Spec \bZ_p``\times''\Spec R$ 并无字面意义，但众所周知 $D_R$ 应当定义为
\[D_R:=\Spec W(R),\quad D^*_R:=\Spec W(R)[1/p],\]
其中 $W(R)$ 是 $R$ 的（$p$-典型）Witt 向量环。它的元素是序列 $(r_0,r_1,r_2,\ldots)\in R^{\bN}$，加法与乘法由某些（复杂的）多项式给出。到第一个分量的投影 $W(R)\to R$，$(r_0,r_1,\ldots)\mapsto r_0$，是满环同态，并带有乘法（但\emph{非加法}）截面 $R\to W(R)$，$r\mapsto [r]=(r,0,0,\ldots)$，称为 $r$ 的 \textbf{Teichm\"uller 提升}。若 $R$ 是完全环，$W(R)$ 中每个元素可唯一写成 $\sum_{i\geq 0} [a_i]p^i$，于是 $W(R)$ 可看作「以 $p$ 为变量、以 $R$ 为系数的幂级数环」。例如 $W(\bF_p)=\bZ_p$。一般介绍见 \cite[\S\ II.5]{Se2}。

一般地，设 $F$ 是 $p$-进域，$\mO$ 是其整数环，$k$ 是其剩余域。则可用 Drinfeld 引入的所谓分歧 Witt 向量类似地定义 $D_{F,R}$ 与 $D^*_{F,R}$。若 $R$ 完全，它们简单地由
\[D_{F,R}= \Spec W(R)\otimes_{W(k)}\mO,\quad D^*_{F,R}= \Spec W(R)\otimes_{W(k)}F\]
给出，其中 $W(k)\to \mO$ 是诱导剩余域恒等映射的唯一环同态。

现在设 $\underline G$ 是 $\mO$ 上光滑仿射群概形。则可完全仿照 \eqref{loc aff} 在 $\on{Aff}_k$ 上定义预层，但有一处需要注意：当 $R$ 不完全时，环 $W(R)$ 病态，例如 $R$ 不既约时 $p$ 可能是 $W(R)$ 的零因子。另一方面，注意

\begin{enumerate}
\item 若 $R$ 是完全 $k$-代数，则 $W(R)$ 性质良好，例如 $W(R)$ 中每个元素都有「幂级数展开」。
\item 概形 $X$ 在完全环上的 $R$-点决定了 $X$ 的完全化\footnote{完全 $k$-概形范畴是完全 $k$-代数上预层范畴的满子范畴，见 \cite[Lemma A.10]{Z14}。}。
\item 概形的（\'etale）拓扑（例如 $\ell$-进上同调）在过渡到完全化时不改变。
\end{enumerate}

因此，我们把 \eqref{loc aff} 中朴素的模空间问题限制到完全 $k$-代数范畴上。这样在该范畴上定义了预层；\'etale 拓扑与 fpqc 拓扑在该范畴上有意义，也可以在此框架下做代数几何（介绍见 \cite[Appendix A]{Z14}）。此时能问的最好问题是：这个函子是否可由完全 $k$-概形的归纳极限表示？下述困难定理（在 \cite{BS} 中证明）给出肯定回答。

\begin{thm}
函子 $\Gr_{\underline G}$ 可由归纳完全射影概形表示，称为 \textbf{Witt 向量仿射 Grassmann 簇}。
\end{thm}

想法是在 $\Gr_{\GL_n}$ 上构造丰沛的「行列式」线丛。困难在于：对 $R$-族格 $\La\subset W(R)^n$，商 $W(R)^n/\La$ \emph{不是} $R$-模，故 \S\ref{determinant line} 的构造不适用。此外，$\Gr_{\underline G}$ 上没有类似的 $\Aut(D)$ 对称性，$p$-进情形当然也没有定理 \ref{glob aff} 的类比。

此前在 \cite[\S~1]{Z14} 中证明了 $\Gr_{\underline G}$ 可由归纳完全真代数空间表示。那是个容易得多的定理，对大多数拓扑问题的应用（例如 \S\ref{Lecture V} 将讨论的几何 Satake 等价）已经够用。

% ============================================================
%  ch05  第二讲 §2.1：Schubert 簇
%  原文行 1255–1466
% ============================================================

\section{第二讲：仿射 Grassmann 簇的更多几何}\label{Lecture II}

本讲研究仿射 Grassmann 簇的一些基本几何。设 $F=k((t))$，$\mO=k[[t]]$。除特别说明外，设 $k$ 代数闭，$\underline G=G\otimes_k \mO$ 是常值约化群概形；在不引起混淆时把 $\Gr_{G}$ 简记作 $\Gr$。关于某些结论对非常值群的推广，见 \cite{PR} 与 \cite[\S\ 6.3]{Z11}；关于 $p$-进群的相应结论，见 \cite{Z14}。

\subsection{Schubert 簇}\label{II:Sch var}

沿用 \S\ref{Intro:group} 的记号。有典范双射
\begin{equation}\label{Cartan decomp2}
G(\mO)\backslash G(F)/G(\mO)\cong \xcoch(T)/W\cong \xcoch(T)^+,
\end{equation}
它可以具体实现如下。本小节余下部分固定嵌入 $T\subset B\subset G$。设 $\mu\in\xcoch(T)$ 是余权，它定义映射 $\mu:F^\times\to T(F)\subset G(F)$。以 $t^\mu=\mu(t)$ 记素元 $t$ 在 $G(F)$ 中的像，并把它看作 $LG$ 的 $k$-点；它依赖于 $t$ 的选取。它在 $\Gr$ 中的像仍记作 $t^\mu$，与素元的选取无关。于是有 Cartan 分解（见 \cite[\S\ 3.3.3]{Ti}）
\begin{equation}\label{Cartan decomp}
G(F)=\bigsqcup_{\mu\in\xcoch(T)^+} G(\mO)t^\mu G(\mO).
\end{equation}
注意双陪集 $G(\mO)t^\mu G(\mO)$ 既不依赖素元 $t$，也不依赖嵌入 $T\subset G$。因此 \eqref{Cartan decomp} 诱导 \eqref{Cartan decomp2}。

现在设 $\mE_1,\mE_2$ 是 $D$ 上两个 $G$-主丛，$\beta: \mE_1|_{D^*}\cong \mE_2|_{D^*}$ 是 $D^*$ 上的同构。选取同构 $\phi_1:\mE_1\cong \mE^0$ 与 $\phi_2:\mE_2\cong \mE^0$，便得到平凡 $G$-主丛的自同构 $\phi_2\beta\phi_1^{-1}\in \Aut(\mE^0|_{D^*})$，从而得到 $G(F)$ 中的一个元素。$\phi_1,\phi_2$ 的不同选取会让该元素左乘、右乘 $G(\mO)$ 中的元素。因此 $\beta$ 给出一个确定的元素
\[\inv(\beta)\in G(\mO)\backslash G(F)/G(\mO)\cong \xcoch(T)^+,\]
有时称为 $\beta$ 的\textbf{相对位置}。

\begin{rmk}\label{non-algebraically closed}
设 $K$ 是包含 $k$ 的域（不必代数闭），$\mE_1, \mE_2$ 是 $D_K$ 上两个 $G$-主丛，$\beta: \mE_1|_{D^*_K}\cong \mE_2|_{D^*_K}$ 是同构。此时仍可如下定义 $\inv(\beta)\in \xcoch(T)^+$。过渡到代数闭包 $\bar K/K$ 之后，上述构造给出
\[\inv(\beta_{\bar K})\in G(\bar K[[t]])\backslash G(\bar K((t)))/G(\bar K[[t]])\cong G(\mO)\backslash G(F)/G(\mO),\]
它与代数闭包 $\bar K$ 的选取无关。因此 $\inv(\beta)=\inv(\beta_{\bar K})$ 有定义。
\end{rmk}

现在设 $\mE_1,\mE_2$ 是 $D_R$ 上两个 $G$-主丛，$\beta: \mE_1|_{D_R^*}\cong \mE_2|_{D_R^*}$。对每个 $x\in \Spec R$，基变换给出
\[(\mE_1|_{D_{k(x)}},\mE_2|_{D_{k(x)}},\beta_{k(x)}: \mE_1|_{D_{k(x)}^*}\simeq \mE_2|_{D_{k(x)}^*}).\]
令 $\inv_x(\beta):=\inv(\beta_{k(x)})$。

\begin{prop}\label{HP}
设 $X=\Spec R$。给定 $\mu\in\xcoch(T)^+$，点集
\[X_{\leq \mu}:=\left\{x\in X\mid \inv_x(\beta)\leq \mu\right\}\]
在 $X$ 中 Zariski 闭。
\end{prop}

\begin{proof}
略述证明。设 $\rho_\chi: G\to \GL(V_\chi)$ 是 $G$ 的最高权 $\chi$ 的有限维最高权表示（例如取 $V_\chi$ 为 $G$ 的最高权 $\chi$ 的 Weyl 模），并令 $V_{\chi,\mE_i}=\mE_i\times^GV_{\chi}$，$i=1,2$，记相应的向量丛。则由定义，所论子集包含于集合
\[X_{V_\chi,\leq \mu}:=\left\{x\in X\mid\rho_\chi(\beta_x)(V_{\chi,\mE_1})\subset t^{-(\chi,\mu)}(V_{\chi,\mE_2})\right\},\]
而后者易见在 $X$ 中闭。事实上，在 $X$ 上局部（Zariski）地可设 $V_{\chi,\mE_i}$（$i=1,2$）都平凡（论证同引理 \ref{et loc triv}）。于是 $\rho_\chi(\beta)$ 由一个方阵给出，其元素为 $a_{ij}=\sum_h a_{ij,h}t^h\in R((t))$。此时 $X_{V_\chi,\leq \mu}$ 由方程 $a_{ij,h}=0$（$h<-(\chi,\mu)$）定义。

若 $G$ 的导群 $G_\der$ 单连通，则
\[X_{\leq \mu}=\bigcap_{V_\chi} X_{V_\chi,\leq \mu},\]
其中交取遍 $G$ 的所有有限维最高权表示。故 $X_{\leq \mu}$ 闭。

一般地，总能找到 $G$ 由环面 $D$ 作的中心扩张 $\tilde G$，使其导群单连通（即所谓的 $z$-扩张，见 \cite[Proposition 3.1]{MS}）。用引理 \ref{et loc triv}，在 $X$ 上 \'etale 局部地把 $\beta$ 提升为 $\tilde G$-主丛的映射 $\tilde \beta$，并把 $\mu$ 提升为 $\tilde G$ 的余权 $\tilde \mu$，便化归到前一种情形。
\end{proof}

定义（球）\textbf{Schubert 簇} $\Gr_{\leq \mu}$ 为 $\Gr_G$ 的闭子集
\[\Gr_{\leq \mu}=\left\{(\mE,\beta)\in \Gr\mid \inv(\beta)\leq \mu\right\},\]
并赋以既约概形结构。由命题 \ref{HP}，
\[\Gr_\mu:=\left\{(\mE,\beta)\in \Gr\mid \inv(\beta)=\mu\right\}=\Gr_{\leq \mu}\setminus \bigcup_{\la<\mu}\Gr_{\leq \la}\]
是 $\Gr_{\leq \mu}$ 的开子集，称为 \textbf{Schubert 胞腔}。

\begin{prop}\label{geom of Sch}~
\begin{enumerate}
\item $\Gr_{\mu}$ 是单个 $L^+G$-轨道，且是维数为 $(2\rho,\mu)$ 的光滑拟射影簇。
\item $\Gr_{\leq \mu}$ 是 $\Gr_\mu$ 的 Zariski 闭包，从而是射影簇。
\end{enumerate}
\end{prop}

\begin{proof}
对 (1)，注意 $t^\mu$ 在 $L^+G$ 作用下的稳定化子是 $L^+G\cap t^\mu L^+G t^{-\mu}$。于是诱导映射
\[L^+G/(L^+G\cap t^\mu L^+G t^{-\mu}) \to LG/L^+G,\quad g \mapsto gt^\mu\]
是局部闭嵌入。由 Cartan 分解 \eqref{Cartan decomp}，其像恰为 $\Gr_\mu$。因此 $\Gr_\mu$ 光滑，且 $\Gr_\mu$ 在 $t^\mu$ 处的切空间可等同于 $\frakg(\mO)/\frakg(\mO)\cap \Ad_{t^\mu}\frakg(\mO)$。但容易看出
\[\frakg(\mO)/\frakg(\mO)\cap \Ad_{t^\mu}\frakg(\mO)=\bigoplus_{(\al,\mu)\geq 0} \frakg_\al(\mO)/t^{(\al,\mu)}\frakg_\al(\mO),\]
其中和取遍 $G$ 的正根 $\al$，$\frakg_\al$ 是对应的根空间。于是简单计算得 $\dim \Gr_\mu=(2\rho,\mu)$。

对 (2)，注意若 $\la\leq \mu$，则存在正余根 $\al$，使得 $\mu-\al$ 支配且 $\la\leq \mu-\al\leq \mu$（例如 \cite[Lemma 2.3]{R3}）。因此只需说明 $t^{\mu-\al}$ 含于 $\Gr_\mu$ 的 Zariski 闭包。为此只需在 $\Gr_{\leq\mu}$ 中构造一条曲线 $C$，使 $t^{\mu-\al}\in C$ 且 $C-\{t^{\mu-\al}\}\subset \Gr_\mu$。这化归为下述 $\SL_2$ 计算。

首先，对每个整数 $m$，令 $t^{\la_m}:=\begin{pmatrix} t^m & 0\\ 0 &1\end{pmatrix}$，看作 $\on{PGL}_2(F)$ 中的元素，并令
\[K_m=\Ad_{t^{\la_m}}(L^+\SL_2)\subset L\SL_2.\]
注意 $K_m$ 作为 $\mO$-模的李代理由 $\{t^me, h, t^{-m}f\}$ 张成，其中 $\{e,h,f\}$ 是标准的 $\fraks\frakl_2$-三元组。又注意
\[\sigma_m:=\begin{pmatrix}0&-t^m\\t^{-m}&0\end{pmatrix}=\begin{pmatrix} t^m & 0\\ 0 &1\end{pmatrix}\begin{pmatrix}0&-1\\1&0\end{pmatrix}\begin{pmatrix} t^{-m} & 0\\ 0 &1\end{pmatrix}\in K_m.\]
以 $L^{>0}\SL_2$ 记求值映射 $\on{ev}:L^+\SL_2\to \SL_2$ 的核（见例 \ref{hyperspacial and Iwahori2}），并令 $K_m^{(1)}=\Ad_{t^{\la_m}}(L^{>0}\SL_2)$。则 $K_m/K_m^{(1)}\cong \SL_2$。

现在设 $i_\al: \SL_2\to G$ 是对应于 $\al$ 的典范同态。取 $m=(\mu,\al)-1$，考虑轨道
\begin{equation}\label{curve C}
C_{\mu,\al}:=Li_\al(K_{m})t^\mu.
\end{equation}
注意 $K_{m}^{(1)}\subset L^+G\cap t^\mu L^+G t^{-\mu}$，故 $C_{\mu,\al}$ 是 $K_{m}/K_m^{(1)}=\SL_2$ 下的齐性空间。容易看出 $C_{\mu,\al}\cong \bP^1$，且 $(L^+G\cap Li_\al(K_m))t^\mu\cong \bA^1\subset \bP^1$。此外 $\bP^1\setminus\bA^1=i_\al(\sigma_m)t^\mu$。由熟知等式
\[\begin{pmatrix}0&-1\\1&0\end{pmatrix}=\begin{pmatrix}1&-t^{-1}\\0&1\end{pmatrix}\begin{pmatrix}1&0\\ t&1\end{pmatrix}\begin{pmatrix}1&-t^{-1}\\0&1\end{pmatrix}\begin{pmatrix}t^{-1}&0\\0&t\end{pmatrix},\]
可见 $i_{\al}(\sigma_m)t^\mu=(t^{\mu-\al} \bmod L^+G)$。因此 $C_{\mu,\al}$ 就是所要的曲线。
\end{proof}

\begin{rmk}
为简单起见，设 $G_\der$ 单连通且 $\cha\ k=0$。此时可定义 $\Gr$ 的闭子概形 $\Gr'_{\leq \mu}$，其 $R$-点分类那些 $(\mE,\beta)$，使得对每个不可约最高权表示 $V_\chi$ 都有 $\beta(V_{\chi,\mE})\subset t^{-(\chi,\mu)}(V_{\chi}\otimes R[[t]])$。由命题 \ref{HP} 的证明知 $\Gr_{\leq \mu}$ 是 $\Gr'_{\leq \mu}$ 的既约子概形。但 $\Gr'_{\leq \mu}=\Gr_{\leq \mu}$ 是否成立仍是个开问题，与 \cite[Conjecture 2.14, Conjecture 2.20]{KWWY} 密切相关。
\end{rmk}

\begin{ex}\label{momega1}
设 $G=\GL_n$，$\mu=(r,0,\ldots,0)$（在标准等同 $\xcoch(T)\cong \bZ^n$ 下）。用 \S\ref{aff GLn} 中 $\Gr_{\GL_n}$ 的格表述，
\[\Gr_{\leq \mu}(R)=\left\{ \Lambda\subset \Lambda_0:=R[[t]]^n\mid \on{rk} \Lambda_0/\Lambda=r\right\}.\]
现设 $r=n$。以 $\mN_n$ 记 $n\times n$ 幂零矩阵构成的簇。有自然映射
\begin{equation}\label{nil to Gr}
\mN_n\to \Gr_{\leq \mu},\quad A\mapsto \La=(t-A)\La_0.
\end{equation}
它确实有定义，因为 $\det(t-A)=t^n$。引理 \ref{Lus obs} 将证明这是开嵌入，特别地 $\Gr_{\leq \mu}$ 给出 $\mN_n$ 的一个紧化。在拓扑空间层面，这最早由 Lusztig 注意到（见 \cite{Lu0}）。注意给出支配余权 $\la\leq \mu=(n,0,\ldots,0)$ 等价于给出 $n$ 的一个分拆，故每个 $\la\leq \mu$ 给出 $\mN_n$ 中的幂零轨道 $\mO_\la$。请读者验证 \eqref{nil to Gr} 把 $\mO_\la$ 送到 $\Gr_\la$。
\end{ex}

继续一般理论。有映射
\begin{equation}\label{parity map}
p:\xcoch(T)\to \bZ/2,\quad \mu\mapsto (-1)^{(2\rho,\mu)}
\end{equation}
它经过 $\xcoch(T)\to \pi_1(G)\to \bZ/2$ 分解，故由定理 \ref{components} 诱导映射 $p:\pi_0(\Gr_G)\to \bZ/2$。

\begin{cor}\label{parity}
Schubert 胞腔 $\Gr_{\mu}$ 位于偶（相应地，奇）分支，即 $p(\Gr_{\mu})=1$（相应地，$-1$），当且仅当 $\dim \Gr_{\mu}$ 为偶（相应地，奇）。
\end{cor}

后文中以 $(LG)_{\leq \mu}$（相应地，$(LG)_\mu$）记 $\Gr_{\leq \mu}$（相应地，$\Gr_\mu$）在 $LG$ 中的原像。注意 $(LG)_{\leq \mu}$ 与 $(LG)_{\mu}$ 都是概形。

对余权 $\mu$，以 $P_\mu$ 记 $G$ 中对应于 $\mu$ 的抛物子群，即由那些满足 $\langle\al,\mu\rangle\leq 0$ 的根 $\al$ 的根子群 $U_\al$ 生成的群。以 $g\mapsto \bar{g}:=\on{ev}(g)$ 记求值映射（见例 \ref{hyperspacial and Iwahori2}）。则有自然投影
\[p_\mu: \Gr_\mu\cong L^+G/(L^+G\cap t^\mu L^+G t^{-\mu})\to G/P_\mu,\quad (gt^\mu \bmod L^+G)\mapsto (\bar{g}\bmod P_\mu).\]
其纤维同构于仿射空间。另一方面，把 $G$ 看作 $L^+G$ 中「常值环路」子群，还有闭嵌入
\[i_\mu: G/P_\mu\cong G\cdot t^\mu\subset \Gr_\mu.\]
注意 $\Aut^+(D)$ 在 $\Gr_G$ 上的作用保持 $\Gr_\mu$\footnote{但 $\Aut(D)$ 在 $\Gr_G$ 上的作用并不保持 Schubert 簇：若 $g\in \Aut(D)$ 把 $t$ 送到 $a_0+t$（$a_0\in R$ 幂零），则 $gt^\mu\not\in \Gr_{\leq \mu}$。}，且在旋转环面 $\bG_m^{\on{rot}}$ 下的不动点子概形是 $G\cdot t^\mu$。

\begin{ex}
设 $G=\GL_n$，$\La_0=k[[t]]^n$，$\mu=(m_1,\ldots,m_n)\in \bZ^n$ 是支配余权（故 $m_i\geq m_{i+1}$）。则映射 $p_\mu: \Gr_\mu\to G/P_\mu$ 把格 $\Lambda\in \Gr_\mu$ 送到 $V=\La_0/t\La_0\cong k^n$ 上如下递减滤过：
\[\on{Fil}^iV=t^{-m_i}\La\cap \La_0/t^{-m_i}\La\cap t\La_0.\]
反之，给定 $V$ 上的递减滤过 $\on{Fil}^\bullet$，映射 $i_\mu$ 把它送到格
\[\Lambda = \sum\on{Fil}^iV\otimes  t^{m_i} \mO\subset V\otimes_kF\cong F^n.\]
\end{ex}

对较小的 $\mu$，下面给出 $\Gr_{\leq \mu}$ 更具体的描述。

回忆支配余权 $\mu$ 称为\textbf{极小}（minuscule）的，若 $\mu\neq 0$ 且对每个正根 $\al$ 有 $(\al,\mu)\leq 1$。注意极小余权在偏序 $\leq$ 下是 $\xcoch(T)^+$ 的极小元。因此

\begin{lem}\label{minusch}
若 $\mu$ 极小，则 $\Gr_{\leq\mu}=\Gr_\mu\cong G/P_\mu$。
\end{lem}

于是 $\Gr_G$ 中的极小 Schubert 簇恰是那些分类对应于极小余权的抛物子群的部分旗簇。例子包括通常的 Grassmann 簇 $\Gr(r,n)$、射影空间中的光滑二次超曲面等。

回忆支配余权 $\mu$ 称为\textbf{拟极小}的，若 $\mu\neq 0$、$\mu$ 不极小，且对每个正根 $\al$ 有 $(\al,\mu)\leq 2$。为简单起见，设 $G_\der$ 单。此时 $\mu$ 就是短支配余根，记作 $\theta$，它是 $\xcoch(T)^+-\{0\}$ 在偏序 $\leq$ 下的极小元。对应于 $\theta$ 的根是最高根 $\theta^\vee$。以 $\mL_\theta=G\times^{P_\theta} k_{\theta^\vee}$ 记 $G/P_\theta$ 上的极丰沛线丛，其中 $P_\theta$ 通过特征 $\theta^\vee$ 作用在一维空间 $k_{\theta^\vee}$ 上。注意 $0\in\xcoch(T)^+$ 是唯一严格小于 $\theta$ 的元素。因此 $\Gr_{\leq \theta}=\Gr_\theta\sqcup \Gr_0$。

\begin{lem}\label{quasiminusch}
设（为简单起见）$G_\der$ 单。若 $\mu=\theta$ 拟极小，则经由投影 $p_\theta$，$\Gr_\theta$ 同构于线丛 $\mL_\theta$ 的全空间，而 $\Gr_{\leq \theta}$ 同构于用 $\mL_\theta$ 给出的 $G/P_\theta$ 的射影嵌入的射影锥。在原点处爆破 $\widetilde{\Gr}_{\leq \theta}= \bP(\mL_{\theta}\oplus \mO)$ 给出奇点消解。
\end{lem}

\begin{proof}见 \cite[\S\ 7]{NP}。
\end{proof}

\begin{rmk}
事实上，$\widetilde{\Gr}_{\leq \theta}$ 是 $\Gr_{\leq \theta}$ 的一种 Demazure 消解，见 \cite[Lemma 2.12]{Z14}。
\end{rmk}

\begin{ex}
若 $G=\SL_2$，则 $\theta=\al$ 是正余根。此时 $\Gr_{\leq \theta}$ 同构于二次曲线 $\{(x,y,z)\mid x^2=yz\}\subset \bP^2$ 的射影锥。
\end{ex}

给定 $G$ 的 $\mu_1,\mu_2\in \xcoch^+$，可用 $L^+G$-主丛 $(LG)_{\leq \mu_1}\to \Gr_{\leq \mu_1}$ 定义 $\Gr_{\leq \mu_1}$ 与 $\Gr_{\leq \mu_2}$ 的扭曲积。也可以把它定义为
\[\Gr_{\leq \mu_1}\tilde{\times}\Gr_{\leq \mu_2}=\left\{(\mE_1,\mE_2,\beta_1,\beta_2)\in\Gr\tilde\times\Gr\mid \inv(\beta_1)\leq \mu_1,\ \inv(\beta_2)\leq \mu_2\right\},\]
它在 $\Gr\tilde\times\Gr$ 中闭，因而可表示。类似地，若 $\mmu=(\mu_1,\ldots,\mu_n)$ 是 $G$ 的支配余权序列，可定义
\[\Gr_{\mmu}:=\Gr_{\mu_1}\tilde\times\cdots\tilde\times\Gr_{\mu_n}\subset \Gr_{\leq\mmu}:=\Gr_{\leq \mu_1}\tilde\times\cdots\tilde\times\Gr_{\leq \mu_n}\subset\Gr\tilde\times\cdots\tilde\times\Gr.\]
令 $|\mmu|=\sum \mu_i$，则卷积映射 \eqref{conv m} 诱导
\begin{equation}\label{conv m2}
m: \Gr_{\leq \mmu}\to \Gr_{\leq |\mmu|}, \quad (\mE_\bullet,\beta_\bullet)\mapsto (\mE_n,\beta_1\cdots\beta_n).
\end{equation}
这个映射类似于 Demazure 消解（只是 $\Gr_{\leq \mmu}$ 一般不光滑）。如 \S\ref{Lecture V} 所示，这些映射的几何极为丰富。这里给出一个例子。

\begin{ex}\label{conv v.s. springer}
沿用例 \ref{momega1} 的记号。设 $\omega_1=(1,0,\ldots,0)$，有卷积映射 $m:\Gr_{\leq(\omega_1,\ldots,\omega_1)}\to\Gr_{\leq \mu}$。另一方面有经典的 Springer 消解 $p:\tilde \mN_r\to\mN_r$（见 \cite[\S~1]{Yu}）。这些映射纳入下述交换图
\begin{equation}\label{sat and springer}
\begin{CD}
\Gr_{\leq(\omega_1,\ldots,\omega_1)}@<<< \Gr_{\leq (\omega_1,\ldots,\omega_1)}^{\square}@>>>\tilde \mN_r\\
@VmVV@VVV@VVpV\\
\Gr_{\leq \mu}@<\pi<<\Gr_{\leq \mu}^{\square}@>\phi>> \mN_r,
\end{CD}
\end{equation}
其中 $\Gr_{\leq \mu}^{\square}(R)$ 分类三元组 $(\mE,\beta,\epsilon)$，这里 $(\mE,\beta)\in\Gr_{\leq \mu}(R)$，而 $\epsilon: \La_0/\La\cong R^r$ 是 $R$-模同构。注意图中两个方块都是笛卡尔方块。事实上，$\Gr_{\leq(\omega_1,\ldots,\omega_1)}$ 中位于 $\La\in \Gr_{\leq \mu}$ 上方的一个点给出链 $\La=\La_r\subset\La_{r-1}\subset\cdots\subset \La_0=k[[t]]^n$；经由标架 $\epsilon$，这样的链给出 $k^r$ 的完全旗，反之亦然。

现设 $r=n$。则 \eqref{nil to Gr} 实际给出 $\phi$ 的一个截面 $\mN_n\to \Gr_{\leq \mu}^{\square}$（因为 $(t-A)\La_0$ 有典范平凡化）。于是得到下述笛卡尔方块
\begin{equation}\label{sat and springer2}
\begin{CD}
\tilde\mN_n@>>>\Gr_{\leq(\omega_1,\ldots,\omega_1)}\\
@VpVV@VVmV\\
\mN_n@>>>\Gr_{\leq \mu}.
\end{CD}
\end{equation}
换言之，此情形下的卷积映射 $m$ 延拓了 $\GL_n$ 的 Springer 消解。
\end{ex}

继续一般讨论。下述基本定理概括了 Schubert 簇奇点的基本事实。

\begin{thm}\label{singularity}
以 $p$ 记 $k$ 的特征指数。若 $p\nmid |\pi_1(G_\der)|$，则 $\Gr_{\leq \mu}$ 正规、Cohen--Macaulay、Gorenstein，且有有理奇点。
\end{thm}

\begin{proof}
这是个困难定理。除 Gorenstein 性外，所有论断都在 \cite[Theorem 8]{Fa} 中证明，另见 \cite[Theorem 0.3]{PR}；关键一步是证明 $\Gr_{\leq \mu}$ 正规。

$\Gr_{\leq \mu}$ 的 Gorenstein 性由 \cite[Equation (241)]{BD} 推出。严格说来，\emph{同上}的论证只在 $k$ 的特征为零（或对 $G$ 为极好素数）时有效。更多细节及向 $p\nmid |\pi_1(G_\der)|$ 情形的推广见 \cite[Theorem 6.11]{Z10}，相关讨论另见注 \ref{symplectic sing}。
\end{proof}

\begin{rmk}
简略讨论 $\underline G$ 为 $G$ 的 Iwahori 群概形时的平行故事，细节见 \cite{Fa,PR}。记 $I=\underline G(\mO)$。此时 Cartan 分解形如
\[I\backslash G(F)/I\simeq \widetilde W,\]
其中 $\widetilde W\simeq \xcoch(T)\rtimes W$ 称为 \textbf{Iwahori Weyl 群}。它是（拟）Coxeter 群，带有 Bruhat 序 $\leq$ 与长度函数 $\ell:\widetilde W\to\bZ$。对 $w\in \widetilde{W}$，以 $S_w$ 记 $\Gr_{\underline G}=\Fl_G$ 上对应的 $L^+\underline G$-轨道，以 $S_{\leq w}$ 记其闭包。则 $S_w$ 同构于维数 $\ell(w)$ 的仿射空间，且 $S_{\leq w}=\cup_{w'\leq w}S_w$。所有 $S_{\leq w}$ 都正规、Cohen--Macaulay 且有有理奇点，但一般不是 Gorenstein。
\end{rmk}

% ============================================================
%  ch06  第二讲 §2.2 漫谈、§2.3 对偶 Schubert「簇」与横截片
%  原文行 1467–1651
% ============================================================

\subsection{插叙：$\Gr$ 中的一些子归纳概形}\label{digression}

这里插叙一种构造 $\Gr$ 中若干重要子簇的方法。本小节在余下部分不再使用。

设 $\varphi: \Gr\to \Gr$ 是 $k$ 上态射。对 $\ga\in LG$ 与 $\mu$，定义 $\Gr$ 的闭子归纳概形
\[X(\mu,\ga\varphi)=\left\{ x\in \Gr\mid \inv(x,\ga\varphi(x))\leq \mu\right\}.\]
确切地说，写 $x=g L^+G\in \Gr$，则
\[X(\mu,\ga\varphi)= \left\{g L^+G\mid g^{-1}\ga\varphi(g) \in (LG)_{\leq \mu}\right\}/L^+G.\]

这里给出几个等价定义。

\begin{enumerate}
\item[(1)] $X(\ga\varphi,\mu)$ 由下述笛卡尔方块定义
\[\begin{CD}
X(\mu,\ga\varphi)@>>> \Gr\tilde\times\Gr_\mu\\
@VVV@VVV\\
\Gr@>1\times \ga\varphi>> \Gr\times \Gr.
\end{CD}\]

\item[(2)] 设 $\varphi$ 由保持 $L^+G$ 的态射 $\varphi:LG\to LG$ 诱导。考虑 $LG$ 被 $\varphi$ 扭变后的共轭作用，即 $\on{Int}(g)\cdot^\varphi g'= g^{-1}g' \varphi(g)$。考虑叠的态射
\[\pi: [(LG)_\mu/^\varphi L^+G]\to [LG/^\varphi LG].\]
则 $X(\mu,\ga\varphi)= \pi^{-1}(\ga)$。
\end{enumerate}

\begin{rmk}
许多时候，商 $[LG/^\varphi LG]$ 虽然性质不佳，仍应被看作某些重要代数结构的模空间。
\end{rmk}

下面是上述构造的几个具体例子。

\begin{enumerate}
\item[(i)] 取 $\varphi=\id$ 为恒等映射。则 $X(\ga,\mu)$ 是群版本仿射 Springer 纤维（见 \cite[\S~2.7.3]{Yu}）。此时 $[LG/LG]$ 分类 $G$ 的共轭类。该构造也有 Lie 代数版本，给出（通常的）仿射 Springer 纤维（见 \cite[\S~2.2]{Yu}）。

\item[(ii)] 设 $k=\bF_q$，$\varphi=\sigma$ 是 $\Gr$ 的 $q$-Frobenius。则 $X(\mu,\ga\varphi)$ 通常记作 $X(\mu,b)$，称为\textbf{仿射 Deligne--Lusztig 簇}（最早在 \cite{Ra} 中引入）。此时 $[LG/^{\sigma}LG]$ 分类 $G(\bar{k}((t)))$ 的 $\sigma$-共轭类，它们与 Drinfeld shtuka 密切相关。

\item[(iii)] 仍设 $k=\bF_p$。设 $\varphi:\Gr\to \Gr$ 由绝对 Frobenius $k[[t]]\mapsto k[[t]]$，$\sum a_it^i\mapsto \sum a_i^pt^{ip}$ 诱导。把 $\ga\varphi$ 记作 $\Phi$。则 $X(\Phi,\mu)$ 称为 \textbf{Kisin 簇}，而 $[LG/^{\sigma}LG]$ 分类 Galois 表示 $\Gal(\overline F/ F)\to G(\bF_p)$（见 \cite{PR2}）。

\item[(iv)] 设 $k=\bC$，$\varphi$ 由环路旋转 $t\mapsto qt$ 诱导，其中 $q\in \bG_m^{\on{rot}}$。则 $[LG/^{\varphi}LG]$ 与椭圆曲线 $E_q=\bC^\times/q^\bZ$ 上的 $G$-丛密切相关（Looijenga、Baranovsky 与 Ginzburg \cite{BG}）。

\item[(v)] 不取 $\varphi\in \bG_m^{\on{rot}}$，也可以让 $\varphi\in \Aut^{++}(D)$。该构造的 Lie 代数版本在 \cite{FZ} 中研究过。
\end{enumerate}

每种情形下都值得问几个基本问题：
\begin{enumerate}
\item 该簇何时有限维？
\item 若有限维，是否有维数公式？各不可约分支是否等维？
\item 如何参数化连通分支？
\item 如何参数化不可约分支？
\end{enumerate}

这些问题在情形 (i) 与 (ii) 中已了解得较好（已知结果的概述见 \cite{Gortz, Yu} 及其所引文献），其他情形则基本未涉及。

\begin{rmk}由于 \S\ref{p-adic} 的讨论，仿射 Springer 纤维与仿射 Deligne--Lusztig 簇现在对 $p$-进群也有定义。\end{rmk}

\subsection{对偶 Schubert「簇」与横截片}

对偶 Schubert「簇」与 \S\ref{II:Sch var} 引入的 Schubert 簇 $\Gr_{\leq \mu}$ 很不一样。事实上，它们并不是真正的簇，而且无穷维（但余维有限）。

先定义群预层
\[L^-G(R)=G(R[t^{-1}]).\]
\S\ref{Kac-Moody} 将说明 $LG$ 与 Kac--Moody 群密切相关；从这个观点看，$L^-G$ 是 $L^+G$ 的「对偶抛物子」。还有另一种看法：$L^-G$ 是 \eqref{Gout} 中将引入的群 $\underline G^{X^*}$ 的特例（对应 $X=\bP^1$、$x=0$）。特别地，由下文引理 \ref{rep of Gout}，群 $L^-G$ 可由归纳概形表示。

如 \S\ref{II:Sch var} 那样固定 $T\subset B\subset G$。则 $G(F)$ 有 \textbf{Birkhoff 分解}
\begin{equation}\label{Birk decom}
G(F)=\bigsqcup_{\mu\in \xcoch(T)^+} G(k[t^{-1}])t^\mu G(\mO),
\end{equation}
它诱导
\begin{equation}\label{Birk decom 2}
L^-G(k)\backslash G(F)/G(\mO)\cong \xcoch(T)^+.
\end{equation}
与 Cartan 分解 \eqref{Cartan decomp2} 类似，\eqref{Birk decom 2} 不依赖任何选取。一般 $G$ 情形的证明见 \cite[Lemma 4]{Fa}。

把 $L^-G$ 在 $\Gr$ 上的轨道称为\textbf{对偶 Schubert 胞腔}。由 Birkhoff 分解 \eqref{Birk decom 2}，它们形如
\[\Gr^\mu=L^-G\cdot t^\mu \subset \Gr,\]
并由 $\xcoch(T)^+$ 参数化。令
\[\Gr^{\geq \mu}=\bigsqcup_{\la\geq \mu}\Gr^\la,\]
称为\textbf{对偶 Schubert 簇}。下述命题说明这个术语是合理的。

\begin{prop}\label{geom of opp Sch}~
\begin{enumerate}
\item $\Gr^\mu$ 是 $\Gr$ 的局部闭子归纳概形。

\item $\Gr_\mu\cap \Gr^\la\neq \emptyset$ 当且仅当 $\mu\geq \la$。此外 $\Gr_\mu\cap \Gr^\mu= G\cdot t^\mu\cong G/P_\mu$，它在旋转环面 $\bG_m^{\on{rot}}$ 下不动。

\item $\Gr^{\geq \mu}$ Zariski 闭，且以 $\Gr^\mu$ 为稠密开子集。

\item $\Gr^{\geq \mu}$ 的余维为 $(2\rho,\mu)-\dim G/P_\mu$。
\end{enumerate}
\end{prop}

\begin{proof}
(1) $t^\mu$ 在 $L^-G$ 中的稳定化子是 $L^-G\cap t^\mu L^+G t^{-\mu}$，它是有限维子群。把 $L^-G$ 写成 $\varinjlim K_i$，使 $K_i\supset L^-G\cap t^\mu L^+G t^{-\mu}$，则每个 $K_it^\mu$ 在 $\Gr$ 中局部闭，故 $\Gr^\mu=\varinjlim K_it^\mu$ 是归纳概形。还要注意
\begin{equation}\label{dim opp Sch}
\dim(L^-G\cap t^\mu L^+G t^{-\mu})=(2\rho,\mu)+\dim P_\mu,
\end{equation}
这可由计算其李代数的维数得到，论证同命题 \ref{geom of Sch}。

(2) 留作练习。

(3) 先说明 $\Gr^0$ 在 $\Gr$ 的单位连通分支中开。注意有求值映射
\[L^-G\to G,\quad g\mapsto g\bmod t^{-1}.\]
以上述映射的核定义 $L^{<0}G=\ker(L^-G\to G)$。注意 $\Gr^0= L^{<0}G\cdot e$。

\begin{lem}\label{big cell}
乘法映射
\[L^{+}G\times L^{<0}G\to LG\]
是开嵌入。特别地，$L^{<0}G\cong L^{<0}G\cdot e\subset \Gr$ 是开浸入。
\end{lem}

\begin{proof}
注意 $L^{<0}G\cap L^+G=\{1\}$，且在李代数层面 $L\frakg=L^+\frakg\oplus L^{<0}\frakg$，其中 $L^+\frakg=\frakg(\mO)$，$L^{<0}\frakg=t^{-1}\frakg\otimes k[t^{-1}]$。因此 $L^{+}G\times L^{<0}G\to LG$ 是开嵌入。
\end{proof}

现在考虑 $L^-G\cdot t^\mu\cdot \Gr^0$。它是开集，且是 $L^-G$-轨道的并。容易看出 $t^\la$ 属于它，当且仅当 $\Gr^\la\cap \Gr_\mu\neq \emptyset$，故由 (2) 当且仅当 $\la\leq \mu$。因此 $\bigsqcup_{\la\leq \mu}\Gr^\la=L^-G\cdot t^\mu \cdot \Gr^0$ 开。于是 $\Gr^{\geq \mu}$ 闭。还需说明对每个 $\la>\mu$，$\Gr^\la$ 在 $\Gr^\mu$ 的闭包中。如命题 \ref{geom of Sch} 那样，可设 $\la=\mu+\al$，其中 $\al$ 是某个正根。取 \eqref{curve C} 定义的曲线 $C_{\mu+\al,\al}$，则容易看出 $t^{\mu+\al}\in C_{\mu+\al,\al}$，而补集 $\bA^1\subset \Gr^\mu$。证毕。

(4) 由 (3) 与 \eqref{dim opp Sch} 推出。
\end{proof}

\begin{rmk}
(i) 开子集 $\Gr^0$ 通常称为 $\Gr$ 的「大开胞腔」。

(ii) 正如刚才所见，$L^+G$ 与 $L^-G$ 扮演的角色不同，因此考虑商 $LG/L^-G$ 也很自然。由下文定理 \ref{uniformization for P1}，它可解释为 $\bP^1$ 上带一个平凡化的 $G$-主丛的模空间，平凡化取在 $D_0$（$0\in \bP^1$ 处的圆盘）上。文献中有时把 $LG/L^-G$ 称为 $G$ 的\textbf{厚仿射 Grassmann 簇}。注意与 $\Gr$ 不同，它具有整体性质（定义涉及一条整体曲线），并且是（无穷型的）概形。
\end{rmk}

理解对偶 Schubert 胞腔与簇的最好办法是下述定理。以 $\Bun_G(\bP^1)$ 记 $\bP^1$ 上 $G$-丛的模叠，它是 $k$ 上局部有限展示的代数叠（例如见 \cite{So}）。

\begin{thm}\label{uniformization for P1}
设 $G$ 在 $k$ 上连通约化。则有典范同构
\[ [L^-G\backslash \Gr]\simeq \Bun_G(\bP^1).\]
事实上，映射 $\Gr\to \Bun_G(\bP^1)$ 把 $\Gr$ 呈现为 $\Bun_G(\bP^1)$ 上 \'etale 拓扑下的 $L^-G$-主丛。
\end{thm}

\begin{proof}
该论断等价于说：经过 \'etale 覆盖 $R\to R'$ 之后，$\bP^1\times \Spec R$ 上每个 $G$-丛在 $(\bP^1-\{0\})\times \Spec R'$ 上平凡。这由两点推出：(i) 经过可能的 \'etale 基变换 $R\to R'$，$X_R$ 上每个 $G$-丛都可约化到 Borel（\cite[Theorem 1]{DS}）；(ii) 由于 $B$ 有滤过、其伴随分次为若干 $\bG_a$ 与 $\bG_m$，$(\bP^1-\{0\})\times \Spec R$ 上每个 $B$-丛经可能的 \'etale 基变换 $R\to R'$ 后平凡。
\end{proof}

在 $k$-点层面，该定理说
\[ \Bun_{G}(\bP^1)(k)\cong G(k[t^{-1}])\backslash G(k((t)))/ G(k[[t]]).\]
特别地，Grothendieck 关于 $\bP^1$ 上向量丛的定理等价于（$G=\GL_n$ 的）Birkhoff 分解 \eqref{Birk decom}。确切地说，对应于 $t^\mu$ 的双陪集的 $G$-丛，是由线丛 $\mO_{\bP^1}(1)$ 沿 $\bG_m\stackrel{\mu}{\to}T\subset G$ 推出的。把这个 $G$-丛记作 $\mO(\mu)$，并把它在 $\Bun_G(\bP^1)$ 上对应的点记作 $[\mO(\mu)]$。于是命题 \ref{geom of opp Sch} 可改述为

\begin{prop}\label{geom of strata}~
\begin{enumerate}
\item 点 $[\mO(\la)]$ 在 $[\mO(\mu)]$ 的闭包中，当且仅当 $\la\geq \mu$；
\item $[\mO(\mu)]$ 的闭包是维数为 $-(2\rho,\mu)-\dim P_\mu$ 的闭子叠。
\end{enumerate}
\end{prop}

作为命题 \ref{geom of opp Sch} 的应用，还可以引入横截片。

\begin{prop}
$L^{<0}G\cdot t^\la$ 是 $\Gr_\la$ 在 $\Gr$ 中的横截片。更确切地说，对每个 $\la\leq \mu$，令 $S^\la_{\leq \mu}=L^{<0}G\cdot t^\la\cap \Gr_{\leq \mu}$，则 $S^\la_{\leq \mu}$ 与每个 $\Gr_\nu$（$\la\leq \nu\leq \mu$）横截相交。
\end{prop}

\begin{proof}由引理 \ref{big cell} 推出。
\end{proof}

\begin{rmk}\label{symplectic sing}
簇 $S^\la_{\leq \mu}$ 是带辛奇点的簇的例子（见 \cite[Theorem 2.7]{KWWY}），该概念的定义见 \cite[Definition 1.1]{Bea}。由 \cite[Proposition 1.3]{Bea}，辛奇点是有理 Gorenstein 的，这与定理 \ref{singularity} 一致。
\end{rmk}

下面讨论几个例子。

\begin{ex}\label{slice GLn}
考虑 $G=\GL_n$，$\mu=(n,0,\ldots,0)$。首先有

\begin{lem}\label{Lus obs}
映射 \eqref{nil to Gr} 是开嵌入。
\end{lem}

\begin{proof}
事实上，映射 $\mN_n\to L^{<0}\GL_n$，$A\mapsto 1-At^{-1}$，是闭嵌入，故由引理 \ref{big cell}，映射 $\mN_n\to \Gr_{\GL_n}$，$A\mapsto (t-A)\La_0$，是局部闭嵌入且像含于 $\Gr_{\leq \mu}$。因此 \eqref{nil to Gr} 本身是局部闭嵌入。由于 $\Gr_{\leq \mu}$ 不可约且既约、$\dim \mN_n=\dim \Gr_{\leq \mu}=n(n-1)$，\eqref{nil to Gr} 必为开嵌入。
\end{proof}

现取 $\la=(1,\ldots,1)$。注意 $\Gr_\la$ 是一个点。于是引理 \ref{Lus obs} 等价于说 \eqref{nil to Gr} 实现了横截片 $S_{\leq \mu}^\la\subset \Gr_{\leq \mu}$。因此 $S_{\leq \mu}^\nu$ 同构于 $\mO_\nu$ 的 Slodowy 片，其中 $\la\leq \nu\leq \mu$。
\end{ex}

\begin{ex}
设 $G_\der$ 单，$\mu=\theta$ 拟极小。由与上面相同的构造，$S^0_{\leq \theta}$ 同构于 $\frakg$ 的极小幂零轨道闭包 $\overline{\mO}_{\on{min}}$。
\end{ex}

\begin{rmk}
注意定理 \ref{singularity} 蕴含：$\GL_n$ 的幂零轨道以及一般半单李代数 $\frakg$ 的极小幂零轨道都正规、Cohen--Macaulay、Gorenstein，且有有理奇点。这对半单李代数中一般的幂零轨道并不成立，这说明上面横截片与幂零轨道之间的关系不能推广到一般情形。
\end{rmk}

% ============================================================
%  ch07  第二讲 §2.4 Picard 群、§2.5 中心扩张与仿射 Kac--Moody 代数
%  原文行 1652–1765
% ============================================================

\subsection{Picard 群}\label{The Picard group}

为简单起见，设 $G$ 单且单连通。以 $\theta$ 记短支配余权，前面已说明它是 $\xcoch(T)^+-\{0\}$ 的极小元。由命题 \ref{geom of opp Sch}，$\Theta:=\Gr^{\geq \theta}$ 余维为一，且 $\Gr=\Gr^0\sqcup \Theta$。下面说明它是 $\Gr$ 上某个丰沛线丛的整体截面的零点集。

\begin{prop}\label{big cell 2}
存在丰沛线丛 $\mL$ 与整体截面 $\sigma\in\Ga(\Gr,\mL)$，使得 $\Gr^0=\{x\in\Gr \mid \sigma(x)\neq 0\}$。
\end{prop}

\begin{proof}
这里说明想法，细节见 \cite[Lemma 2, Corollary 3]{Fa}。

先考虑 $G=\SL_n$。此时 $L^{<0}G\cdot e$ 是 $\mL_{\det}$ 的某个整体截面 $\sigma$ 的非零集，$\sigma$ 构造如下。令 $L_0=(t^{-1}k[t^{-1}])^n$，考虑 $\Spec R$ 上的两项复形
\[0\to R\otimes L_0\oplus\Lambda\stackrel{\Upsilon}{\to} R\otimes L_0\oplus R[[t]]^n\to 0,\quad \Upsilon(v,w)=(v,0)+\beta(w),\]
其中 $\beta: \La\to R((t))=R\otimes L_0\oplus R[[t]]$ 是自然包含。由定义，该复形（更确切地说，把两项都除以 $ R\otimes L_0\oplus t^NR[[t]]^n$（$N\gg 0$）后所得的复形）的行列式上同调恰是 \S\ref{determinant line} 中构造的行列式线丛 $\mL_{\det}$，而 $\sigma=\det(\Upsilon)$ 定义 $\mL_{\det}$ 的整体截面。把点 $x\in \Gr$ 写成 $g\La_0$（$g\in LG$），则由定义 $\sigma(g\La_0)\neq 0$ 当且仅当 $\Upsilon$ 在该点处是同构。而后一条件又等价于 $g^{-1} L_0=L_0$（见 \cite[Lemma 2]{Fa}）。这就完成了 $\SL_n$ 情形的证明。

其次，取表示 $\rho:G\to \GL(V)$，使得存在 $v\in V$，其在 $\GL(V)$ 中的稳定化子为 $G$（这样的表示存在，因为 $G$ 半单）。设 $f_\rho: \Gr_G\to \GL_{\SL_n}$ 是诱导映射，可以证明 $f_\rho^{-1}(L^{<0}\SL_n\cdot e)=L^{<0}G\cdot e$。于是 $\Theta_G$ 是 $\Theta_{\SL_n}$ 的拉回（作为闭子集）。
\end{proof}

回忆预叠上线丛的定义（见 \S\ref{Intro:sp}）。设 $e\in X(k)$ 是点，定义 $X$ 上在 $e$ 处刚性化的线丛的 Picard 群胚 $\Pic^e(X)$ 为由 $(\mL,\epsilon)$ 构成的群胚，其中 $\mL$ 是 $X$ 上的线丛，$\epsilon: e^*\mL\simeq k$ 是平凡化。注意若 $X$ 是真代数簇，则 $\Pic^e(X)$ 离散，退化为 $X$ 的通常 Picard 群。

\begin{thm}\label{Pic of Gr}
设 $G$ 单且单连通。给 $\Theta$ 赋以既约闭归纳概形结构，则 $\Theta$ 是有效 Cartier 除子，相应的线丛 $\mO(1):=\mO(\Theta)$ 丰沛。此外 $\Pic^e(\Gr)=\bZ\mO(1)$。
\end{thm}

\begin{rmk}\label{generator}
若 $\rho: G\to \GL_n$ 是表示，则 $\mL_{\det}$ 沿 $f_\rho:\Gr_G\to \Gr_{\GL_n}$ 的拉回是 $\mO(1)$ 的正倍数，即 $f_\rho^*\mL_{\det}=\mO(d_\rho)$，其中正整数 $d_\rho$ 称为表示 $\rho$ 的 \textbf{Dynkin 指数}。$d_\rho$ 的一般公式见 \cite[Lemma 5.2]{KNR}，另见 \cite[Lemma 10.6.1]{So} 与 \cite[Lemma 4.2]{YZ}（注意上述文献都只考虑 $\cha\ k=0$，但 \cite[Lemma 4.2]{YZ} 的论证对任意特征都适用）。

例如，若 $\rho$ 是伴随表示，则 $d_\rho=2h^\vee$，其中 $h^\vee$ 是 $G$ 的对偶 Coxeter 数。若 $\rho$ 是 $\SL_n$ 或 $\on{Sp}_{2n}$ 的标准表示，则 $d_\rho=1$。若 $G$ 是正交群 $\on{SO}_n$（设 $\cha\ k>2$）、$\rho$ 是标准表示，或 $G=G_2$、$\rho$ 是唯一的 $7$ 维表示，则 $d_\rho=2$，其唯一的平方根可由 $\on{OGr}(V)$ 上 Pfaffian 线丛 $\mL_{\on{Pf}}$ 的拉回得到（见注 \ref{Pfaffians}）。

换个角度看，上述讨论本质上给出了 $G$ 为典型群或 $G=G_2$ 时 $\mO(1)$ 的构造。但若 $G=E_8$，则对任意 $\rho$ 都有 $60\mid d_\rho$，因此在此情形下找出 $\mO(1)$ 的显式构造是个有趣的问题。
\end{rmk}

\begin{rmk}\label{square root of adjoint}
由此可知，对伴随表示 $\Ad: G\to \GL(\frakg)$，$f_{\Ad}^*\mL^{-1}_{\det}$ 有唯一的刚性化平方根，记作 $\mL_{c}$。这就是仿射 Grassmann 簇上所谓的\textbf{临界线丛}，它在几何朗兰兹纲领中起重要作用（见 \cite{BD}）。
\end{rmk}

证明该定理的一条途径是把 $\Gr$ 等同于相应 Kac--Moody 群的部分旗簇（见下文 \S\ref{Kac-Moody}），然后引用 \cite{Ma} 的结果。下面是更直接的途径：$\mO(1)$ 应当是与除子 $\Theta$ 相配的线丛。但 $\Gr$ 不是有限维的，一般也不能写成光滑有限维子簇的递增并，故其真正的定义依赖下述引理。

\begin{lem}\label{O1 on Bun}
设 $G$ 单且单连通。$\Bun_G(\bP^1)$ 有开子叠 $\mathring{\Bun}_G(\bP^1)$，分类 $\bP^1$ 上可平凡化的 $G$-主丛，且边界 $\Bun_G\setminus \mathring{\Bun}_G$ 纯余维一。因此它是有效 Cartier 除子，在 $\mathring{\Bun}_G(\bP^1)$ 上定义了线丛 $\mO(1)$。
\end{lem}

\begin{proof}
有 $\mathring{\Bun}_G(\bP^1)= [L^-G\backslash L^{<0}G\cdot e]$，而 $\Bun_G\setminus \mathring{\Bun}_G=[L^-G\backslash \Theta]$ 纯余维一。也可以用形变理论证明。
\end{proof}

于是拉回在 $\Gr$ 上定义了线丛，暂记作 $\mL(\Theta)$。注意 $\mL(\Theta)$ 丰沛，因为由命题 \ref{big cell}，它的某个倍数是行列式线丛的拉回。

其次，不难证明 \eqref{curve C} 中构造的有理曲线 $C_{\theta,\theta}$ 在 $t^\theta$ 处与 $\Theta$ 横截相交（例如看切空间），其中 $\theta$ 仍是短支配余根。此外可以证明限制到 $C_{\theta,\theta}$ 诱导 Picard 群的单射（用 Demazure 消解）。于是有包含 $\bZ\mL(\Theta)\subset \Pic^e(\Gr)\subset\Pic(\bP^1)$，它把 $\mL(\Theta)$ 送到 $\mO_{\bP^1}(1)$。定理随之成立。更多细节见 \cite[Corollary 12]{Fa} 与 \cite[\S~10.2]{So}。

下述论断留作练习。

\begin{lem}\label{lifting Aut action}
$\Aut^+(D)$ 在 $\Gr$ 上的作用提升为 $\mO(1)$ 上的唯一作用，它在 $\mO(1)$ 于 $e$ 处的纤维上为平凡作用。
\end{lem}

\subsection{中心扩张与仿射 Kac--Moody 代数}\label{Kac-Moody}

仍设 $k$ 代数闭，$F=k((t))$。回忆仿射 Dynkin 图分类了两类不同的群。设 $\Ga$ 是连通仿射 Dynkin 图，$S$ 是 $\Ga$ 的顶点子集。Bruhat--Tits 的工作对 $\Ga$ 给出 $F$ 上单且单连通的约化群 $G_\Gamma$，对 $S$ 给出 $\mO$ 上的抛孔群概形 $\underline G_S$（概述见 \cite{Ti}，分类见\emph{同上} \S 4.2）。另一方面，对 $\Ga$ 有 Kac--Moody 代数 $\frakg_\Ga$（\cite[Chap. 4]{Kac}）与相应的 Kac--Moody 群 $G_{\Ga}$，对 $S$ 有 $G_{\Ga}$ 的抛物子群 $P_S$\footnote{实际情况要复杂得多：由于 Moody--Teo、Marcuson、Garland、Kac--Peterson、Slodowy、Tits、Mathieu 等人的工作，附着于 $\frakg_\Ga$ 的无穷维群有多种版本（极小的、形式的等等）。概述与文献见 \cite{Ti3}。}。

下面通过解释下图中虚线箭头，说明 Bruhat--Tits 理论与 Kac--Moody 理论的关系。
\begin{center}
\begin{tikzpicture}
  \node (a) at (0,1.7) {$\{\text{$F$ 上的约化群}\}$};
  \node (b) at (-3.2,0) {$\{\text{仿射 Dynkin 图}\}$};
  \node (c) at (0,-1.7) {$\{\text{$k$ 上的仿射 Kac--Moody 代数/群}\}$};
  \draw[dashed, ->] (a) -- (c);
  \draw[->] (b) -- (a);
  \draw[->] (b) -- (c);
\end{tikzpicture}
\end{center}

这里假定 $\Ga$ 是非扭仿射 Dynkin 图\footnote{扭情形更复杂，尤其当 $k$ 特征很小时。}，即 \cite[Chap. 4, TABLE Aff 1]{Kac} 所列的仿射 Dynkin 图。它们可由有限 Dynkin 图（\cite[Chap 4, TABLE Fin]{Kac} 所列）添加一个特殊顶点 $v_0$ 得到。对应于这些 Dynkin 图的 $F$ 上约化群，是 $k$ 上单且单连通约化群作基变换所得。若 $S=\{v_0\}$，则相应的抛孔群概形超特殊，形如 $G\otimes_kk[[t]]$（见例 \ref{hyperspacial and Iwahori}）。

此时，$\Gr_G$ 上的线丛 $\mO(1)$ 可用来定义 $LG$ 被 $\bG_m$ 的中心扩张
\begin{equation}\label{central extension}
1\to\bG_m\to \widehat{LG}\to LG\to 1.
\end{equation}
确切地说，$\widehat{LG}$ 的 $R$-点是 $LG(R)$ 中的 $R$-点 $g$ 连同一个同构
\[g^*\mO(1)\simeq \mO(1).\]
旋转环面的作用典范地提升为 $\widehat{LG}$ 上的作用（由引理 \ref{lifting Aut action}）。令
\[\widetilde{LG}:=\widehat{LG}\rtimes\bG_m^{\on{rot}}.\]

\begin{thm}
归纳群 $\widetilde{LG}$ 同构于由相应非扭仿射 Dynkin 图构造的 Kac--Moody 群 $G_\Ga$ 的（形式版本）。
\end{thm}

这可由下述更精确的定理推出。设 $\underline G$ 是 $G$ 的抛孔群概形，对应于 $G$ 的仿射 Dynkin 图的顶点子集 $S$；设 $P_S\subset G_\Ga$ 是仿射 Kac--Moody 群中相应的抛物子群。则

\begin{thm}
仿射 Grassmann 簇 $\Gr_{\underline G}$ 同构于 \cite{Ma,Ku} 中定义的部分旗簇 $G_\Ga/P_S$。
\end{thm}

\begin{proof}
$G=\SL_n$ 的情形见 \cite{BL}，一般 $G$ 见 \cite{Fa}（另见 \cite[\S 10]{PR}）。想法是先证明 $\Gr_{\underline G}$ 中的 Schubert 簇同构于 Kac--Moody 理论中出现的（先验地不同的）Schubert 簇。
\end{proof}

\begin{rmk}\label{2-cocycle}
$\widehat{LG}$ 的李代数 $\widehat{L\frakg}$ 是 $\frakg(F)$ 被 $k$ 的中心扩张。若 $\cha\ k=0$，它可显式写成 $\frakg\otimes F\oplus k\mathbf{1}$，其李括号由
\[[X\otimes f,\mathbf{1}]=0,\quad [X\otimes f,Y\otimes g]=[X,Y]\otimes fg- (\Res_{t=0}fdg) (X,Y)\mathbf{1}\]
给出，其中 $X,Y\in \frakg$，$f,g\in F$，$(-,-): \frakg\times\frakg\to k$ 是 $\frakg$ 上规范化的不变双线性型，$\Res: Fdt\to k$ 是留数映射。细节见 \cite[\S~7]{Kac}。\zy{关于 nLab：其词条 central extension（\url{https://ncatlab.org/nlab/show/central+extension}，2025-02-16 修订）目前只有导引性内容、未给可引用的定义，故本条不以该词条为依据，而依 \cite[Chap.~4, \S~7]{Kac}：$\widehat{L\frakg}=\frakg(F)\oplus k\mathbf{1}$ 的括号由留数给出，这与 $\widehat{LG}$ 的中心扩张 \eqref{central extension} 在无穷小层面相对应。[来源：A 级 \cite{Kac}；nLab central extension 仅记其现状。]}
\end{rmk}

结合 \cite{Ku,Ma} 的主要结果，可得 Borel--Weil 型定理。为简单起见设 $\cha\ k=0$，并对非扭仿射 Dynkin 图 $\Ga$ 与 $S=\{v_0\}$ 叙述。

\begin{thm}\label{Borel-Weil}
整体截面空间的对偶 $\Ga(\Gr_G,\mO(m))^*$ 同构于 $L_m(\frakg_\Ga)$，即 $\frakg_{\Ga}$ 的第 $m$ 层可积表示。
\end{thm}

Kac--Moody 群的可积表示是约化代数群有限维表示的类比，可用纯表示论方法构造。下面简要回忆 $L_m(\frakg_\Ga)$ 的定义，沿用 \cite{Kac} 的记号。以 $\Lambda_0$ 记 $\frakg_\Ga$ 的基本权（定义见\emph{同上} \S\ 6），以 $M(m\Lambda_0)$ 记 $\frakg_\Ga$ 的最高权 $m\Lambda_0$ 的 Verma 模（定义见\emph{同上} \S\ 9）。则 $L_m(\frakg_\Ga)$ 是 $M(m\Lambda_0)$ 对由\emph{同上} 公式 (10.1.1) 中那些向量生成的子模的商。

% ============================================================
%  ch08  第三讲 §3.1：Beilinson--Drinfeld Grassmann 簇
%  原文行 1766–2013
% ============================================================

\section{第三讲：Beilinson--Drinfeld Grassmann 簇与分解结构}\label{Lecture III}

本讲引入曲线上的仿射 Grassmann 簇的版本，即 Beilinson--Drinfeld Grassmann 簇，并讨论它们的一些基本几何性质，特别是分解性质。

\subsection{Beilinson--Drinfeld Grassmann 簇}

Beilinson--Drinfeld 的重大发现是：经由定理 \ref{glob aff} 对仿射 Grassmann 簇的解释可以有极为广泛的推广。

首先，可以让定理 \ref{glob aff} 中的点 $x$ 变动。对点 $x\in X(R)$，以 $\Ga_x$ 记映射 $x:\Spec R\to X$ 的图像。于是可定义预叠
\[\Gr_{\underline G, X}(R)=\left\{(x,\mE,\beta)\ \left|\ \begin{array}{l}
x\in X(R),\quad \mE \text{ 是 } X_R \text{ 上的 } \underline{G}\text{-主丛},\\
\beta: \mE|_{X_R-\Ga_x}\cong \mE^0|_{X_R-\Ga_x} \text{ 是平凡化}
\end{array}\right.\right\}.\]
但事实上还可以做得更多：允许有多个点。设 $I$ 是有限非空集。对点 $x=\{x_i,i\in I\}\in X^I(R)$，以
\[\Ga_x:=\bigcup_{i\in I} \Ga_{x_i}\]
记各映射 $x_i: \Spec R\to X$ 的图像之并（作为 $X_R$ 的闭子概形）。

\begin{dfn}
在曲线自积 $X^I$ 上定义预层
\begin{equation}\label{BD Grass def}
\Gr_{\underline G, X^I}(R)=\left\{(x,\mE,\beta)\ \left|\ \begin{array}{l}
x\in X^I(R),\quad \mE \text{ 是 } X_R \text{ 上的 } \underline{G}\text{-主丛},\\
\beta: \mE|_{X_R-\Ga_x}\cong \mE^0|_{X_R-\Ga_x} \text{ 是平凡化}
\end{array}\right.\right\}.
\end{equation}
以 $q_I:\Gr_{\underline G,X^I}\to X^I$ 记结构映射。\zy{$\Gr_{\underline G,X^I}$ 是「把抛物结构放在 $I$ 个点上」的族；\S\ref{Ran space} 用 Ran 空间把这些族按点集之并粘成 $\Gr_{\Ran}$，分解性质即在该对象上表述。nLab 词条 Ran space 与 factorization algebra 记的正是这一框架。[来源：nLab Ran space（2020-07-27）、factorization algebra（2026-08-16）。]}
\end{dfn}

\begin{thm}\label{representability of BD}
预层 $\Gr_{\underline G, X^I}$ 可由 $X^I$ 上有限型的归纳概形表示。若 $\underline G$ 约化，则它归纳射影。
\end{thm}

\begin{proof}
同样可用类似命题 \ref{lin rep} 与 \ref{gen G} 的论证化归到 $\underline G=\GL_n$ 的情形。然后只需证明对正整数 $N$，
\[\Gr_{\GL_n,X^I}^{(N)}=\left\{(x,\mE,\beta)\in\Gr_{\GL_n,X^I}\mid \mO_{X_R}^n(-N\Ga_x)\subset \mE\subset \mO_{X_R}^{n}(N\Ga_x)\right\}\]
可由 $X^I$ 上射影概形表示。用引理 \ref{projectivity} 中的论证，映射 $\mE\mapsto \mF:=\mO_{X_R}^{n}/\mE(-N\Ga_x)$ 把它等同于下述函子：给 $R$ 配上对 $(x,\mF)$，其中 $x\in X(R)$，而 $\mF$ 是 $R$-平坦的 $\mO_{X_R}$-模，且是 $\mO_{X_R}^n/\mO_{X_R}^n(-2N\Ga_x)$ 的商。其可表示性随即由商概形（代替引理 \ref{Embedding} 中的 Grassmann 簇）的可表示性推出。
\end{proof}

在不引起混淆时，把 $\Gr_{\underline G,X^I}$ 简记作 $\Gr_{X^I}$。

\begin{rmk}\label{projectivity of BD}
定理 \ref{representability of BD} 的最后一个论断在更一般的情形下也成立。称 $X$ 上光滑仿射群概形 $\underline G$ 为 \textbf{Bruhat--Tits 群概形}（或\textbf{抛孔群概形}），若它在每个 $D_x$（$x\in |X|$）上的限制都是 Bruhat--Tits 意义下的抛孔群概形。用命题 \ref{indproper}，类似 \cite{PZ} 的论证表明 $\Gr_{X^I}$ 在 $X^I$ 上归纳真。
\end{rmk}

类似地可定义多点版本的射流群与环路群。对 $x\in X^I(R)$，以 $\hat{\Ga}_x$ 记图形 $\Ga_x$ 在 $X_R$ 中的形式完备化（它是形式概形）。有典范映射 $\pi:\hat{\Ga}_x\to X_R$。
\begin{equation}
(L^+\underline G)_{X^I}=\left\{(x,\beta)\mid x\in X^I(R),\ \beta\in
\underline{G}(\hat{\Ga}_x)\right\}.
\end{equation}

\begin{prop}
$(L^+\underline{G})_{X^I}$ 可由 $X^I$ 上仿射概形表示。
\end{prop}

为严格定义整体环路群，需要下述引理。

\begin{lem}
形式概形 $\hat{\Ga}_x$ 归纳仿射。以 $A_x$ 记其函数环的拓扑代数，$\hat{\Ga}'_x=\Spec A_x$。有典范映射 $p:\hat{\Ga}_x\to\hat{\Ga}'_x$，则存在唯一态射 $i: \hat{\Ga}'_x\to X_R$ 使 $\pi=ip:\hat{\Ga}_x\to X_R$。
\end{lem}

\begin{proof}见 \cite[Proposition 2.12.6]{BD}。
\end{proof}

于是 $\hat{\Ga}_x^\circ=\hat{\Ga}'_x\setminus \Ga_x$ 有定义，且
\begin{equation}
(L\underline G)_{X^I}=\left\{(x,\beta)\mid x\in X^I(R),\ \beta\in
\underline G(\hat{\Ga}^\circ_x)\right\}.
\end{equation}

由定理 \ref{descent lemma} 与推论 \ref{global interpretation of loop} 中类似的论证，有

\begin{prop}
作为 $k$-空间，$\Gr_{\underline G,X^I}\cong[(L\underline G)_{X^I}/(L^+\underline G)_{X^I}]$。特别地，
\begin{enumerate}
\item $(L\underline G)_{X^I}$ 可由 $X^I$ 上归纳概形表示；
\item $(L\underline G)_{X^I}$ 自然作用在 $\Gr_{\underline G,X^I}$ 上。
\end{enumerate}
\end{prop}

Beilinson--Drinfeld Grassmann 簇确实是极不寻常的对象。先分析 $I=\{1\}$ 与 $I=\{1,2\}$ 的情形。

设 $I=\{1\}$，并先假定 $\underline G=G\times X$ 常值。对每个 $x\in X$，选取同构 $\mO_x\simeq \mO$，便由定理 \ref{glob aff} 得到典范同构
\[\Gr_{X}|_x=\Gr_{\underline G, x}\cong \Gr_{G}.\]
然而整体上一般没有典范等同 $\Gr_X\neq \Gr_G\times X$。相反，$\Gr_X$ 可看作扭曲积 $X\tilde\times\Gr_G$（扭曲积的定义见 \S\ref{Intro:torsor}）。确切地说，$X$ 上有 $\Aut(D)$-主丛 $\hat{X}\to X$，分类 $x\in X$ 连同一个 $D_x\cong D$：
\[\hat{X}(R)=\left\{(x,\al)\mid x\in X(R),\ \al: \hat{\Ga}'_x\cong \Spec R[[t]]\right\}.\]
同样有 $\Aut^+(D)$-主丛 $\hat{X}^+$，由那些限制为典范同构 $\Ga_x\simeq \Spec R$ 的 $\al$ 组成。回忆 $\Aut(D)$ 作用在 $\Gr_G$ 上（由 \eqref{auto of disc} 定义）。于是有同构
\begin{equation}\label{III: twisted constr}
\Gr_X\cong \hat{X}\times^{\Aut(D)}\Gr_G\cong \hat{X}^+\times^{\Aut^+(D)}\Gr_G.
\end{equation}
第一个同构实际上蕴含 $\Gr_X$ 上存在联络（另见命题 \ref{III:property}）。第二个同构允许把
\begin{equation}\label{III:SV over X}
\Gr_{\leq \mu,X}: =X\tilde\times \Gr_{\leq \mu}:=\hat{X}^+\times^{\Aut^+(D)}\Gr_{\leq \mu}
\end{equation}
定义为闭子概形。等价地，$\Gr_{\leq \mu,X}$ 可等同于 $\Gr_X$ 中分类那些满足 $\inv_x(\beta)\leq \mu$ 的点 $(x,\mE,\beta)$ 的闭子集（赋既约概形结构）。

\begin{rmk}
也可以用上述扭曲手续描述 $(L^+G)_X$ 与 $(LG)_X$。
\end{rmk}

仍设 $I=\{1\}$，但 $\underline G$ 不必常值。例如可以构造 $X$ 上的群概形 $\underline G$，使得 $\underline G|_{X-x}\simeq G\times (X-x)$，而 $\underline G|_{D_x}$ 是 $G$ 的 Iwahori 群概形。此时 $\Gr_{\underline G, X}$ 是从仿射 Grassmann 簇到 $G$ 的仿射旗簇的形变（术语见注 \ref{terminology}）：$\Gr_{\underline G, X}$ 在 $X$ 上点 $y\neq x$ 处的纤维同构于 $\Gr_G$，在 $x$ 处的纤维同构于 $\Fl_G$。这个形变在 Gaitsgory 构造中心层时起关键作用 \cite{Ga}。类似的 $p$-进整数环上的形变在 \cite{PZ} 中被用来定义志村簇的局部模型。

下面转向 $\Gr_{X^2}$。设 $\Delta:X\to X^2$ 是对角。

\begin{prop}\label{III:Gr2}
有典范同构
\[\Delta:\Gr_X\cong \Gr_{X^2}\times_{X^2,\Delta}X,\quad c:\Gr_{X^2}|_{X^2-\Delta}\cong (\Gr_X\times\Gr_X)|_{X^2-\Delta}.\]
此外 $\Gr_{X^2}$ 上有自然 $S_2$-作用，使 $c$ 为 $S_2$-等变，其中 $S_2$ 通过交换两个因子作用在 $\Gr_X\times\Gr_X$ 上。
\end{prop}

\begin{proof}
第一个同构显然。为构造第二个，设 $(x,y)\in (X^2-\Delta)(R)$。则 $(X_R-\Ga_x)$ 与 $(X_R-\Ga_y)$ 构成 $X_R$ 的 Zariski 覆盖。定义映射
\[\Gr_{X^2}\times_{X^2}\Spec R\to \Gr_X\times_{X,x}\Spec R\]
为把 $(\mE,\beta)$ 送到 $(\mE_x,\beta_x)$：其中 $\mE_x$ 由 $\mE|_{X_R-\Ga_y}$ 与 $\mE^0|_{X_R-\Ga_x}$ 经 $\beta$ 粘成，因而带有平凡化 $\beta_x: \mE|_{X_R- \Ga_x}\cong \mE^0|_{X_R-\Ga_x}$。同样有映射 $\Gr_{X^2}\times_{X^2}\Spec R\to \Gr_X\times_{X,y}\Spec R$。合起来定义映射 $\Gr_{X^2}|_{X^2-\Delta}\to (\Gr_X\times\Gr_X)|_{X^2-\Delta}$。

反之，若给定 $(\mE_x,\beta_x)\in \Gr_X\times_{X,x}\Spec R$ 与 $(\mE_y,\beta_y)\in \Gr_X\times_{X,y}\Spec R$，可用 $\beta^{-1}_y\beta_x$ 把 $\mE_{x}|_{X_R- \Ga_y}$ 与 $\mE_y|_{X_R-\Ga_x}$ 粘成 $\mE$，按构造它在 $X_R-\Ga_x-\Ga_y$ 上自然带有平凡化 $\beta$。

最后，非平凡元 $\sigma\in S_2$ 通过把 $(x,y, \mE,\beta)$ 送到 $(y,x,\mE,\beta)$ 作用在 $\Gr_{X^2}$ 上。容易验证在 $c$ 之下，该作用对应于 $\Gr_X\times\Gr_X$ 两个因子的交换。
\end{proof}

这些同构与直觉相悖：$\Gr_{X^2}\to X^2$ 在对角上的纤维比一般纤维「更小」！这是纯粹的无穷维现象，因为我们都知道代数簇族的纤维维数在特殊化时会上升。

事实上并不矛盾，有下述结果。

\begin{prop}\label{III:closure}
设 $\underline G=G\times X$ 常值，$p\nmid \pi_1(G_\der)$。经由上述等同，把 $\Gr_{\leq \mu,X}\times \Gr_{\leq \la,X}|_{X^2-\Delta}$ 看作 $\Gr_{X^2}|_{X^2-\Delta}$ 的子集，并以 $\Gr_{\leq (\mu,\la)}$ 记它在 $\Gr_{X^2}$ 中的闭包。则 $\Gr_{\leq (\mu,\la)}|_{\Delta}= \Gr_{\leq \mu+\la,X}$。
\end{prop}

于是只要限制到有限维子簇，直觉就恢复了。该结果的略弱版本在 \cite{Z07} 中证明（另见 \cite{Z10}），但大致相同的论证在此也适用。

下面考虑线丛。设 $\mL$ 是 $\underline G$-主丛的预叠 $\Bun_{\underline G}$ 上的线丛，且在平凡 $\underline G$-主丛处刚性化。注意有自然遗忘映射
\begin{equation}\label{III:I-un}
u_{X^I}: \Gr_{X^I}\to \Bun_{\underline G}.
\end{equation}
以 $\mL_{X^I}$ 记 $\mL$ 到 $\Gr_{X^I}$ 的拉回。

\begin{prop}\label{III:line}
设 $\underline G=G\times X$ 常值，$G$ 单连通。则在命题 \ref{III:Gr2} 给出的等同下，有典范同构
\[\Delta^*\mL_{X^2}\cong \mL_{X}, \quad \mL_{X^2}|_{X^2-\Delta}\cong c^*(\mL_X\boxtimes\mL_X|_{X^2-\Delta}).\]
\end{prop}

\begin{proof}
第一个同构显然，说明第二个。由命题 \ref{III:Gr2} 的证明，下图交换
\[\begin{tikzcd}
(\Gr_X\times X)|_{X^2-\Delta}\arrow[rr, "c^{-1}(\id\times e)"]\arrow[dr] & & \Gr_{X^2}|_{X^2-\Delta}\arrow[dl] \\
& \Bun_G &
\end{tikzcd}\]
故 $\mL_{X^2}|_{X^2-\Delta}$ 在 $\Gr_X\times X|_{X^2-\Delta}$ 上的限制典范同构于 $\mL_X\boxtimes\mO_{X}|_{X^2-\Delta}$。利用 $\Pic^e(\Gr_G)\simeq \bZ$ 离散（故代数等价于零的线丛都平凡）以及 $\Gr_X$ 连通，可见对 $\Gr_X$ 的每个点 $s$，$\mL_{X^2}|_{X^2-\Delta}$ 在 $\Gr_X\times\{s\}$ 上的限制同构于 $\mL_X$；同理它在 $\{s\}\times \Gr_X$ 上的限制也同构于 $\mL_X$。

由于 $\Gr_X$ 在 $X$ 上归纳真，可用 See-Saw 原理（见 \cite[\S 5]{Mum}）推出 $\mL_{X^2}|_{X^2-\Delta}$ 典范同构于 $\mL_X\boxtimes\mL_X|_{X^2-\Delta}$，且与刚性化相容。
\end{proof}

注意若 $G$ 不单连通，命题 \ref{III:line} 对 $\Bun_G$ 上任意线丛不能成立\footnote{不过，若要求 $\Bun_G$ 上的 $\mL$ 在适当意义下「可乘」，命题大概可以推广到不必单连通的群。}。此外，该论证对非常值群概形 $\underline G$ 也不适用，即使其一般纤维单连通，因为我们关键用到了 $\Gr_{X^I}\to X^I$ 归纳真（以便应用 See-Saw 原理）。多数应用可用下文命题 \ref{det factorization} 代替。

确切地说，\S\ref{determinant line} 中 $\Gr_{\GL_n}$ 上行列式线丛的构造立刻推广，给出 $\Gr_{\GL_n,X^I}$ 上的行列式线丛 $\mL_{\det, X^I}$。

\begin{prop}\label{det factorization}
对行列式线丛 $\mL_{\det,X^I}$，有与命题 \ref{III:line} 类似的同构。
\end{prop}

特别地，给定表示 $\rho:\underline G\to \GL_n$，沿 $f_\rho: \Gr_{\underline G,X^I}\to\Gr_{\GL_n,X^I}$ 拉回，可在 $\Gr_{\underline G,X^I}$ 上构造带命题 \ref{III:line} 那种同构的线丛。

\begin{proof}
第一个同构的存在性显然。第二个同构由典范同构
\begin{equation}\label{det functor}
\det(M\oplus N)\cong \det(M)\otimes_R \det(N)
\end{equation}
推出，其中 $M,N$ 是有限投射 $R$-模。
\end{proof}

\begin{rmk}\label{graded factorization}
然而命题 \ref{III:line} 中的同构 $\mL_{X^2}|_{X^2-\Delta}\cong c^*(\mL_X\boxtimes\mL_X|_{X^2-\Delta})$ 与命题 \ref{det factorization} 中相应同构有本质差别。注意 $\mL_{X^2}$ 与 $\mL_{\det,X^2}$ 都带有自然 $S_2$-作用：命题 \ref{III:line} 的同构是 $S_2$-等变的，而命题 \ref{det factorization} 的同构只差一个符号才 $S_2$-等变。符号出现是因为在 \eqref{det functor} 下，典范同构 $\det(M\oplus N)\cong\det(N\oplus M)$ 对应于典范同构 $\det(M)\otimes\det(N)\cong\det(N)\otimes\det(M)$ 差符号 $(-1)^{\on{rk}M\on{rk}N}$。这反映了：取行列式线是从有限投射 $R$-模范畴 $\on{Proj}_R$（张量积取直和）到\emph{分次}可逆 $R$-模范畴 $\Pic^\bZ_R$ 的对称幺半函子。进一步阐发见 \cite{OZ1}。
\end{rmk}

再列出 Beilinson--Drinfeld Grassmann 簇的两条基本性质。

\begin{prop}\label{III:property}~
\begin{enumerate}
\item $\Gr_{X^I}$ 上存在典范联络，即：设 $x,y:S\to X^I$ 满足 $x|_{S_{red}}=y|_{S_{red}}$，则有典范同构 $\Gr_{X^I}\times_{X^I,x}S=\Gr_{X^I}\times_{X^I,y}S$。
\item Beilinson--Drinfeld Grassmann 簇与 \'etale 基变换交换，即：设 $U\to X$ 是 \'etale 态射，则有典范同构 $\Gr_{\underline G,X^I}\times_{X^I}U^I\cong\Gr_{\underline G_U, U^I}$。
\end{enumerate}
\end{prop}

\begin{proof}
(1) 见 \cite[\S 2.8]{BD}；(2) 见 \cite[Lemma 3.2]{Z10}。
\end{proof}

本小节最后引入 Beilinson--Drinfeld Grassmann 簇的几种变体。首先是 $X^J$ 上的卷积 Grassmann 簇。

设 $\phi:J\twoheadrightarrow \{1,2,\ldots,n\}$ 是有限集的满射，$J_i=\phi^{-1}(i)$。对 $x\in X^J$，以 $x_i\in X^{J_i}$ 记相应分量。定义卷积 Grassmann 簇
\begin{equation}\label{III:conv}
\Gr_{\underline G, \phi}=\left\{(x,(\mE_{i},\beta_i)_{i=1,2,\ldots,n})\ \left|\ \begin{array}{l}
x\in X^J(R),\quad \mE_i \text{ 是 } X_R \text{ 上的 } \underline{G}\text{-主丛},\\
\beta_i: \mE_{i}|_{X_R-\Ga_{x_i}}\cong \mE_{i-1}|_{X_R-\Ga_{x_i}},\ \text{其中 } \mE_0=\mE^0
\end{array}\right.\right\}.
\end{equation}
它有时也记作 $\Gr_{X^{J_1}}\tilde\times\Gr_{X^{J_2}}\tilde\times\cdots\tilde\times\Gr_{X^{J_n}}$，因为有投影
\[\pr_i:\Gr_{X^{J_1}}\tilde\times\Gr_{X^{J_2}}\tilde\times\cdots\tilde\times\Gr_{X^{J_i}}\to \Gr_{X^{J_1}}\tilde\times\Gr_{X^{J_2}}\tilde\times\cdots\tilde\times\Gr_{X^{J_{i-1}}}\]
把 $\Gr_{X^{J_1}}\tilde\times\cdots\tilde\times\Gr_{X^{J_i}}$ 实现为 $\Gr_{X^{J_1}}\tilde\times\cdots\tilde\times\Gr_{X^{J_{i-1}}}$ 与 $\Gr_{X^{J_i}}$ 的扭曲积（定义见 \S\ref{Intro:torsor}）。确切地说，有 $(L^+\underline G)_{X^{J_n}}$-主丛
\begin{equation}\label{torsor E}
\bE\to \Gr_{X^{J_1}}\tilde\times\Gr_{X^{J_2}}\tilde\times\cdots\tilde\times\Gr_{X^{J_{n-1}}}\times X^{J_n},
\end{equation}
分类点 $(x,(\mE_{i},\beta_i)_{i=1,2,\ldots,n-1})\in \Gr_{X^{J_1}}\tilde\times\cdots\tilde\times\Gr_{X^{J_{n-1}}}$、点 $y\in X^{J_n}$，以及 $\mE_{n-1}$ 沿 $\hat{\Ga}_y$ 的平凡化；再用该主丛构造扭曲积。

还有卷积映射
\[m_{\phi}: \Gr_{\underline G, \phi}\to \Gr_{X^J},\]
它把 $(x,(\mE_{i},\beta_i)_{i=1,2,\ldots,n})$ 送到 $(x, \mE_n, \beta_1\cdots\beta_n)$。用类似命题 \ref{III:Gr2} 的论证可得

\begin{lem}\label{generic conv}
$m_\phi$ 在 $X^{(\phi)}$ 上限制为同构。
\end{lem}

其次是 Beilinson--Drinfeld Grassmann 簇的对称化版本。（相对）有效 Cartier 除子的定义见 \cite[Tag01WQ, Tag056P]{St}。回忆给定 $d$，第 $d$ 个对称积 $X^{(d)}=X^d/\!\!/S_d$ 分类 $X$ 上的有效 Cartier 除子，即对每个 $k$-代数 $R$，$X^{(d)}(R)$ 是 $X_R$ 上相对有效 Cartier 除子的集合。于是可如下定义对称化的 Beilinson--Drinfeld Grassmann 簇 $\Gr^{(d)}\to X^{(d)}$：
\begin{equation}\label{sym BD Grass}
\Gr^{(d)}=\left\{(D,\mE,\beta)\ \left|\ \begin{array}{l}
D\subset X_R \text{ 是 } d \text{ 次相对有效 Cartier 除子},\\
\mE \text{ 是 } X_R \text{ 上 } \underline{G}\text{-主丛},\ \beta: \mE|_{X_R-D}\cong \mE^0|_{X_R-D} \text{ 是平凡化}
\end{array}\right.\right\}.
\end{equation}
容易看出自然映射
\[\Gr_{X^d}\to \Gr^{(d)}\times_{X^{(d)}}X^d\]
是同构。给定余权 $\mu$，有闭子簇 $\Gr^{(d)}_{\leq \mu}$，且自然映射
\[\Gr_{\leq (\mu,\ldots,\mu)}\to\Gr^{(d)}_{\leq \mu}\times_{X^{(d)}}X^d\]
是同构。卷积映射诱导 $\Gr_{\leq \mu}\tilde\times\cdots\tilde\times\Gr_{\leq \mu}\to \Gr_{\leq \mu}^{(d)}$。

\begin{ex}\label{momega1BD}
设 $G=\GL_n$，$\omega_1=(1,0,\ldots,0)$。则
\[\Gr^{(m)}_{\leq \omega_1}=\left\{ \mE\subset \mO_{X_R}^n\mid \mO_{X_R}^n/\mE \text{ 是秩 } m \text{ 的投射 } R\text{-模} \right\}.\]
进一步设 $X=\bA^1=\Spec k[t]$。则例 \ref{momega1} 与 \ref{conv v.s. springer} 有整体对应物。首先，公式 \eqref{nil to Gr} 实际定义了
\[\gl_n\to \Gr_{\leq \omega_1}^{(n)}.\]
在整体情形构造 \eqref{sat and springer} 的对应物留作练习。此外，\eqref{sat and springer2} 的整体类比是
\[\begin{CD}
\widetilde{\gl}_n@>>> \Gr_{\omega_1,X}\tilde\times\cdots\tilde\times\Gr_{\omega_1,X}@>>> X^n\\
@VVV@VVV@VVV\\
\gl_n@>>>\Gr^{(n)}_{\leq \omega_1}@>>>X^{(n)},
\end{CD}\]
其中 $\widetilde{\gl}_n\to\gl_n$ 是 Grothendieck--Springer 消解（\cite[\S~1]{Yu}），左侧方块是笛卡尔方块。把 $X^n\cong \bA^n$ 与 $\gl_n$ 的标准 Cartan 子代数 $\frakt_n$ 等同，把 $X^{(n)}$ 与 $\frakt_n/\!\!/S_n$ 等同，则下排两个映射的复合可等同于把矩阵送到其特征多项式的 Chevalley 映射。这个图最早出现在 \cite{MVy}。
\end{ex}

\begin{rmk}
$p$-进群的 Beilinson--Drinfeld 仿射 Grassmann 簇版本也存在，最近由 P. Scholze 构造 \cite{Sch}，但已超出本讲义范围。
\end{rmk}

% ============================================================
%  ch09  第三讲 §3.2 分解性质、§3.3 Ran 空间、§3.4 $\Gr_{\Ran}$ 上的刚性化线丛
%  原文行 2014–2194
% ============================================================

\subsection{分解性质}

下面用定理 \ref{III:factorization} 概括 Beilinson--Drinfeld Grassmann 簇 $\{\Gr_{X^I}\}$ 最重要的性质，它是命题 \ref{III:Gr2} 的极大推广，但证明方式相似。

\begin{thm}\label{III:factorization}
设 $\underline G$ 是 $X$ 上光滑仿射群概形。

(i) 对每个 $I$，有单位截面 $e_I: X^I\to \Gr_{X^I}$，由平凡 $\underline G$-主丛给出。

(ii) 对每个满射 $\phi: J\twoheadrightarrow I$，以 $\Delta(\phi):X^I\to X^J$ 记相应的对角嵌入，它把 $\{x_i, i\in I\}$ 送到 $\{y_j, j\in J\}$，其中 $y_j=x_{\phi(j)}$。则有典范同构
\begin{equation}\label{BD restriction}
\Delta(\phi):\Gr_{X^I}\cong \Gr_{X^J}\times_{X^J,\Delta(\phi)}X^I.
\end{equation}
为简记，复合 $\Gr_{X^I}\cong \Gr_{X^J}\times_{X^J,\Delta(\phi)}X^I\to \Gr_{X^J}$ 也记作 $\Delta(\phi)$。

(iii) 给定满射 $\phi: J\twoheadrightarrow  I$，以 $J=\sqcup_{i\in I}J_i$ 记 $J$ 的分拆，以 $X^{(\phi)}\subset X^J$ 记由那些满足「$\phi(j)\neq \phi(j')$ 时 $x_{j}\cap x_{j'}=\emptyset$」的 $\{x_j, j\in J\}$ 组成的开子集。则有典范同构（通常称为分解同构）
\begin{equation}\label{BD factorization}
c_\phi: \Gr_{X^J}\times_{X^J}X^{(\phi)}\cong \Big(\prod_{i\in I}\Gr_{X^{J_i}}\Big)\times_{X^J}X^{(\phi)}.
\end{equation}
这些映射满足下述条件：
\begin{enumerate}
\item $\Delta(\phi)\circ e_I=e_J|_{\Delta(\phi)}$，且 $c_\phi (e_J|_{X^{(\phi)}})=(\prod_{i\in I} e_{J_i})|_{X^{(\phi)}}$。
\item 给定 $K\stackrel{\phi}{\twoheadrightarrow}  J\stackrel{\psi}{\twoheadrightarrow}  I$，有 $\Delta(\phi)\Delta(\psi)=\Delta(\psi\phi)$。
\item 给定 $K\stackrel{\phi}{\twoheadrightarrow}  J\stackrel{\psi}{\twoheadrightarrow}  I$，下图交换
\[\begin{CD}
\Gr_{X^J}|_{X^{(\psi)}}@>\Delta(\phi)|_{X^{(\psi)}}>>\Gr_{X^K}|_{X^{(\psi\phi)}}\\
@Vc_{\psi} VV@VV c_{\psi\phi}V\\
(\prod_i\Gr_{X^{J_i}})|_{X^{(\psi)}}@>(\prod_{i}\Delta(\phi_i))|_{X^{(\psi)}}>> (\prod_i \Gr_{X^{K_i}})|_{X^{(\psi\phi)}},
\end{CD}\]
其中 $\phi_i:K_i\to J_i$ 是 $i\in I$ 上的映射。
\item 给定 $K\stackrel{\phi}{\twoheadrightarrow}  J\stackrel{\psi}{\twoheadrightarrow}  I$，有 $(\prod_{i\in I}c_{\phi_i})(c_{\psi\phi}|_{X^{(\phi)}})=c_{\phi}$。
\end{enumerate}
\end{thm}

\begin{ex}\label{FA ex}
看看这个定理如何包含命题 \ref{III:Gr2} 的全部信息。首先，把 \eqref{BD restriction} 用于映射 $\phi:\{1,2\}\to \{1\}$，便得同构 $\Delta: \Gr_X\cong \Gr_{X^2}|_{\Delta}$；把它用于交换 $\sigma:\{1,2\}\to \{1,2\}$，便得覆盖 $X^2\to X^2$ 的交换映射的典范同构 $\Delta(\sigma):\Gr_{X^2}\to \Gr_{X^2}$。另一方面，把 \eqref{BD factorization} 用于映射 $\id: \{1,2\}\to \{1,2\}$，便得同构 $c:\Gr_{X^2}|_{X^2-\Delta}\cong \Gr_X\times\Gr_X|_{X^2-\Delta}$。把性质 (3) 用于 $\{1,2\}\stackrel{\sigma}{\to}\{1,2\}\stackrel{\id}{\to}\{1,2\}$、把性质 (4) 用于 $\{1,2\}\stackrel{\id}{\to}\{1,2\}\stackrel{\sigma}{\to}\{1,2\}$，可见同构 $c$ 与 $\Delta(\sigma)$ 以及 $\Gr_{X}\times\Gr_X$ 的自然交换相容，这正是命题 \ref{III:Gr2} 所断言者。
\end{ex}

命题 \ref{III:line} 的推广是

\begin{thm}\label{III:factor line}
设 $G$ 单连通，$\underline G=G\times X$ 常值。设 $\mL$ 是 $\Bun_{G}$ 上刚性化线丛，$\mL_{X^I}$ 是它到 $\Gr_{X^I}$ 的拉回。则在定理 \ref{III:factorization} 中的同构下，有典范同构

(i) $e_I^*\mL_{X^I}\cong \mO_{X^I}$；(ii) $\Delta(\phi)^*\mL_{X^J}\cong \mL_{X^I}$；(iii) $\mL_{X^J}|_{X^{(\phi)}}\cong c_\phi^*(\boxtimes \mL_{X^{J_i}}|_{X^{(\phi)}})$。

此外，这些同构彼此相容，即定理 \ref{III:factorization} 中的 (1)--(4) 有相应的陈述。
\end{thm}

上述 Beilinson--Drinfeld Grassmann 簇的结构与性质也为其他几何对象所共有，值得概括为下述定义。

\begin{dfn}\label{FA space}
$X$ 上的\textbf{分解空间} $\mY$ 是指：对每个非空有限集 $I$ 指定空间 $\mY_I\to X^I$，并给出一族如定理 \ref{III:factorization} (ii)--(iii) 的映射，满足该定理的相容条件 (2)--(4)。若还存在一族如定理 \ref{III:factorization} (i) 的映射并满足相容条件 (1)，则称该分解空间是\textbf{带单位的}。

带单位分解空间 $\mY$ 上的\textbf{线丛} $\mL$ 是指：对每个 $I$ 指定 $\mY_I$ 上的线丛 $\mL_I$，并给出一族如定理 \ref{III:factor line} (i)--(iii) 的同构，满足该定理中的自然相容条件。
\end{dfn}

\begin{rmk}
分解空间（相应地，带单位分解空间）在 \cite[\S\ 3.10.16]{BD0} 中称为手征半群（相应地，手征幺半群）。进一步讨论与例子见\emph{同上}。
\end{rmk}

\begin{rmk}\label{det factorization2}
也可以把命题 \ref{det factorization} 推广为：族 $\mL_{\det, X^I}$ 构成 $\Gr_{\GL_n, X^I}$ 上的分次分解线丛，它们之间有定理 \ref{III:factor line} (i)--(iii) 那样的一族同构，但这些同构之间的相容条件须作修改。更确切地说，定理 \ref{III:line} 中类似的陈述 (3) 须修改。
\end{rmk}

\subsection{Ran 空间}

把所有 $\{\Gr_{X^I}\}$ 组织成一个对象、并用它重述上述定理，在概念上更简洁。

\begin{dfn}
$X$ 的 \textbf{Ran 空间} $\Ran(X)$ 是这样的预层：它给每个 $R$ 配上 $X(R)$ 的有限非空子集构成的集合。
\end{dfn}

注意对任意有限非空集 $I$ 都有重言映射 $X^I\to \Ran(X)$，它诱导同构
\[\varinjlim_I X^I\to \Ran(X),\]
其中余极限取在所有有限非空集（态射为满射）的范畴上的预层范畴中。由于余极限所过的范畴并不滤过，$\Ran(X)$ 不是空间（虽然每个 $X^I$ 都是）。事实上，它连 \'etale 拓扑下的层都不是。

对 $x=\{x_i,i\in I\}\in \Ran(X)(R)$，以 $\Ga_x$ 记所有图形 $\Ga_{x_i}$ 之并（作为 $X_R$ 中的闭子概形）。

\begin{dfn}\label{dfn Ran Grass}
定义 \textbf{Ran Grassmann 簇} $\Gr_{\underline G,\Ran(X)}$ 为
\[\Gr_{\underline G,\Ran(X)}=\left\{(x,\mE,\beta)\ \left|\ \begin{array}{l}
x\in \Ran(X)(R),\quad \mE \text{ 是 } X_R \text{ 上的 } \underline{G}\text{-主丛},\\
\beta: \mE|_{X_R-\Ga_x}\cong \mE^0|_{X_R-\Ga_x} \text{ 是平凡化}
\end{array}\right.\right\}.\]
\end{dfn}

注意我们并不声称 $\Gr_{\underline G,\Ran(X)}$ 可表示，因为它\emph{不}可表示。不过有自然投影
\[q_{\Ran}: \Gr_{\underline G,\Ran(X)}\to \Ran(X)\]
与单位截面
\[e_{\Ran}:\Ran(X)\to  \Gr_{\underline G,\Ran(X)},\]
它们相对可表示。事实上
\[\Gr_{\underline G,X^I}=X^I\times_{\Ran(X),q_{\Ran}}\Gr_{\underline G,\Ran(X)}.\]
若曲线 $X$ 与群概形 $\underline G$ 在上下文中明确，为简洁把 $\Gr_{\underline G,\Ran(X)}$ 记作 $\Gr_{\Ran(X)}$ 或 $\Gr_{\Ran}$。

下面是编码 Beilinson--Drinfeld Grassmann 簇分解结构的简洁方式。首先，Ran 空间 $\Ran(X)$ 有由点集之并给出的半群结构
\[\on{union}: \Ran(X)\times \Ran(X)\to \Ran(X),\quad (\{x_i\},\{x_j\})\mapsto \{x_i,x_j\}.\]
以 $(\Ran\times \Ran)_{disj}\subset \Ran\times\Ran$ 记由满足 $\{x_i\}\cap \{x_j\}=\emptyset$ 的 $\{x_i\},\{x_j\}$ 组成的开子预层，类似地定义 $(\Ran\times \Ran\times \Ran)_{disj}$。于是定理 \ref{III:factorization} 可重述为

\begin{thm}\label{Ran Grass factorization}
有典范同构
\[a: \Gr_{\Ran}\times_{\Ran}(\Ran\times\Ran)_{disj}\simeq (\Gr_{\Ran}\times\Gr_{\Ran})|_{(\Ran\times\Ran)_{disj}},\]
它与单位截面 $e_{\Ran}:\Ran\to\Gr_{\Ran}$ 相容，并在 $(\Ran\times\Ran\times \Ran)_{disj}$ 上满足自然的余圈条件。
\end{thm}

验证定理 \ref{III:factorization} 与定理 \ref{Ran Grass factorization} 等价，留作练习。

\subsection{$\Gr_{\Ran}$ 上的刚性化线丛}\label{line bundle on BD Grass}

Ran Grassmann 簇 $\Gr_{\Ran}$ 是极引人入胜的几何对象。下一讲将说明它编码了 $X$ 上 $\underline G$-丛模空间的许多信息，因此研究 $\Gr_{\Ran}$ 的几何很重要。这里用一个例子说明方法：计算 $\Gr_{\underline{G}, \Ran(X)}$ 上刚性化线丛的 Picard 群胚，其中 $\underline{G}\to X$ 是 $X$ 上单连通单群概形。

对带截面 $e:S\to X$ 的预叠态射 $X\to S$，以 $\Pic^e(X/S)$ 记 $X$ 上沿线丛的 Picard 群胚，附带沿 $e$ 的平凡化。有下述定理。

\begin{thm}\label{Rigid line}
设 $G$ 单且单连通，$\underline G=G\times X$ 常值。则对任意点 $x\in X$，映射 $\Gr_x\to \Gr_{\Ran}$ 诱导同构 $\Pic^e(\Gr_{\Ran}/\Ran)\simeq \Pic^e(\Gr_x)$。
\end{thm}

注意
\[\Pic^e(\Gr_{\Ran}/\Ran)=\varinjlim\ \Pic^e(\Gr_{X^I}/X^I),\]
且 Picard 群胚 $\Pic^e(\Gr_{X^I}/X^I)$ 离散。此外它是 $(X^I)_{et}$ 上某个 \'etale 层的整体截面群
\[\mP ic^e(\Gr_{X^I}/X^I)(U)= \Pic^e((\Gr_{X^I})_U/U).\]
先研究 $\mP ic^e(\Gr_{X^I}/X^I)$。

\begin{lem}
设 $X$ 是光滑曲线，$\underline G$ 是 $X$ 上连通约化群概形，逐纤维连通且单连通。则有典范同构
\[\bZ_X\simeq \mP ic^e(\Gr_X/X),\]
其中 $\bZ_X$ 是 $X$ 上的常值层 $\bZ$。特别地，对每个点 $x\in X$，茎 $\mP ic^e(\Gr_X/X)_x=\Pic^e(\Gr_x)=\bZ$。
\end{lem}

\begin{proof}
先处理 $\underline G=G\times X$ 常值且存在 \'etale 态射 $X\to \bA^1$ 的情形。由于仿射 Grassmann 簇与 \'etale 基变换交换（见命题 \ref{III:property}），化归到 $X=\bA^1$。此时有典范同构 $\Gr_{\bA^1}\cong \Gr_0\times {\bA^1}$：因为 $\bA^1$ 上有整体坐标，$\Aut(D)$-主丛 $\widehat{\bA^1}\to \bA^1$ 平凡。引理随即成立。注意同构 $\bZ\cong \Pic^e(\Gr_{\bA^1})$ 是典范的，因为有丰沛生成元 $\mO(1)$。

一般地，存在 \'etale 映射 $f:U\to X$，使得 $\underline G\times_XU\cong G\times U$（对某个常值群 $G$），且存在 \'etale 映射 $U\to \bA^1$。同构 $\bZ_U\cong \mP ic^e(\Gr_U/U)$ 的典范性蕴含它下降为典范同构 $\bZ_X\to \mP ic^e(\Gr_X/X)$。
\end{proof}

在上述同构下，把 $n\in \bZ_X$ 对应的 $\Gr_X$ 上的线丛记作 $\mO(n)_X$。

关键一步是理解 $\mP ic^e(\Gr_{X^2}/X^2)$。注意由于 $\Gr_X\to X$ 归纳射影，由 See-Saw 原理（见命题 \ref{III:line} 的证明）有典范同构
\[ \mP ic^e(\Gr_X/X)\boxtimes \mP ic^e(\Gr_X/X)\cong\mP ic^e(\Gr_X\times\Gr_X/X\times X).\]
特别地
\[\mP ic^e(\Gr_{X^2}/X^2)|_{X^2-\Delta}\cong \mP ic^e(\Gr_X/X)\boxtimes \mP ic^e(\Gr_X/X)|_{X^2-\Delta}\cong(\bZ_X\boxtimes\bZ_X)|_{X^2-\Delta}.\]

\begin{lem}\label{Pic X2}
上述同构可延拓为正合列
\[0\to \mP ic^e(\Gr_{X^2}/X^2)\to \bZ_X\boxtimes\bZ_X\to \Delta_*\bZ_X\to 0, \]
其中映射 $\bZ_X\boxtimes\bZ_X\to \Delta_*\bZ_X$ 由 $(a,b)\mapsto a-b$ 给出。
\end{lem}

\begin{proof}
先证 $\mP ic^e(\Gr_{X^2}/X^2)\to  \bZ_X\boxtimes\bZ_X$ 单射，即它没有沿对角支撑的非零截面。设 $(x,x)$ 是对角上的点。如上，在 $X$ 上 \'etale 局部地可设 $\underline G=G\times X$ 常值，且 $\Gr_{X}\to X$ 同构于 $\Gr\times X\to X$。设 $V\to X\times X$ 是 $(x,x)$ 的 \'etale 开邻域，为简洁把 $V\cap\Delta$ 记作 $\Delta$。设 $\mL$ 是 $\Gr_V:=\Gr_{X^2}\times_{X^2}V$ 上的刚性化线丛，在 $\Gr_{V-\Delta}:=\Gr_V|_{V-\Delta}$ 上平凡。要证 $\mL$ 平凡，只需证它对每个 Schubert 簇 $\Gr_{\leq(\la,\mu)}|_V\subset \Gr_V$ 的限制平凡（定义见命题 \ref{III:closure}），为简洁记该限制为 $\mL$。

由刚性化线丛的定义，$\mL$ 带有平凡化 $\theta: \mO_{V}\cong e^*\mL$；又 $\mL|_{\Gr_{V-\Delta}}$ 平凡。由于沿映射 $\Gr_{\leq(\la,\mu)}|_{V-\Delta}\to V-\Delta$ 的拉回诱导正则函数空间的同构，存在唯一平凡化
\[s:  \mO_{\Gr_{\leq(\la,\mu)}|_{V-\Delta}}\cong \mL|_{V-\Delta},\]
使 $e^*(s)=\theta|_{V-\Delta}$。由命题 \ref{III:closure} 与定理 \ref{singularity}，$\Gr_{\leq(\la,\mu)}$ Cohen--Macaulay，故 $\mL((n-1)\Delta)$ 是 $\mL(n\Delta)$ 的子层。于是存在整数 $n$ 使
\[s\in \Gamma(\Gr_{\leq(\la,\mu)},\mL(n\Delta))\setminus \Gamma(\Gr_{\leq(\la,\mu)},\mL((n-1)\Delta)).\]
把 $s$ 看作 $\mL(n\Delta)$ 的正则截面。此时 $e^*(s)$ 与 $\mO_V\simeq e^*\mL\to e^*\mL(n\Delta)$ 作为 $(e^*\mL)(n\Delta)$ 的正则截面一致，特别地 $n\geq 0$。但若 $n>0$，则 $s$ 在 $e(\Delta)$ 处消失，从而必在整个 $\Gr_{\leq (\la,\mu)}|_\Delta$ 上消失（因为 $s$ 的零点集纯余维一且含于 $\Gr_{\leq (\la,\mu)}|_\Delta$，而后者由命题 \ref{III:closure} 仍是整概形）。这将蕴含 $s\in \Gamma(\Gr_{\leq(\la,\mu)},\mL((n-1)\Delta))$，矛盾。故 $n=0$。同理 $s$ 不沿 $e(\Delta)$ 消失，从而处处不消失，于是 $s$ 给出 $\mL$ 的平凡化。单射性证毕。

现设 $\mL$ 是 $\Gr_{X^2}$ 上的刚性化线丛，其在 $\Gr_{X^2}|_{X^2-\Delta}$ 上的限制为 $\mO(a)_X\boxtimes\mO(b)_X|_{X^2-\Delta}$。仍由 See-Saw 原理，它沿卷积映射 $m: \Gr_X\tilde\times\Gr_X\to \Gr_{X^2}$ 的拉回 $m^*\mL$ 形如 $\pr_1^*\mO(a)_X\otimes \mL'$，其中 $\mL'$ 是 $\Gr_X\tilde\times\Gr_X$ 上某个线丛，其在每个 $\pr_1^{-1}(y)(\simeq \Gr_X)$ 上的限制同构于 $\mO(b)_X$。在对角上的一个点处，整体卷积映射限制为局部卷积映射（见 \eqref{conv m0}）。于是简单计算 \eqref{conv m0} 之下 $\Gr$ 上线丛的拉回即可知 $a=b$。这表明复合 $\Pic^e(\Gr_{X^2}/X^2)\to \bZ_X\boxtimes\bZ_X\to \Delta_*\bZ_X$ 为零。

最后要证 $\Pic^e(\Gr_{X^2}/X^2)$ 满射到 $\ker(\bZ_X\boxtimes\bZ_X\to\Delta_*\bZ_X)$。同样只需考虑沿对角上的茎，故如前可过渡到 $(x,x)$ 的 \'etale 邻域，从而设 $X=\bA^1\subset \bP^1$，$\underline G=G\times X$ 常值。此时把 $\Bun_G(\bP^1)$ 上的线丛 $\mO(1)$（见引理 \ref{O1 on Bun}）拉回到 $\Gr_{X^2}$ 即可。
\end{proof}

同样的论证表明 $\mP ic^e(\Gr_{X^I}/X^I)$ 适合下述正合列
\[0\to \mP ic^e(\Gr_{X^I})\to \boxtimes_{i\in I}\bZ_X\to \bigoplus_{\phi:I\twoheadrightarrow J/\sim}\Delta(\phi)_*\boxtimes_{j\in J}\bZ_X.\]
这里 $|J|=|I|-1$，且 $\phi\sim \phi'$ 指存在同构 $J\cong J'$ 使 $\phi$ 与 $\phi'$ 相容。换言之，和取遍把 $I$ 分成 $|I|-1$ 个子集的所有分拆。映射 $d_{\phi}:\boxtimes_{i\in I}\bZ_X\to \Delta(\phi)_*\boxtimes_{j\in J}\bZ_X$ 如下给出：可写 $I=I'\cup\{1,2\}\to J=I'\cup\{1\}$，$d_\phi$ 在 $\boxtimes_{i\in I'}\bZ_X$-因子上为恒等，在对应于 $\{1,2\}\to\{1\}$ 的因子上为引理 \ref{Pic X2} 中的映射 $\bZ_X\boxtimes\bZ_X\to \Delta_*\bZ_X$。

\begin{cor}
对每个点 $x\in X$ 与每个满射 $\phi: J\twoheadrightarrow I$，映射 $\Pic^e(\Gr_X)\to \Pic^e(\Gr_x)$ 与 $\Pic^e(\Gr_{X^J}/X^J)\to\Pic^e(\Gr_{X^I}/X^I)$ 都是同构。
\end{cor}

定理 \ref{Rigid line} 随即成立。

% ============================================================
%  ch10  第四讲 §4.1 单点单值化、§4.2 线丛与共形块
%  原文行 2199–2292
% ============================================================

\section{第四讲：对 $G$-丛模空间的应用}\label{Lecture IV}

\subsection{单点单值化}\label{one point uniformization}

设 $X$ 是 $k$ 上的光滑射影曲线，$\underline G$ 是 $X$ 上逐纤维连通的平滑仿射群概形。则 $\underline G$-主丛的模空间 $\Bun_{\underline G}$ 是 $k$ 上局部有限型的代数叠（例如见 \cite[Proposition 1]{He}）。

\begin{ex}
设 $\underline G|_{X-\{x_1,\ldots,x_n\}}\cong G\times (X-\{x_1,\ldots,x_n\})$，且在每个 $x_i$ 处 $\underline G|_{D_{x_i}}$ 是适当的抛孔群概形，则 $\Bun_{\underline G}$ 就是 $G$ 的抛物丛的模叠（抛物结构放在 $\{x_1,\ldots,x_n\}$ 处）。
\end{ex}

由 $\Gr_{\underline G,x}$ 的定义（参见 \eqref{aff glob}），有自然态射
\begin{equation}\label{one pt uni map}
u_x: \Gr_{\underline G,x}\to \Bun_{\underline G}.
\end{equation}
设 $X^*=X-\{x\}$，以 $\underline G^{X^*}$ 记截面空间，即
\begin{equation}\label{Gout}
\underline G^{X^*}(R)=\Gamma(X_R^*,\underline G).
\end{equation}

\begin{lem}\label{rep of Gout}
空间 $\underline G^{X^*}$ 可被一个群归纳概形表示。
\end{lem}

\begin{proof}
仍可化归到 $\underline G=\GL_n$，此时容易（例如见 \cite[(1.9)]{BL}）。
\end{proof}

通过限制到 $D_x^*\subset X^*$，可把 $\underline G^{X^*}$ 看作 $L\underline G_x$ 的子群。于是 $\underline G^{X^*}$ 沿 $u_x$ 的纤维传递作用，因此 $u_x$ 分解为
\begin{equation}\label{1pt uniform}
[\underline G^{X^*}\backslash\Gr_{\underline G,x}]\to \Bun_{\underline G}.
\end{equation}

该映射未必满射，$G=\bG_m$ 的情形即可见（开曲线 $X^*$ 上的线丛未必平凡）。但若 $\underline G=G\otimes X$ 半单，则它满射。事实上还有下述更强的结果。

\begin{thm}
设 $\underline G$ 是 $X$ 上逐纤维连通的平滑仿射群概形，其在有理函数域 $k(X)$ 上的通用纤维半单。则 \eqref{1pt uniform} 是同构。
\end{thm}

在 $k$-点层面，它说的是
\[\Bun_{\underline G}(k)= \Gamma\backslash G(F_x)/G(\mO_x),\]
其中 $\Gamma=\underline G^{X^*}(k)=\Gamma(X^*,\underline G)$。这归功于 Harder。$\SL_n$ 情形归功于 Beauville--Laszlo（见 \cite{BL}）。对常值半单群概形，这归功于 Drinfeld--Simpson（见 \cite{DS}）；\emph{同上}证明了：若 $p\nmid|\pi_1(G)|$，则 $\underline G^{X^*}$-主丛 $\Gr_{\underline G}\to \Bun_{\underline G}$ 在 \'etale 局部平凡。一般情形由 Heinloth 证明（见 \cite{He}）。因为有这个定理，可把 \eqref{one pt uni map} 看作单值化映射。

\begin{ex}
用最简单的例子 $G=\SL_2$ 说明该定理。此时 $\Bun_G$ 分类 $X$ 上行列式平凡的二阶向量丛。下面解释为什么这样的丛 $\mE$ 在 $X^*$ 上都变平凡。首先存在（足够负的）线丛 $\mL$ 使有短正合列 $0\to \mL\to \mE\to \mL^{-1}\to 0$。由于 $X^*$ 仿射，$\on{H}^1(X^*,\mL^{-2})=0$，故 $\mE|_{X^*}\cong (\mL\oplus\mL^{-1})|_{X^*}$。其次，可取足够大的整数 $n$，以及 $\mL(nx)$ 的整体截面 $s_1$ 与 $\mL^{-1}(nx)$ 的整体截面 $s_2$，使 $s_1$ 与 $s_2$ 在 $X^*$ 上无公共零点。令 $s=(s_1,s_2)|_{X^*}:\mO_{X^*}\to (\mL\oplus\mL^{-1})|_{X^*}$。则 $s$ 在 $X^*$ 上处处非零，于是有 $0\to \mO_{X^*}\to (\mL\oplus\mL^{-1})|_{X^*}\to \mO_{X^*}\to 0$。再由 $\on{H}^1(X^*,\mO_{X^*})=0$，得 $\mE|_{X^*}\cong (\mL\oplus\mL^{-1})|_{X^*}\cong \mO_{X^*}^2$。
\end{ex}

为简单起见，本小节余下部分假设 $\underline G=G\times X$ 常值，并把 $\underline G^{X^*}$ 记作 $G^{X^*}$。

\begin{cor}
有典范同构 $c:\pi_0(\Bun_G)\cong \pi_1(G)$。特别地，若 $G$ 单连通，则 $\Bun_G$ 连通。
\end{cor}

\begin{ex}
当 $G=\GL_n$ 时，映射 $\Bun_G\to \pi_0(\Bun_G)\cong\pi_1(G)\cong \bZ$ 由给 $X$ 上的向量丛配上其 Chern 类给出。
\end{ex}

\begin{proof}
要点是证明 $G^{X^*}$ 含于 $LG$ 的单位连通分支，于是 $\pi_0(\Bun_G)\cong \pi_0(\Gr_{x})$，再用定理 \ref{components}。该断言之证明见 \cite[(1.2)]{BLS}。
\end{proof}

\begin{rmk}
(i) 注意 $G^{X^*}$ 本身未必连通。事实上 \cite{BLS} 证明了 $\pi_0(G^{X^*})\simeq H^1(X^*,\pi_1(G))$（在 $\cha\ k=0$ 时）。

(ii) 该结果对非常值群概形的推广见 \cite{He}。
\end{rmk}

\subsection{线丛与共形块}

单值化定理对研究 $\Bun_G$ 上的线丛尤为有用。下面是主要定理。

\begin{thm}\label{uniformize line}
设 $G$ 单连通。则拉回诱导同构 $u_x^*:\Pic^e(\Bun_G)\to \Pic^e(\Gr_x)$。
\end{thm}

特别地，若 $G$ 单，则 $\Pic^e(\Bun_G)\cong\bZ$；$\Gr_G$ 的丰沛生成元 $\mO(1)$ 下降为 $\Pic^e(\Bun_G)$ 的生成元，仍记作 $\mO(1)$。

该映射的单射性容易说明：由单值化定理，$\Gr_G$ 上的刚性化线丛等同于 $\Gr_G$ 上 $G^{X^*}$-等变线丛，因此该推论等价于「任何特征 $G^{X^*}\to \bG_m$ 都平凡」；后者见 \cite[Corollary 5.2]{LS}。然而由于 $G^{X^*}$ 的几何仍很复杂，从单点单值化定理推出定理的满射部分并不容易。对典型群与 $G_2$，定理 \ref{uniformize line} 在 \cite{LS} 中通过显式构造拉回为 $\mO(1)$ 的 $\Bun_G$ 上线丛而得证（类似于注 \ref{generator} 的讨论）。完整定理先后由 Sorger \cite{So0}（$\cha\ k=0$，用到共形块空间的非消没结果）与 Faltings \cite[Theorem 17]{Fa}（一般情形）证明。下一小节将从另一个单值化定理推出它。

现在设 $k=\bC$（或任意特征零的代数闭域），$\mL$ 是 $\Bun_G$ 上的线丛。$\mL$ 的整体截面通常称为\textbf{非阿贝尔 theta 函数}（因为 $G=\bG_m$ 时就是经典 theta 函数）。设 $\frakg^{X^*}=\frakg\otimes\mO_X(X^*)$，这是 $G^{X^*}$ 的李代数。

\begin{lem}
Kac--Moody 中心扩张（即 \eqref{central extension} 的李代数）在 $\frakg^{X^*}$ 上唯一分裂，因此 $\frakg^{X^*}$ 作为子代数嵌入 Kac--Moody 代数 $\frakg_\Ga$（参见 \S\ref{Kac-Moody}）。
\end{lem}

\begin{proof}
由注 \ref{2-cocycle} 的显式公式与留数定理，中心扩张在 $\frakg^{X^*}$ 上分裂；又因 $[\frakg,\frakg]=\frakg$，这样的分裂唯一。
\end{proof}

在物理中，可积表示 $L_m(\frakg_\Ga)$ 是 WZW 共形场论的真空态空间，而 $\Hom_{\frakg^{X^*}}(L_m(\frakg_\Ga),k)$ 称为\textbf{共形块}空间。结合上述单值化定理与 Borel--Weil 型定理 \ref{Borel-Weil}，得

\begin{cor}
设 $G$ 单且单连通。$\Bun_G$ 上 $m$ 级非阿贝尔 theta 函数空间同构于共形块的对偶，即
\[\Gamma(\Bun_G,\mO(m))^*\simeq L_m(\frakg_\Ga)/\frakg^{X^*} L_m(\frakg_\Ga),\]
其中 $\mO(m)=\mO(1)^{\otimes m}$。
\end{cor}

共形块空间的维数由所谓的 Verlinde 公式计算，由此也得到 $\Bun_G$ 上非阿贝尔 theta 函数空间的维数。例如

\begin{cor}
设 $G=\SL_n$。则
\[\dim \Gamma(\Bun_G,\mO(m))= \Big(\frac{n}{n+m}\Big)^g \sum_{\substack{S\subset \{1,\ldots,m+n\}\\|S|=n}}\prod_{\substack{s\in S\\ t\not\in S}}\left|2\sin\pi\frac{s-t}{m+n}\right|^{g-1}.\]
\end{cor}

特别地，右端是正整数！见 \cite[Corollary 8.6]{BL}。

\begin{cor}
设 $G$ 是 $A,D,E$ 型，则 $\dim \Ga(\Bun_G,\mO(1))=|Z(G)|^g$，其中 $Z(G)$ 是 $G$ 的中心，$g$ 是 $X$ 的亏格。
\end{cor}

特别地，若 $G=E_8$，则 $\Bun_G$ 上定义 $\mO(1)$ 的除子唯一。在几何上找出这个除子将非常有趣（此问题由 Faltings 提出，见 \cite{Fa2}）。

单点单值化定理的进一步讨论与应用见 \cite{BL, LS, BLS, So0, So, He}，另见 \cite{KNR, KN}。

% ============================================================
%  ch11  第四讲 §4.3 adèle 单值化
%  原文行 2293–2457
% ============================================================

\subsection{adèle 单值化}

我们已看到用 $\Gr_x$ 单值化 $\Bun_{\underline G}$ 的一些重要应用。但在实践中，该单值化定理有几处局限。首先它只对半单群有效；更严重的是，$u_x$ 的纤维（它们是 $\underline G^{X^*}$-主丛）的几何仍然复杂，因此关于 $\Bun_{\underline G}$ 的某些问题难以通过研究 $\Gr_x$ 得到（如定理 \ref{uniformize line}）。本小节引入 $\Bun_{\underline G}$ 的另一种单值化来克服这些困难。

回忆 A. Weil 的下述结果。设 $k$ 代数闭，$X$ 是 $k$ 上光滑连通射影曲线，$|X|$ 记闭点集。设 $\bO=\prod_{x\in |X|} \mO_x$ 是整 adèle 环，$\bA=\prod' F_x$ 是 adèle 环，即 $F_x$（$x\in |X|$）关于所有闭点的\emph{限制}乘积：$\bA$ 中的元素是族 $(f_x\in F_x)_{x\in |X|}$，其中除有限多个 $x$ 外均有 $f_x\in \mO_x$。设 $\eta$ 是 $X$ 的通用点，$F_\eta$ 是 $X$ 的函数域。注意自然嵌入 $F_\eta\subset F_x$ 给出对角嵌入 $F_\eta\subset \bA$。

设 $\underline G$ 是 $X$ 上逐纤维连通的平滑仿射群概形，记 $G=\underline G\otimes F_\eta$。则

\begin{lem}
有自然双射
\begin{equation}\label{Weil}
\Bun_{\underline G}(k)\simeq G(F_\eta)\backslash G(\bA)/\underline G(\bO).
\end{equation}
\end{lem}

\begin{proof}
设 $\mE$ 是 $X$ 上的 $\underline G$-主丛。由于 $F_\eta$ 是上同调维数为一的域，由 Steinberg 定理，$G$-主丛 $\mE\otimes F_\eta$ 平凡。取这样的平凡化 $\beta_\eta: \mE\otimes F_\eta\simeq \mE^0\otimes F_\eta$，它可延拓为某个开子集 $U\subset X$ 上的平凡化。

另一方面，对每个 $x\in |X|$，$\mE|_{D_x}$ 也平凡，固定这样的平凡化 $\beta_x: \mE|_{D_x}\simeq \mE^0|_{D_x}$。则 $(\beta|_{D_x^*})\circ(\beta_x^{-1}|_{D_x^*})$ 是 $D_x^*$ 上平凡 $\underline G$-主丛的自同构，从而定义一个元素 $g_x\in G(F_x)$。注意对 $x\in |U|$ 有 $g_x\in G(\mO_x)$，于是由三元组 $(\mE,\beta_\eta,\{\beta_x\}_{x\in |X|})$ 构造出元素 $(g_x)\in G(\bA)$。$\beta_\eta$ 与 $\beta_x$ 的不同选取，会把该元素分别左乘 $G(F_\eta)$ 中元素、右乘 $G(\bO)$ 中元素加以改变。因此 $\mE$ 给出 $G(F_\eta)\backslash G(\bA)/\underline G(\bO)$ 中一个确定的元素，这就定义了映射 $\Bun_{\underline G}(k)\to G(F_\eta)\backslash G(\bA)/\underline G(\bO)$。

反之，给定 $g=(g_x)\in G(\bA)$，存在开子集 $U\subset X$ 使 $x\in |U|$ 时 $g_x\in G(\mO_x)$。于是可用 $g_x\in G(F_x)=\Aut(\mE^0|_{D_x^*})$ 作为过渡函数，把 $U$ 上与 $D_x$（$x\in |X|\setminus |U|$）上的平凡主丛粘起来，得到 $X$ 上的 $\mE$。该构造给出逆映射。
\end{proof}

下面描述 \eqref{Weil} 的代数几何版本，称之为 $\Bun_{\underline G}$ 的\textbf{adèle 单值化}。第一个问题是如何赋予集合 $G(\bA)/\underline G(\bO)$ 代数几何结构。事实上，这个问题有两种自然的解决方案。

第一种方案用 Ran Grassmann 簇 $\Gr_{\Ran}$（定义 \ref{dfn Ran Grass}）。$G(\bA)/\underline G(\bO)$ 中的元素可解释为一对 $(\mE,\beta_\eta)$，其中 $\mE$ 是 $X$ 上的 $\underline G$-主丛，$\beta_\eta$ 是 $\mE$ 在 $\eta$ 上（等价地，在某个开子集 $U\subset X$ 上）的截面；而 $\Gr_{\Ran}$ 的 $k$-点还指定了 $X$ 上的一组点，$\beta_\eta$ 在这组点上可能不再是正则映射。可按如下方式去掉这种冗余（沿 \cite[\S~2]{Ga3} 的做法）。

回忆 Ran 空间 $\Ran=\Ran(X)$ 有点集之并给出的半群结构。有「作用」映射
\[\on{act}: \Ran\times \Gr_{\Ran}\to \Gr_{\Ran},\]
把 $(x, (y,\mE,\beta))\in \Ran(R)\times\Gr_{\Ran}(R)$ 送到 $(x\cup y, \mE, \beta|_{X_R-\Ga_{x\cup y}})$。所谓「商」预层记作 $\Gr_{\Ran,\on{indep}}$，定义为下图的上等化子
\[\Ran\times\Gr_{\Ran}\substack{\on{act}\\  \xrightarrow{\hspace*{.6cm}} \\[-0.6ex]  \xrightarrow{\hspace*{.6cm}} \\ \pr_2 } \Gr_{\Ran}.\]
具体说来，$\Gr_{\Ran,\on{indep}}(R)$ 是 $\Gr_{\Ran}(R)$ 按如下等价关系的商：$(x,\mE,\beta)\sim (x',\mE',\beta')$，若存在 $X_R$ 上 $\underline G$-主丛的同构 $\al: \mE\cong \mE'$，使限制到 $X_R-\Ga_{x\cup x'}$ 时与 $\beta$、$\beta'$ 相容。由定义即得
\[\Gr_{\Ran,\on{indep}}(k)= G(\bA)/\underline G(\bO).\]
但一般地 $\Gr_{\Ran,\on{indep}}$ 不是空间，除非 $\underline G$ 是平凡群（此时 $\Gr_{\Ran,\on{indep}}$ 是一点）。所幸 $\Gr_{\Ran,\on{indep}}$ 的 fpqc 层化存在，并可显式描述如下。

称开子概形 $U\subset X_R$ \textbf{极稠密}，若它逐纤维稠密，即复合 $U\to X_R\to \Spec R$ 忠实平坦。注意极稠密开子集的有限交仍极稠密。

定义 $\Gr_{\underline G,\Rat(X)}$ 为给每个 $k$-代数 $R$ 配上下述群胚的预叠：对象是三元组 $(\mE,U,\beta)$，其中 $\mE$ 是 $X_R$ 上的 $\underline G$-主丛，$U\subset X_R$ 是极稠密开子概形，$\beta$ 是 $\mE$ 在 $U$ 上的平凡化；从 $(\mE,U,\beta)$ 到 $(\mE',U',\beta')$ 的态射是使 $\beta'\alpha|_{U\cap U'}=\beta|_{U\cap U'}$ 的同构 $\alpha:\mE\cong \mE'$。

若 $\underline G$ 与 $X$ 在上下文中明确，为简洁把 $\Gr_{\underline G,\Rat(X)}$ 记作 $\Gr_{\Rat}$。需要

\begin{lem}\label{rat sheaf}
上述群胚离散，且 $\Gr_{\Rat}$ 实际上是层。
\end{lem}

\begin{proof}
要证任何对象都没有非平凡自同构，即：设 $R$ 是 $k$-代数，$\mE$ 是 $X_R$ 上的 $\underline G$-主丛，$\al: \mE\to \mE$ 是自同构，若 $\al|_U=\id$，则 $\al=\id$。把 $\mE$ 嵌入 $X_R$ 上某个仿射空间，便化归为：若 $f$ 是 $X_R$ 上的正则函数，且在某个极稠密开子集上为零，则 $f=0$。注意这并不完全平凡，因为没有假设 $R$ 既约。它由下述引理推出。
\end{proof}

\begin{lem}\label{very dense vs Cartier}
设 $X$ 是拟射影簇，$Z\subset X$ 是有限（可为空）闭点集，看作闭子概形。设 $U\subset X_R$ 是开稠密子集，满足 (i) $Z_R\subset U$；(ii) $U\to \Spec R$ 满射。则存在 $\Spec R$ 由 $\Spec R_i$ 组成的开覆盖与相对 Cartier 除子 $D_i\subset X_{R_i}$，使 $Z_R\cap D_i=\emptyset$ 且 $X_{R_i}\setminus D_i\subset U\otimes_RR_i$。
\end{lem}

回忆在仿射空间 $X=\Spec A$ 上，Cartier 除子按定义是由某个\emph{非}零因子 $f\in A$ 的方程 $f=0$ 定义的闭子概形（参见 \cite[Tag01WQ]{St}）。于是由定义，局部化映射 $A\to A_f$ 单射，因此在 $X\setminus D$ 上为零的函数在整个 $X$ 上为零。

\begin{proof}
化归到如下情形：$X=\Spec A$ 仿射，$S=\Spec R$，$U=\Spec A\otimes_k R\setminus V(\fraka)$，其中 $\fraka$ 是 $A\otimes_k R$ 的理想，且 $V(\fraka)\cap Z=\emptyset$。

$U\to S$ 的满射性等价于：对 $R$ 的每个素理想 $\frakp$，令 $K=\on{Frac}(R/\frakp)$，有 $\on{Im}(\fraka\otimes_RK\to A\otimes_kK)\neq 0$。固定这样的 $\frakp$，并可自由缩小 $S$ 使其围绕 $x=V(\frakp)$。

把 $R$ 写成有限生成 $k$-代数之并，并缩小 $U$，可设 $R$ 有限生成且 $U\to \Spec R$ 的像含 $x$。取 $\on{Im}(\fraka\otimes_R R/\frakp\to A\otimes_k R/\frakp)$ 中这样的非零元 $\bar{f}$，使 $V(\bar{f})\cap Z_K=\emptyset$，并取 $\bar f$ 的提升 $f\in \fraka$。则由 \cite[Tag00ME]{St}，乘法映射 $f: A\otimes_k R_\frakp\to A\otimes_k R_\frakp$ 单射且商为 $R_\frakp$-平坦。因 $R$ noether，必要时缩小 $S$ 使其围绕 $x$，可设 $R\to R_\frakp$ 单射，于是 $f:A\otimes_k R\to A\otimes_k R$ 单射。再次必要时缩小 $S$，可设 $V(f)$ 不含竖直除子且 $V(f)\cap Z=\emptyset$。于是对 $R$ 的每个素理想 $\frakq$，$f: A\otimes_k R/\frakq\to A\otimes_k R/\frakq$ 单射。因此 $(A\otimes_k R/f)_\frakq$ 对每个 $R_\frakq$ 平坦，从而 $A\otimes_kR/f$ 在 $R$ 上平坦。则 $D=V(f)$ 即所求的相对 Cartier 除子。
\end{proof}

显然，映射 $(x,\mE,\beta)\mapsto (\mE,U=X_R-\Ga_x,\beta)$ 把 $\Gr_{\Ran,\on{indep}}$ 实现为 $\Gr_{\on{Rat}}$ 的子预层。

\begin{lem}\label{sheafify Ran}
空间 $\Gr_{\on{Rat}}$ 是 $\Gr_{\Ran,\on{indep}}$ 的 fpqc 层化。
\end{lem}

\begin{proof}
只需证明 fpqc 局部地，每个点 $\Spec R\to \Gr_{\on{Rat}}$ 都能提升为 $\Gr_{\Ran,\on{indep}}$ 的点。设 $(\mE,U, \beta)$ 代表 $\Gr_{\on{Rat}}$ 的 $R$-点。由引理 \ref{very dense vs Cartier}，在 $\Spec R$ 的 Zariski 局部可设 $U=X_R\setminus D$，其中 $D\subset X_R$ 是相对 Cartier 除子。回忆给每个 $R$ 配上 $X$ 上相对 Cartier 除子集的预层由 $X$ 的对称幂表示。另一方面，对称化映射 $X^d\to X^{(d)}$ 有限忠实平坦。因此经 $R$ 的有限忠实平坦基变换后可设 $U$ 是某个 $x:\Spec R\to X^d$ 的 $\Ga_x$ 的补集。
\end{proof}

其次需要 $G(F_\eta)$ 上的代数几何结构，同样有两种方案。这里只介绍 ${\underline G}_{\Rat}:={\underline G}_{\Rat(X)}$，即对应于 $\Gr_{\on{Rat}}$ 的那种。它给每个 $R$ 配上对 $(U,\al)$ 的群胚，其中 $U$ 是 $X_R$ 上的极稠密开子集，$\al\in \underline G(U)$。这仍是群层。下述引理容易。

\begin{lem}\label{generic dependence}
${\underline G}_{\Rat(X)}$ 只依赖 $\underline G\otimes F_\eta$，且 ${\underline G}_{\Rat(X)}(k)=G(F_\eta)$。
\end{lem}

现在可陈述 \eqref{Weil} 的代数几何版本。

\begin{thm}\label{ran uniformization of BunG}
设 $\underline G$ 的通用纤维约化。则态射 $u_{\Rat}:\Gr_{\Rat}\to\Bun_{\underline G}$ 是 ${\underline G}_{\Rat}$-主丛，从而诱导同构
\[[{\underline G}_{\Rat(X)}\backslash \Gr_{\Rat}]\cong \Bun_{\underline G}.\]
确切地说，对每个 $\Spec R\to \Bun_{\underline G}$，存在忠实平坦映射 $\Spec R'\to \Spec R$，使 $\Spec R'\times_{\Bun_{\underline G}}\Gr_{\Rat}$ 与 $\Spec R'\times {\underline G}_{\Rat}$ 在 ${\underline G}_{\Rat}$-等变意义下同构。
\end{thm}

\begin{proof}
基本要素与定理 \ref{uniformization for P1} 相同：$X_R$ 上每个 $\underline G$-主丛在忠实平坦基变换 $\Spec R'\to\Spec R$ 后，都在某个极稠密开子集上可约化为 Borel。为证此，由引理 \ref{generic dependence} 可修改 $\underline G$ 的某些纤维使其成为 $X$ 上的 Bruhat--Tits 群概形（定义见注 \ref{projectivity of BD}），再应用 \cite[Corollary 26]{He}。又因 $\bG_m$-与 $\bG_a$-主丛可 Zariski 局部平凡，定理得证。
\end{proof}

引入上述 $\Bun_{\underline G}$ 的 adèle 单值化定理，一个原因是 $u_{\Rat}$ 的纤维的几何出人意料地「简单」：群空间 ${\underline G}_{\Rat(X)}$ 在某种意义上表现得像一点。下面用定理来支持这一断言。回忆对 $k$-空间 $\mF$（或 $k$-预叠），可以谈论其正则函数环与线丛的 Picard 群胚（见 \eqref{Pic line}）；\cite[\S~2]{GL} 还解释了如何定义预叠的（$\ell$-进）（上）同调。

\begin{thm}
设 $R$ 是有限型 $k$-代数。
\begin{enumerate}
\item[(1)] $\underline G_{\Rat}\otimes R$ 上每个正则函数都是 $\Spec R$ 上某正则函数的拉回。
\item[(2)] $\underline G_{\Rat}\otimes R$ 上每个线丛都是 $\Spec R$ 上某线丛的拉回。
\item[(3)] 映射 $\underline G_{\Rat}\otimes R\to \Spec R$ 诱导（$\ell$-进）同调的同构。
\end{enumerate}
\end{thm}

(1) 与 (2) 归功于 Beilinson--Drinfeld（见 \cite[Proposition 4.3.13]{BD} 与 \cite[\S~4.9.1]{BD0}）；(3) 归功于 Gaitsgory--Lurie（见 \cite{Ga3} 与 \cite[Theorem 3.3.1]{GL}）。

\begin{proof}
(3) 参见 \cite{GL}；这里仅对 $G=\bG_m$ 或 $\bG_a$ 的情形解释 (1)，一般情形化归到该特殊情形（用引理 \ref{generic dependence}）。

取 $X$ 上的丰沛线丛 $\mL$。对 $n\gg 0$，令 $V_n=\Gamma(X,\mL^n)$，$V'_n= V_n\setminus \{0\}$。则有映射
\[p_n:V'_n\times V'_n\to (\bG_m)_{\on{rat}(X)},\quad  p_n:V_n\times V'_n\to (\bA^1)_{\on{rat}(X)},\]
由 $p_n(f,g)\mapsto f/g$ 给出。由引理 \ref{very dense vs Cartier}，每个映射 $S\to (\bG_m)_{\on{rat}(X)}$（或 $S\to (\bA^1)_{\on{rat}(X)}$）在 $S$ 的 Zariski 局部，对某个足够大的 $n$ 通过 $p_n$ 分解。

给定 $Y_{\on{rat}(X)}\otimes R$ 上的正则函数 $\varphi$，其拉回给出 $(V'_n\times V'_n)\otimes R$（或 $(V_n\times V'_n)\otimes R$）上的正则函数 $p_n^*\varphi$，它关于 $\bG_m$ 在 $V_n'\times V'_n$（或 $V_n\times V'_n$）上的显然作用不变。取 $n$ 足够大使 $\dim V_n>1$。则 $p_n^*\varphi$ 延拓为 $(V_n\times V_n)\otimes R$ 上 $\bG_m$-不变的正则函数，而它必来自 $R$。这就证明了 (1)。(2) 可类似论证。
\end{proof}

为使该定理更有用，需要下述配套陈述。按定义有映射序列
\[\Gr_{\Ran}\stackrel{w}{\to} \Gr_{\Rat}\stackrel{u_{\on{Rat}}}{\to} \Bun_{\underline G},\]
且 $u_{\Ran}=w\cdot u_{\on{Rat}}$。有

\begin{thm}\label{Ran=Rat}~
\begin{enumerate}
\item[(1)] 沿 $w$ 的拉回诱导正则函数空间的同构。
\item[(2)] 沿 $w$ 的拉回诱导刚性化线丛的 Picard 群胚的等价。
\item[(3)] 沿 $w$ 的推前诱导 $\ell$-进同调的同构。
\end{enumerate}
\end{thm}

注意即使 $\underline G$ 是平凡群，该定理也不平凡；此时把它陈述为下述命题。

\begin{prop}\label{Ran pt}~
\begin{enumerate}
\item[(1)] $\Ran X\otimes R$ 上每个正则函数都是 $\Spec R$ 上某正则函数的拉回。
\item[(2)] 自然映射 $\Ran X\otimes R\to \Spec R$ 诱导（$\ell$-进）同调的同构。
\end{enumerate}
\end{prop}

\begin{proof}
(2) 仍参见 \cite{Ga3} 与 \cite[\S~2]{GL}。为证 (1)，用下述引理。
\begin{lem}
设 $f\in R[[x,y,z]]$ 是系数在 $R$ 中的关于 $x,y,z$ 的幂级数。若 $f$ 关于 $x,y,z$ 的置换不变，且 $f(x,x,y)=f(x,y,y)$，则 $f(x,x,x)$ 为常数。
\end{lem}
该引理留作有趣的习题。把 $X^I\to \Ran(X)$ 分解为 $X^I\stackrel{\Delta}{\to} X^{I\sqcup I\sqcup I}\to \Ran(X)$，便知 $\Ran\otimes R$ 上每个函数拉回到 $X^I\otimes R$ 时都来自 $\Spec R$ 上的正则函数，于是 (1) 成立。
\end{proof}

下面解释定理 \ref{Ran=Rat} 的 (1)。回忆预叠上正则函数空间的定义（参见 \eqref{reg fun}）。由于 $\bG_a$ 是层，由引理 \ref{sheafify Ran}，
\[\Gamma(\Gr_{\Rat},\mO)=\Hom(\Gr_{\Rat},\bG_a)=\Hom(\Gr_{\Ran,\on{indep}},\bG_a)=\Gamma(\Gr_{\Ran,\on{indep}},\mO),\]
而按定义后者同构于下述序列的等化子
\[\Gamma(\Gr_{\Ran},\mO)\substack{\on{act}^*\\  \xrightarrow{\hspace*{.6cm}}\\[-0.6ex] \xrightarrow{\hspace*{.6cm}} \\ \pr_2^*} \Gamma(\Ran\times \Gr_{\Ran},\mO).\]
因此要证 (1)，只需说明 $\on{act}^*=\pr_2^*$。有自然映射 $(q_{\Ran},\id):\Gr_{\Ran}\to \Ran\times\Gr_{\Ran}$，把 $(x,\mE,\beta)$ 送到 $(x,(x,\mE,\beta))$；此外
\[ \on{act}\circ(q_{\Ran},\id)=\pr_2\circ(q_{\Ran},\id)=\id.\]
由命题 \ref{Ran pt} (1)，$\pr_2^*$ 是同构，故 $(q_{\Ran},\id)^*$ 是同构，从而 $\on{act}^*=\pr_2^*$。定理 \ref{Ran=Rat} 的 (2) 可类似证明。

\begin{cor}\label{Gr=Bun}~
\begin{enumerate}
\item 沿 $u_{\Ran}$ 的拉回诱导正则函数空间的同构。
\item 它诱导刚性化线丛的 Picard 群胚的等价。
\item 推前诱导 $\Gr_{\Ran}$ 与 $\Bun_{\underline G}$ 的同调的同构。
\end{enumerate}
\end{cor}

\begin{rmk}
(i) 注意把定理 \ref{Rigid line} 与推论 \ref{Gr=Bun} 的 (2) 结合，可得定理 \ref{uniformize line} 的自然证明。

(ii) 推论 \ref{Gr=Bun} 的 (3) 是 Gaitsgory--Lurie 证明整体函数域上 Weil 关于 Tamagawa 数猜想的关键步骤之一（参见 \cite{GL}）。
\end{rmk}

% ============================================================
%  ch12  第五讲 §5.1 Satake 范畴、§5.2 $\Sat_G$ 作为 Tannaka 范畴
%  原文行 2464–2643
% ============================================================

\section{几何 Satake 等价}\label{Lecture V}

本讲讨论几何 Satake 等价，它是几何朗兰兹纲领的基石之一。除另有声明外，本节假设 $k$ 代数闭。设 $\mO=k[[t]]$、$F=k((t))$，$G$ 是 $k$ 上的约化群，$\underline G= G\otimes_k\mO$。若群 $G$ 明确，为简洁把 $\Gr_G$ 记作 $\Gr$。

固定一个与 $k$ 的特征不同的素数 $\ell$。层指 $\Ql$-层，所有推前与拉回都是\emph{导出的}。需要指出，几何 Satake 等价还可在另外两种层范畴中表述：(1) 若 $\cha\ k=0$，可用 D-模（见 \cite[\S\ 5.3]{BD}）；(2) 若 $k=\bC$，可用 $\Gr$ 上解析拓扑下 $E$-向量空间层，其中 $E$ 是特征零的域。下面给出的途径对这两种范畴同样适用。另一方面，几何 Satake 等价还有 $\bZ_\ell$-与 $\bF_\ell$-层版本，或者当 $k=\bC$ 时 $\Lambda$-模层版本，其中 $\Lambda$ 是有限整体维数的交换 noether 环（例如 $\Lambda=\bZ$）。此时下述途径不够用，一般系数的推广需参见 \cite{MV}。

还要指出，$p$-进群也有几何 Satake 等价，细节见 \cite{Z14}。

\subsection{Satake 范畴 $\Sat_G$}\label{Satake category}

首先，$L^+G$ 在 $\Gr$ 上的作用满足 \S\ref{A:ind} 中的条件 (S)：每个轨道形如 $\Gr_\mu=L^+G/(L^+G\cap t^{\mu} L^+G t^{-\mu})$，且 $L^+G\cap t^{\mu} L^+G t^{-\mu}$ 连通。因此 \S\ref{A:ind} 的构造给出 $\Gr$ 上 $L^+G$-等变反常层范畴 $\on{P}_{L^+G}(\Gr)$。以 $\IC_\mu$ 记 $\Gr_{\leq \mu}$ 上的相交上同调层。则 $\IC_\mu|_{\Gr_\mu}=\Ql[(2\rho,\mu)]$，且它在每个层片 $\Gr_\la$ 上的限制是常值的。又因 $L^+G\cap t^{\mu} L^+G t^{-\mu}$ 连通，$\on{P}_{L^+G}(\Gr)$ 的不可约对象恰为这些 $\on{IC}_\mu$。\zy{反常层之名有误导：nLab 词条 perverse sheaf（\url{https://ncatlab.org/nlab/show/perverse+sheaf}，2026-10-16 修订）开篇即指出这类对象「既非 perverse 也非 sheaf」，而 $\on{P}(X)$ 是 $D^b_c(X)$ 的全 Abel 子范畴。$\IC_\mu$ 与相交上同调的关系见 nLab 词条 intersection cohomology（2023-05-07 修订）。[来源：nLab perverse sheaf、intersection cohomology。]}

\begin{prop}\label{semisimplicity}
范畴 $\on{P}_{L^+G}(\Gr)$ 半单。
\end{prop}

\begin{proof}
要证对任意 $\la,\mu$ 有 $\on{Ext}^1(\IC_\la,\IC_\mu)=0$。本质情形是 $\la\leq \mu$ 或 $\mu\leq \la$。此时消没由以下两点推出：(i)（仿射）Schubert 簇的 $\on{IC}$-层茎上同调的纯性（参见 \cite{KL}）；(ii) $\Gr$ 每个连通分支中 $L^+G$-轨道维数的奇偶性（见引理 \ref{parity}）。细节见 \cite[Proposition 1]{Ga}（另见 \cite[Proposition 3.1]{Ri}）。
\end{proof}

由此得到

\begin{cor}
$\on{P}_{L^+G}(\Gr)$ 中的对象都是 $\Aut^+(D)$-等变的。
\end{cor}

\begin{proof}
因为每个 IC 层都必须是 $\Aut^+(D)$-等变的（由引理 \ref{A:forgetful}）。
\end{proof}

$\on{P}_{L^+G}(\Gr)$ 上有自然的幺半结构，由 Lusztig 的层卷积积给出。即对 $\mA_1,\mA_2\in \on{P}_{L^+G}(\Gr)$，以 $\mA_1\tilde\boxtimes\mA_2$ 记 $\mA_1$ 与 $\mA_2$ 在 $\Gr\tilde\times\Gr$ 上的「外扭曲积」（见 \S\ref{A:twist product} 与 \S\ref{A:ind}）。例如若 $\mA_i=\IC_{\mu_i}$，则 $\IC_{\mu_1}\tilde\boxtimes\IC_{\mu_2}$ 典范同构于 $\Gr_{\leq \mu_1}\tilde\times\Gr_{\leq\mu_2}$ 的相交上同调层。于是 $\mA_1$ 与 $\mA_2$ 的卷积积定义为
\begin{equation}\label{Lus conv}
\mA_1\star\mA_2:=m_!(\mA_1\tilde\boxtimes\mA_2),
\end{equation}
其中 $m:\Gr\tilde\times\Gr\to \Gr$ 是卷积映射（由 \eqref{conv m0} 定义）。这是 $\Gr$ 上的 $L^+G$-等变 $\ell$-进复形。类似地，用 $n$ 重卷积映射 \eqref{conv m} 可定义 $n$ 重卷积积 $\mA_1\star\cdots\star\mA_n=m_!(\mA_1\tilde\boxtimes\cdots\tilde\boxtimes\mA_n)$。注意因 $m$ 归纳 proper，故 $m_!=m_*$。

\begin{prop}\label{conv prod}
卷积积 $\mA_1\star\mA_2$ 反常。
\end{prop}

\cite{BD} 称之为奇迹定理。它使我们能停留在 Abel 幺半范畴中，而不必处理更麻烦的三角幺半范畴概念。$\Sat_G$ 的另一个奇迹是引理 \ref{left=right}。

该命题本质上以两种不同途径证明。第一种基于仿射 Grassmann 簇的某种半无限几何（下面 \S\ref{MV semiinfinite} 将引入），Lusztig 用它证明了仿射 Hecke 代数的一个数值结果 \cite{Lu}，该结果与此命题等价（见 \cite[Proposition 2.2.1]{Gi1}）；Mirkovi\'c--Vilonen 用它证明了卷积映射的半小性，这与命题等价（见 \cite[Lemma 4.3, 4.4]{MV} 与下面注 \ref{semismall}）；\cite[Corollary 9.7]{NP}（另见 \cite[\S\ 2.2]{Z14}）也用它给出该命题的直接证明。

第二种途径基于 Beilinson--Drinfeld Grassmann 簇。Gaitsgory 借助邻近圈函子证明了更强的结果：在命题 \ref{conv prod} 中 $\mA_1$ 可以是 $\Gr$ 上任意反常层（注意定义外扭曲积并不需要 $\mA_1$ 的等变性条件）。见下面注 \ref{conv perverse}。

需要命题 \ref{conv prod} 的一个推论。设 $\mmu=(\mu_1,\ldots,\mu_n)$ 是支配余权序列。由命题 \ref{semisimplicity} 与 \ref{conv prod}，可写
\[\IC_{\mu_1}\star\cdots\star\IC_{\mu_n}=\bigoplus_{\la} V_{\mmu}^\la\otimes \IC_\la,\]
其中 $V_{\mmu}^\la=\Hom_{\on{P}(\Gr)}(\IC_\la, \IC_{\mu_1}\star\cdots\star\IC_{\mu_n})$。令 $|\mmu|=\sum \mu_i$。

\begin{cor}\label{Satake fiber}
有典范同构
\[V_{\mmu}^\la\cong \on{H}^{(2\rho,|\mmu|-\la)}_c(\Gr_{\mmu}\cap m^{-1}(t^\la),\Ql),\]
且后者有由 $\Gr_{\mmu}\cap m^{-1}(t^\la)$ 的维数为 $(2\rho,|\mmu|-\la)$ 的不可约分支给出的基。
\end{cor}

\begin{proof}
令 $d=\dim \Gr_{\mmu}\cap m^{-1}(t^\la)$。把 $\Gr_{\leq \mu_\bullet}$ 分层为 $\sqcup_{\mu'_\bullet}\Gr_{\mu'_\bullet}$，一个谱序列论证表明 $\IC_{\mu_1}\star\cdots\star\IC_{\mu_n}$ 在 $t^\la$ 处次数 $2d- (2\rho,|\mmu|)$ 的茎上同调为
\[\on{H}^{2d}_c(\Gr_{\mmu}\cap m^{-1}(t^\la),\Ql).\]
由 $\IC_{\mu_1}\star\cdots\star\IC_{\mu_n}$ 的反常性，$2d- (2\rho,|\mmu|)\leq -(2\rho,\la)$，等价于 $d\leq (\rho,|\mmu|-\la)$。又 $V_{\mmu}^\la$ 等于 $\IC_{\mmu}$ 在 $t^\la$ 处次数 $-(2\rho,\la)$ 的茎上同调。推论得证。
\end{proof}

\begin{rmk}\label{semismall}
(i) 注意该论证还表明
\[\dim  \Gr_{\leq \mmu}\cap m^{-1}(t^{\la})\leq (\rho,|\mmu|-\la).\]
于是卷积映射半小。但一般来说，$\Gr_{\leq \mmu}\cap m^{-1}(t^{\la})$ 可能存在维数严格小于 $(\rho,|\mmu|-\la)$ 的不可约分支，例子见 \cite[Remark 4.3, 8.3]{Ha}。

(ii) 这些 $\Gr_{\leq \mmu}\cap m^{-1}(t^{\la})$ 有时称为 \textbf{Satake 纤维}。当 $G=\GL_n$ 时，它们有时可等同于 Springer 纤维（见例 \ref{conv v.s. springer}）。特别地，这蕴含 Springer 消解 $\tilde\mN_n\to \mN_n$ 的半小性。
\end{rmk}

把 $(\mA_1\star\mA_2)\star\mA_3$ 与 $\mA_1\star(\mA_2\star\mA_3)$ 都等同于 $\mA_1\star\mA_2\star\mA_3$，就赋予卷积积自然结合约束。称幺半范畴 $(\on{P}_{L^+G}(\Gr),\star)$ 为 \textbf{Satake 范畴}，为简洁有时记作 $\Sat_G$。

\subsection{$\Sat_G$ 作为 Tannaka 范畴}\label{mon str of H}

本小节先按 \cite[\S\ 2.3]{Z14} 赋予超上同调函子 $\on{H}^*: \Sat_G\to \on{Vect}_{\Ql}$ 幺半结构，再说明如何把 $(\Sat_G,\on{H}^*)$ 上的幺半结构升级为对称幺半结构。真正的证明将基于后面对卷积积作为融合积的解释。\zy{这里用的是 Tannaka 对偶的「纤维函子」形式：由张量函子 $F:\mC\to\on{Vect}$ 的张量自同构群重建群概形，一般形式见 nLab 词条 Tannaka duality（\url{https://ncatlab.org/nlab/show/Tannaka+duality}，2025-07-28 修订）；本文把它用于 $\mC=\Sat_G$。[来源：nLab Tannaka duality，2025-07-28。]}

\begin{prop}\label{fiber functor}
超上同调函子 $\on{H}^*:\Sat_G\to \on{Vect}_{\Ql}$ 有典范幺半结构。
\end{prop}

想法如下。先设 $k=\bC$，从而可如 \S\ref{top loop group} 那样把 $\Gr$ 与 $\Gr\tilde\times\Gr$ 看作无穷维复解析空间。注意范畴 $\on{P}_{L^+G}(\Gr_G)$ 也可用解析拓扑定义。容易验证下述事实（见 \cite[Proposition 2.1.1]{Gi1}）。

\begin{lem}
在定理 \ref{top triv} 的同构下，$\Gr\tilde\times\Gr$ 上的 $\mA_1\tilde\boxtimes\mA_2$ 等同于 $\Omega K\times\Omega K$ 上的 $\mA_1\boxtimes\mA_2$。
\end{lem}

由 K\"unneth 公式可知 $\on{H}^*(-):=\on{H}^*(\Gr,-):\Sat_G\to \on{Vect}_{\Ql}$ 具有命题 \ref{fiber functor} 所述的幺半结构。

若 $k$ 一般，定理 \ref{top triv} 与上述引理并不字面成立。有两种途径使这些拓扑同构有意义。一种基于 Beilinson--Drinfeld Grassmann 簇把 $\Gr\tilde\times \Gr$ 形变到 $\Gr\times \Gr$，将在 \S\ \ref{fusion} 解释。这里解释另一种用等变上同调的途径，其优点是完全局部（即不依赖曲线的选取）。

回忆对 $\mA\in\Sat_G$，其 $L^+G$-等变上同调 $\on{H}^*_{L^+G}(\mA)$ 是自由 $R_{G}$-模（参见 \S\ref{A:char class}）。把 $\Gr$ 上 $L^+G$-等变层（非正式地）看作 $L^+G\backslash LG/L^+G$ 上的层，则 $\on{H}^*_{L^+G}(\mA)$ 上应还有另一个 $R_{G}$-模结构；它可以如下字面构造。以 $L^+G^{(m)}\subset L^+G$ 记第 $m$ 个同余子群（即投影 $L^+G\to L^mG$ 的核），以 $\Gr^{(m)}=LG/L^+G^{(m)}$ 记 $\Gr$ 上的通用 $L^m G$-主丛。则 $\Gr^{(m)}$ 带有 $L^+G\times L^mG$ 的作用，且投影
\[\pi_m:\Gr^{(m)}\to \Gr\]
是 $L^+G$-等变的。于是由 \eqref{coh free action}，
\[\on{H}^*_{L^+G}(\mA)= \on{H}^*_{L^+G\times L^mG}(\pi_m^*\mA),\]
从而给 $\on{H}^*_{L^+G}(\mA)$ 一个 $R_{G}$-双模结构。注意只要 $m>0$，该结构与 $m$ 无关。

回忆（含幺）代数 $R$ 的双模范畴有标准幺半结构 $(M,N)\mapsto M\otimes_RN$。

\begin{lem}\label{soergel monoidal}
对任意 $\mA_1,\mA_2,\ldots,\mA_n$，存在 $R_{G}$-双模的典范同构
\[\on{H}^*_{L^+G}(\mA_1\star\cdots\star\mA_n)\cong \on{H}^*_{L^+G}(\mA_1)\otimes_{R_{G}}\cdots\otimes_{R_{G}}\on{H}^*_{L^+G}(\mA_n),\]
并满足自然的相容条件。因此 $\on{H}^*_{L^+G}(-):\Sat_G\to (\on{R}_{G}\otimes\on{R}_{G})\on{-mod}$ 上有自然幺半结构。
\end{lem}

这是 Soergel 想法的一个变体：把（仿射）旗簇上的等变层对应到 Soergel 双模。这里只解释 $n=2$、$\mA_i=\IC_{\mu_i}$ 时如何构造映射，一般情形类似（例如见 \cite[\S~2.3]{Z14}）。设 $L^+G$ 在 $\Gr_{\leq \mu_2}$ 上的作用通过 $L^mG$ 分解，于是 $\Gr_{\leq \mu_1}\tilde\times\Gr_{\leq \mu_2}= \Gr_{\leq\mu_1}^{(m)}\times^{L^mG}\Gr_{\leq \mu_2}$。由卷积积的定义与 \eqref{coh free action}，
\[\on{H}^*_{L^+G}(\IC_{\mu_1}\star\IC_{\mu_2})=\on{H}^*_{L^+G\times L^mG}(\pi_m^*\IC_{\mu_1}\boxtimes \IC_{\mu_2}).\]
另一方面，由等变 K\"unneth 公式 \eqref{equiv Kunneth}，有典范同构
\[\on{H}^*_{L^+G}(\IC_{\mu_1})\otimes_{R_{G}} \on{H}^*_{L^+G}(\IC_{\mu_2})\cong \on{H}^*_{L^+G\times L^mG}(\pi_m^*\IC_{\mu_1}\boxtimes \IC_{\mu_2}).\]
合起来便得所需同构。

还需要一个事实。

\begin{lem}\label{left=right}
$\on{H}^*_{L^+G}(\mA)$ 上的两个 $R_{G}$-结构相同。
\end{lem}

\begin{proof}
这里只说明思路，细节见 \cite[\S\ 2.3]{Z14}。用引理 \ref{minusch} 与引理 \ref{quasiminusch}，容易对 $\mA=\IC_\mu$（$\mu$ 极小或拟极小）证明本引理。再由引理 \ref{soergel monoidal}，结论对这些 $\IC_\mu$ 的卷积成立。最后，每个 $\mA\in\Sat_G$ 都作为这些卷积的直和项出现（这是所谓 PRV 猜想的几何版本，见 \cite[Proposition 9.6]{NP}）。引理得证。
\end{proof}

\begin{rmk}
本引理是 $\underline G=G\otimes \mO$ 情形所特有的。例如它对仿射旗簇 $\Fl_G$ 上的等变反常层不成立（术语见注 \ref{terminology}）。
\end{rmk}

注意把命题 \ref{soergel monoidal} 与 \ref{left=right} 结合，实际上赋予等变超上同调函子
\[\on{H}^*_{L^+G}: \Sat_G\to \on{Proj}_{R_G}\]
典范幺半结构，其中如前 $\on{Proj}_{R_G}$ 记有限投射 $R_G$-模范畴，配以通常张量积。由 \eqref{de-equiv}，沿增广映射 $R_{G}\to \Ql$ 特殊化后，便得 $\on{H}^*$ 上所需的幺半结构。

留给读者验证：当 $k=\bC$ 时，上述 $\on{H}^*$ 上的幺半结构与用同胚 $\Gr\tilde\times\Gr\cong \Gr\times \Gr$ 构造者一致。

赋予 $\on{H}^*$ 幺半结构后，下面说明 $\Sat_G$ 上交换约束的构造，它使 $\on{H}^*$ 成为对称幺半函子。关键陈述如下。

\begin{prop}\label{existence of comm}
对任意 $\mA_1,\mA_2\in\Sat_G$，存在唯一同构 $c'_{\mA_1,\mA_2}:\mA_1\star\mA_2\cong\mA_2\star\mA_1$ 使下图交换
\[\begin{CD}
\on{H}^*(\mA_1\star\mA_2)@>\on{H}^*(c'_{\mA_1,\mA_2})>>\on{H}^*(\mA_2\star\mA_1)\\
@V\cong VV@VV\cong V\\
\on{H}^*(\mA_1)\otimes\on{H}^*(\mA_2)@>\cong>c_{\on{gr}}>\on{H}^*(\mA_2)\otimes\on{H}^*(\mA_1),
\end{CD}\]
其中竖直同构来自命题 \ref{monoidal functor}，而 $c_{\on{gr}}$ 是分次向量空间的交换约束，即
\[c_{\on{gr}}(v\otimes w)=(-1)^{\deg(v)\deg(w)}(w\otimes v).\]
\end{prop}

\begin{rmk}
读者应注意：命题 \ref{fiber functor} 给 $(\Sat_G, \on{H}^*)$ 添了额外结构，而命题 \ref{existence of comm} 建立的是 $(\Sat_G,\on{H}^*)$ 的某些性质。
\end{rmk}

因 $\on{H}^*:\Sat_G\to\on{Vect}_{\Ql}$ 忠实（由 $\Sat_G$ 半单），$c'_{\mA_1,\mA_2}$ 的唯一性显然，内容在其存在性。它将在把卷积积解释为融合积之后于 \S\ref{fusion} 证明。注意 \cite[\S~2.4]{Z14} 提示本命题可能存在纯局部（即不用整体曲线）的途径。

下面是所需的推论。首先可按 Koszul 符号改变来修改约束 $c'_{\mA_1,\mA_2}$（另见 \cite[\S~5.3.21]{BD}；更优雅的处理见 \cite{MV} 的注 6.2 之后）。即：范畴 $\on{P}_{L^+G}(\Gr)$ 带有由 \eqref{parity map} 诱导的 $\bZ/2$-分次。称 $\mA$ 具有纯奇偶性，若 $p(\on{Supp}(\mA))$ 为 $1$ 或 $-1$，此时定义 $p(\mA)=p(\on{Supp}(\mA))$，其中 $\on{Supp}(\mA)$ 记 $\mA$ 的支集。于是当 $\mA_1,\mA_2$ 具有纯奇偶性 $p(\mA_1),p(\mA_2)$ 时，令
\begin{equation}\label{koszul sign change}
c_{\mA_1,\mA_2}:=(-1)^{p(\mA_1)p(\mA_2)}c'_{\mA_1,\mA_2}.
\end{equation}
用修改后的交换约束，下图交换
\[\begin{CD}
\on{H}^*(\mA_1\star\mA_2)@>\on{H}^*(c_{\mA_1,\mA_2})>>\on{H}^*(\mA_2\star\mA_1)\\
@V\cong VV@VV\cong V\\
\on{H}^*(\mA_1)\otimes\on{H}^*(\mA_2)@>\cong>c_{\on{vect}}>\on{H}^*(\mA_2)\otimes\on{H}^*(\mA_1),
\end{CD}\]
其中同构 $c_{\on{vect}}$ 是向量空间的通常交换约束，即 $c_{\on{vect}}(v\otimes w)=w\otimes v$。

\begin{thm}\label{Tannakian structure}
幺半范畴 $\Sat_G$ 连同上述约束 $c_{\mA_1,\mA_2}$ 构成对称幺半范畴。超上同调函子 $\on{H}^*$ 是张量函子。事实上，$(\on{P}_{L^+G}(\Gr),\star,c,\on{H}^*)$ 构成中性 Tannaka 范畴。
\end{thm}

\begin{proof}
需验证 $c_{\mA_2,\mA_1}c_{\mA_1,\mA_2}=\id$ 与六边形公理。由 $\on{H}^*$ 的忠实性，取上同调后证明即可。由命题 \ref{existence of comm} 与 $c_{\on{vect}}^2=\id$ 知 $\on{H}^*(c_{\mA_2,\mA_1}c_{\mA_1,\mA_2})=\id$，故 $c_{\mA_2,\mA_1}c_{\mA_1,\mA_2}=\id$。六边形公理类似可证。第二个陈述显然。

为证 $(\on{P}_{L^+G}(\Gr),\star,c,\on{H}^*)$ 是中性 Tannaka 范畴，应用 \cite[Proposition 1.20, Remark 2.18]{DM}。即：仍由半单性，$\on{H}^*$ 正合（且忠实）。容易验证\emph{同上}命题 1.20 的条件 (a)--(d)（单位对象取 $\IC_0$）。至于条件 (e)，注意 $\dim \on{H}^*(\Gr,\IC_\mu)=1$ 当且仅当 $\dim \Gr_{\leq \mu}=0$，此时 $\mu$ 是中心 $Z_G\subset G$ 的余特征；于是 $\IC_{-\mu}\star\IC_\mu\cong \IC_0$。
\end{proof}

% ============================================================
%  ch13  第五讲 §5.3 Langlands 对偶群（半无限轨道与 MV 理论）
%  原文行 2644–2929
% ============================================================

\subsection{Langlands 对偶群}\label{MV semiinfinite}

前面已见 $\Sat_G$ 连同交换约束 $c$ 与纤维函子 $\on{H}^*$ 构成中性 Tannaka 范畴（在命题 \ref{existence of comm} 之后成立，该命题将在下一小节证明）。令
\[\tilde G=\Aut^\otimes \on{H}^*\]
为 Tannaka 群。下一步是用 $G$ 显式描述这个仿射群概形。首先容易证明

\begin{lem}
$\Ql$ 上的仿射群概形 $\tilde G$ 连通且约化。
\end{lem}

\begin{proof}
首先注意 $\IC_{\mu_1}\star\IC_{\mu_2}$ 以 $\IC_{\mu_1+\mu_2}$ 为直和项，因为卷积映射 $\Gr_{\leq\mu_1}\tilde\times\Gr_{\leq \mu_2}\to \Gr_{\leq \mu_1+\mu_2}$ 双有理。由此可知 $\Sat_G$ 作为张量范畴由有限多个 $\IC_\mu$ 生成（例如由 $\mu$ 为基本余权者生成），且不存在只含有限多个 $\IC$ 的直和的张量子范畴（例如 $\IC_\mu$ 的 $n$ 重自卷积含 $\IC_{n\mu}$ 为直和项）。于是由 \cite[Proposition 2.20, Corollary 2.22]{DM}，$\tilde G$ 代数且连通。最后，因 $\Sat_G$ 半单，$\tilde G$ 约化（\cite[Proposition 2.23]{DM}）。
\end{proof}

回忆对一个连通约化群 $G$，可附上四元组 $(\xch,\xcoch,\Phi,\Phi^\vee)$，称为 $G$ 的\textbf{根数据}（参见 \S\ref{Intro:group}）。令人惊讶的事实是：代数闭域上的连通约化群在同构意义下由这个组合数据唯一确定。此外，若 $(\xch,\xcoch,\Phi,\Phi^\vee)$ 来自某个 $G$，则 $(\xcoch,\xch,\Phi^\vee,\Phi)$ 也来自某个 $\hat G$：即 $\xcoch$（相应地 $\xch$）是 $\hat G$ 的 Cartan 环面 $\hat T$ 的权格（相应地余权格），而 $\Phi^\vee\subset \xcoch$（相应地 $\Phi\subset\xch$）是相应的根系（相应地余根系）。综述见 \cite[\S~1,2]{Sp}。约化群 $\hat G$ 通常称为 $G$ 的 \textbf{Langlands 对偶群}。\zy{nLab 词条 Langlands dual group（\url{https://ncatlab.org/nlab/show/Langlands+dual+group}，2014-09-07 修订）目前只列参考文献、未给定义，故不引为定义依据；本条依原文所引综述 \cite{Sp}（A 级）。[来源：A 级 \cite{Sp}；nLab 词条现状如上。]}

注意从 $G$ 到 $\hat G$ 的过渡完全是组合性的，依赖约化群的分类定理。下表给出若干例子。

\begin{center}
\begin{tabular}{| c | c | c | c | c | c |}
\hline 
$G$ & $\GL_n$ & $\on{SL}_n$ & $\on{SO}_{2n+1}$ & $\on{SO}_{2n}$ & $E_8$  \\ \hline 
$\hat G$ & $\GL_n$ & $\on{PGL}_{n}$ & $\on{Sp}_{2n}$ & $\on{SO}_{2n}$ & $E_8$ \\ \hline 
\end{tabular}
\end{center}

\begin{thm}\label{birth of dual group}
Tannaka 群 $\tilde G$ 的根数据为 $(\xcoch,\xch,\Phi^\vee,\Phi)$，因此同构于 $\hat G$。
\end{thm}

\begin{rmk}
在某种意义上，本定理可看作 $G$ 的 Langlands 对偶群的无分类定义。
\end{rmk}

本小节余下部分解释该定理的证明思路与一些应用。在证明过程中将给 $\tilde G$ 添上更多结构（见命题 \ref{canonical pinning}）。注意 Richarz 给出了本定理的另一途径（参见 \cite[\S~3]{Ri}）。

首先容易验证

\begin{lem}
设 $G=T$ 是环面。则定理成立。
\end{lem}

\begin{proof}
事实上，$(\Gr_T)_{\on{red}}$ 是典范同构于 $\xcoch(T)$ 的离散点集。于是容易看出 $\Sat_T$ 等价于 $\xcoch(T)$-分次有限维 $\Ql$-向量空间范畴，而 $\on{H}^*$ 就是遗忘分次的函子。引理得证。
\end{proof}

为从这一特殊情形推出一般情形，模仿经典构造定义「范畴 Satake 变换」，即对称幺半函子
\[\on{CT}:\Sat_G\to \Sat_T,\]
其中 $T$ 是 $G$ 的抽象 Cartan。

沿用 \S\ref{Intro:group} 的记号。固定 Borel 子群 $B\subset G$，令 $T=B/U$ 为抽象 Cartan，其中 $U$ 是 $B$ 的幂幺根。

考虑下述诱导图
\begin{equation}\label{CT diagram}
\Gr_{T}\stackrel{q}{\leftarrow} \Gr_{B}\stackrel{i}{\to} \Gr_{G}.
\end{equation}
注意 $\Gr_{B}\to \Gr_{G}$ 在点层面是双射，但远非同构于归纳概形；事实上态射 $i$ 本身就很病态。不过 $(\Gr_B)_{\on{red}}$ 的每个连通分支在 $\Gr_G$ 中局部闭。由于只关心这些空间上的 $\ell$-进层，可安全地把涉及的仿射 Grassmann 簇换为其约化子归纳概形（但为简化记号，下文省略下标 $_{\on{red}}$）。于是如前所述，$\Gr_T\simeq \xcoch(T)$ 是离散点集。对 $\la\in\xcoch(T)$，定义
\[S_\la= i(q^{-1}(\la))\subset \Gr_G.\]
注意它是 $LU$ 在 $\Gr_G$ 上作用（经由嵌入 $LU\subset LG$）下过 $t^\la$ 的轨道。

回忆 $\xcoch(T)$ 上的偏序 $\leq$。首先有

\begin{prop}\label{geom of Sla}
并 $S_{\leq \la}:=\cup_{\la'\leq \la} S_{\la'}$ 闭，且 $S_\la\subset S_{\leq \la}$ 开稠密。特别地，$S_\la$ 在 $\Gr_G$ 中局部闭，而 $S_{\leq \la}$ 是其闭包 $\bar{S}_\la$。
\end{prop}

\begin{proof}
证明与命题 \ref{HP} 及命题 \ref{geom of Sch} 类似。先证 $\cup_{\la'\leq \la} S_{\la'}$ 闭。先设 $G_\der$ 单连通。则 $S_{\leq \la}$ 上的点有如下模解释：对 $G$ 的最高权表示 $V_\chi$，以 $\ell_\chi$ 记相应的最高权直线，则
\[S_{\leq \la}=\left\{(\mE,\beta)\in\Gr\mid \beta^{-1}(\ell_\chi)\subset t^{-(\chi,\la)}(V_{\chi,\mE}),\ \mbox{对一切 } V_\chi\right\}.\]
于是如命题 \ref{HP} 的论证，可把这些条件写成一组矩阵方程，从而定义闭子集。对一般 $G$，可通过 $z$-扩张化归到 $G_\der$ 单连通的情形。

剩下要证：若 $\al$ 是正余根，则 $S_{\la-\al}$ 在闭包 $\overline{S}_\la$ 中。这由 \eqref{curve C} 中曲线 $C_{\la,\al}$ 的构造推出。
\end{proof}

这些 $LU$-轨道与 Schubert 胞腔、对偶 Schubert 胞腔之间有下述关系。

\begin{lem}\label{intersection of semi and Sch}
若交 $S_\la \cap \Gr_\mu$ 非空，则 $t^\la\in \Gr_{\leq \mu}$；若交 $S_\la\cap \Gr^{\mu}$ 非空，则 $t^\la\in \Gr^{\geq \mu}$。
\end{lem}

\begin{proof}
以 $2\rho^\vee$ 记 $G$ 的正余根之和。选取嵌入 $T\subset B$。由复合
\[\bG_m\stackrel{2\rho^\vee}{\to} T\subset G\subset L^+G,\]
$L^+G$ 在 $\Gr_G$ 上的作用诱导 $\bG_m$ 在 $\Gr_G$ 上的作用。则 $\Gr_B$ 成为该作用的吸引子，而 $\Gr_T\subset \Gr_G$ 是不动点集，即
\[S_\la=\left\{x\in \Gr_G\mid \lim\limits_{s\to 0}2\rho^\vee(s)x=t^\la\right\}.\]
于是因 $\Gr_\mu$ 是 $2\rho^\vee(\bG_m)$-不变的，$S_\la\cap \Gr_\mu\neq\emptyset$ 仅当 $t^\la\in \overline{\Gr_\mu}=\Gr_{\leq \mu}$。同一论证给出第二个断言。
\end{proof}

\begin{cor}\label{boundary of semi}
集合 $S_{<\la}:=\cup_{\la'<\la} S_{\la'}$ 是 $\overline{S}_\la$ 与丰沛线丛 $\mO(1)$ 的某个超平面截影之交。
\end{cor}

当 $\cha\ k=0$ 时，Mirkovi\'c--Vilonen 用相应 Kac--Moody 代数（参见 \S\ref{Kac-Moody}）的基本表示 $L_1(\frakg_\Ga)$ 证明了这一点（参见 \cite[Proposition 3.1(b)]{MV}）。下面给出一个不依赖 Kac--Moody 理论（且对任意特征有效）的论证。

\begin{proof}
经平移 $t^{-\la}$，可设 $\la=0$。容易看出 $S_0\subset \Gr^0$。由引理 \ref{intersection of semi and Sch}，$S_{<0}\cap \Gr^0=\emptyset$。换言之，$S_{<0}=S_{\leq 0}\cap \Theta$，其中由定理 \ref{Pic of Gr} 回忆 $\Theta$ 是 $\Gr_G$（的中性连通分支）上对应于 $\mO(1)$ 的有效 Cartier 除子。证毕。
\end{proof}

现在可以叙述 Mirkovi\'c--Vilonen 关于半无限轨道的主要几何结果（参见 \cite[\S\ 3]{MV}）。

\begin{thm}\label{MV theory}~
\begin{enumerate}
\item 若 $t^\la\in \Gr_{\leq \mu}$，则 $S_\la\cap \Gr_\mu$ 非空，且 $S_\la \cap \Gr_{\leq \mu}$ 纯维数为 $(\rho,\la+\mu)$。

\item 设 $\mA\in\Sat_G$。则 $(q_!i^*\mA)_\la=\on{H}_c^*(S_\la, \mA)$ 集中在次数 $(2\rho,\la)$。此外有典范同构
\[\on{H}_c^{(2\rho,\la)}(S_\la, \IC_\mu)\otimes\Ql((\rho,\la+\mu))\cong \Ql[\on{Irr}(S_\la\cap\Gr_{\leq \mu})],\]
其中右端是 $S_\la\cap\Gr_{\leq \mu}$ 的不可约分支张成的 $\Ql$-空间。

\item 有函子的典范同构
\[\on{H}^*(-)\cong \on{CT}:=\bigoplus_\la \on{H}_c^*(S_\la,-):\Sat_G\to \on{Vect}_{\Ql}.\]
\end{enumerate}
\end{thm}

\begin{rmk}
Mirkovi\'c--Vilonen 对一般系数环 $\Lambda$ 建立了该定理（当 $k=\bC$）。在这种一般性下，(1) 中的 $\IC_\mu$ 须换为相应的标准或余标准对象。
\end{rmk}

下面说明该定理证明的基本思路，细节见 \cite[\S\ 3]{MV}。\cite[\S~2.2]{Z14} 基于 \cite{NP} 的某些想法给出了另一途径。这一替代途径对建立 $p$-进群的几何 Satake 是必要的，因为（目前）我们在那种情形下有相应的推论 \ref{boundary of semi}（那里既没有 Kac--Moody 理论，也没有 Cartier 除子 $\Theta$）。

首先，由推论 \ref{boundary of semi} 可知：若 $\al$ 是单余根，则 $\dim(S_{\la-\al}\cap\Gr_{\leq \mu})=\dim(S_\la\cap\Gr_{\leq \mu})-1$。于是由下述起点归纳即得 (1)：
\begin{equation}\label{extreme semi}
S_\la\cap \Gr_{\leq \mu}=S_\la\cap \Gr_\mu=\left\{\begin{array}{ll}t^{w_0(\mu)} & \la=w_0(\mu), \\
                               L^+U t^\mu\cong \bA^{(2\rho,\mu)} & \la=\mu, \end{array}\right.
\end{equation}
其中 $w_0$ 是 $G$ 的 Weyl 群的最长元。

至于 (2)，注意函子 $q_!i^*$ 恰是 Braden 的双曲局部化函子（参见 \cite{Br}）。更确切地说，设 $B^-$ 是满足 $B\cap B^-=T$ 的对偶 Borel，则 $\Gr_{B^-}$ 是 $2\rho^\vee(\bG_m)$-作用的排斥子，且有图
\[\Gr_T\stackrel{q^-}{\leftarrow}\Gr_{B^-}\stackrel{i^-}{\to}\Gr_G.\]
由 Braden 定理，对 $\mA\in\Sat_G$ 有
\[q_!i^*\mA=(q^-)_*(i^-)^!\mA.\]
此外，用 $\{S_\la\cap\Gr_\mu, \mu\in\xcoch(T)^+\}$ 对 $S_\la$ 分层，诱导一个 $E_1$-项为 $\on{H}_c^*(S_\la\cap\Gr_\mu,\mA)$、收敛到 $\on{H}_c^*(S_\la,\mA)$ 的谱序列。这些考虑连同 (1) 的维数公式即得 (2)。

为证 (3)，用半无限轨道 $\{S_\la, \la\in\xcoch(T)\}$ 对 $\Gr$ 分层，诱导一个 $E_1$-项为 $\on{H}_c^*(S_\la,-)$、收敛到 $\on{H}^*(-)$ 的谱序列。由 (ii) 中建立的理由，它在 $E_1$-项退化。于是 $\on{H}^*$ 上有自然滤过，其伴随分次为 $\oplus_\la \on{H}_c^*(S_\la,-)$。具体地说，这是以 $(\xcoch(T),\leq)$ 为指标、定义为
\[\on{Fil}_{\geq \mu}\on{H}^*(\mA)=\ker(\on{H}^*(\mA)\to \on{H}^*(S_{<\la},\mA))\]
的滤过，其中 $S_{<\la}=\bar{S}_\la-S_\la$。另一方面，用
\[\left\{S^-_\la:=i^-(q^-(\la)),\ \la\in\xcoch(T)\right\}\]
对 $\Gr$ 分层给出 $\on{H}^*$ 的另一个滤过
\[\on{Fil}'_{<\la}\on{H}^*(\mA)=\on{Im}(\on{H}^*_{S^-_{<\la}}(\mA)\to \on{H}^*(\mA)),\]
其中 $\on{H}^*_{S^-_{<\la}}(\mA)$ 记 $\mA$ 到 $S^-_{<\la}$ 的 $!$-限制的上同调。这两个滤过彼此互补（由 Braden 定理与 (2)），共同给出分解 $\on{H}^*=\oplus_\la \on{H}_c^*(S_\la,-)$。

\begin{rmk}
(i) 由此可知 $\dim (S_\la\cap \Gr_{\leq \mu})=(\rho,\la+\mu)$ 与 $\dim \Gr_{\leq \mu}-\dim (S_\la\cap \Gr_{\leq \mu})=(\rho,\mu-\la)$ 都随 $\mu$ 趋于无穷而趋于无穷。换言之，$S_\la$ 在 $\Gr_G$ 中既无穷维又无穷余维。另一方面它是 $LU$-轨道，因此 $S_\la$ 通常称为\textbf{半无限轨道}。

(ii) 交 $S_\la\cap S^-_\mu$ 有限维，且只依赖 $\la-\mu$，故通常记作 $Z_{\la-\mu}$，称为 \textbf{Zastava 空间}。注意 $Z_{\la-\mu}$ 非空当且仅当 $\la\geq \mu$。

(iii) \cite[Theorem 3.6]{MV} 还证明了 $\on{H}^*$ 的滤过与分解不依赖嵌入 $B\subset G$ 与 $T\subset B$ 的选取。另见推论 \ref{canonical TB}。
\end{rmk}

可把 $\on{CT}$ 看作函子 $\Sat_G\to \Sat_T$，把定理 \ref{MV theory} (iv) 中的同构看作
\begin{equation}\label{geom sat tran}
\on{H}^*\circ \on{CT}\cong \on{H}^*: \Sat_G\to \on{Vect}_\Ql.
\end{equation}

\begin{prop}
$\on{CT}$ 上有唯一幺半结构使同构 \eqref{geom sat tran} 幺半。
\end{prop}

$\on{CT}$ 上的幺半结构若存在则显然唯一。其存在性由 \cite[Proposition 6.4]{MV} 用卷积积的融合积解释证明。这里用等变上同调勾勒另一途径。

\begin{proof}
存在性如下。前已说明 $\on{H}^*$ 上有滤过 $\on{Fil}^*_{\geq \mu}$ 使 $\on{CT}=\on{gr}\on{H}^*$。只需证明命题 \ref{fiber functor} 给出的幺半结构与该滤过相容。于是由 Verdier 对偶论证知互补滤过也与幺半结构相容；再由上述讨论中 $\on{H}^*\cong \on{CT}$ 的构造即得本命题。

显然的困难是：$\on{H}^*$ 的幺半结构是用 $L^+G$-等变上同调构造的，而 $S_\la$ 并非 $L^+G$-等变。不过 $S_\la$ 在常值环面 $T\subset L^+T\subset L^+G$ 的作用下稳定。于是分层 $\{S_\la\}$ 也诱导一个 $E_1$-项为 $\on{H}_{T,c}^*(S_\la,-)$、收敛到 $\on{H}^*_{T}$ 的谱序列；由 (ii) 的理由它在 $E_1$-项退化，故可把滤过 $\on{Fil}_{\geq \mu}\on{H}^*$ 提升为 $\on{H}^*_{T}$ 的由自由 $R_{T}$-模构成的滤过
\[\on{Fil}_{\geq \mu}\on{H}_{T}^*(\mA)=\ker(\on{H}_{T}^*(\mA)\to \on{H}_{T}^*(S_{<\la},\mA)),\]
使伴随分次为 $\oplus_\la\on{H}_{T,c}(S_\la,\mA)$（特别地是自由 $R_{T}$-模）。此外，沿增广映射 $R_{T}\to\Ql$ 基变换后即恢复 $\on{H}^*$ 上的滤过。细节见 \cite[\S\ 2.2]{YZ}。

因 $\on{H}_{T}^*(\mA)\cong \on{H}^*_{L^+G}(\mA)\otimes_{R_{G}}R_{T}$，命题 \ref{monoidal functor} 中构造的 $\on{H}_{L^+G}^*$ 上的幺半结构经基变换给出 $\on{H}_{T}^*$ 上的幺半结构。因此只需证明典范同构
\begin{equation}\label{equiv monoidal}
\on{H}^*_{T}(\mA_1)\otimes_{R_{T}}\on{H}^*_{T}(\mA_2)\cong \on{H}^*_{T}(\mA_1\star\mA_2)
\end{equation}
与滤过相容；而这一点只需在 $\Spec R_{T}$ 的通用点处验证。令 $\on{H}_\eta=\on{H}^*_{T}\otimes_{R_{T}}Q$，其中 $Q$ 是 $R_{T}$ 的分式域。应用等变局部化定理（参见定理 \ref{equiv localization}）得
\[\on{H}_\eta(\mA)\cong \bigoplus_\nu \on{H}_\eta(\mA|_{t^\nu}),\quad \on{H}_\eta(S_{<\la},\mA)\cong\bigoplus_{\nu<\la}\on{H}_\eta(\mA|_{t^\nu}).\]
于是 $\on{Fil}_{\geq \la}\on{H}_\eta(\mA)=\oplus_{\nu\geq \la}\on{H}_\eta(\mA|_{t^\nu})$。此外，由 $\on{H}^*_{T}$ 上幺半结构的定义可知，在这些等同下 \eqref{equiv monoidal} 由典范同构
\[\begin{split}\bigoplus_{\nu_1+\nu_2=\nu} \on{H}_\eta(\mA_1|_{t^{\nu_1}})\otimes_Q \on{H}_\eta(\mA_2|_{t^{\nu_2}})\cong &\bigoplus_{\nu_1+\nu_2=\nu}\on{H}_\eta((\mA_1|_{t^{\nu_1}})\otimes (\mA_2|_{t^{\nu_2}})) \\
\cong &\quad \on{H}_\eta(\mA_1\star\mA_2|_{t^\nu})\end{split}\]
诱导，其中第一个同构是等变 K\"unneth 公式 \eqref{equiv Kunneth}，第二个仍是局部化定理。合起来即得命题。
\end{proof}

\begin{rmk}\label{Parabolic MV}
Mirkovi\'c--Vilonen 的理论可推广到抛物子群。设 $P\subset G$ 是抛物子群，$L_P$ 是其 Levi 商，则有诱导图
\[\Gr_L\stackrel{q}{\leftarrow}\Gr_P\stackrel{i}{\to}\Gr_G.\]
给定 $\mA\in\Sat_G$，可定义 $q_!i^*\mA$ 为 $\Gr_L$ 上的层复形。通过选取 $G$ 的适当单参数子群，也可把它看作双曲局部化函子，因此 $q_!i^*\mA$ 反常（差若干平移）。于是有 $\on{CT}=\on{CT}^G_T$ 的抛物版本
\[\on{CT}^G_L: \Sat_G\to \Sat_L,\quad \mA\mapsto q_!i^*\mA[\mbox{平移}],\]
满足 $\on{CT}^G_T=\on{CT}^L_T\circ\on{CT}^G_L$。
\end{rmk}

应用命题 \ref{existence of comm} 可见 $\on{CT}$ 事实上尊重对称幺半结构，从而是两个 Tannaka 范畴之间的张量函子，因此诱导同态
\[\hat{T}\cong \tilde T\to \tilde G.\]
这给出 $\tilde G$ 的一个子环面。

\begin{lem}
它是极大环面。
\end{lem}

\begin{proof}
事实上，令 $L_\mu=\on{H}^*(\IC_\mu)$ 为 $\tilde G$ 的不可约表示。则在定理 \ref{MV theory} (3) 给出的分解下，$\on{H}_c^{(2\rho,\la)}(S_\la,\IC_\mu)$ 是 $L_\mu$ 的 $\la$-权空间。由 \eqref{extreme semi}，$\dim \on{H}_c^{(2\rho,\mu)}(S_\mu,\IC_\mu)=1$，这说明 $L_\mu$ 的所有特征线性无关。因此 $\tilde G$ 的表示环是 $\tilde T$ 的表示环的子环，从而 $\tilde T$ 极大。
\end{proof}

\begin{rmk}
注意 $T$（定义为抽象 Cartan）不是 $G$ 的子群，但它的 Langlands 对偶 $\hat T$ 典范地是 $\tilde G$ 的子群。
\end{rmk}

继续定理的证明。令 $\ell_\mu=\on{H}_c^{(2\rho,\mu)}(S_\mu,\IC_\mu)$。注意自然映射（差一个纯量）
\[L_{\mu_1}\otimes L_{\mu_2}\to L_{\mu_1+\mu_2}\]
把 $\ell_{\mu_1}\otimes\ell_{\mu_2}$ 送到 $\ell_{\mu_1+\mu_2}$，因为 $\on{H}^*$ 上的幺半结构与 $\on{H}^*$ 上的滤过 $\on{Fil}$ 相容。于是由 Pl\"ucker 关系，族 $\{\ell_\mu\subset L_\mu, \mu\in\xcoch(T)^+\}$ 定义一个 Borel 子群 $\tilde B\subset \tilde G$，且 $\tilde T\subset \tilde B$。按定义，$\mu$ 是表示 $L_\mu$ 的（关于 $\tilde B$ 的）最高权。因此支配余权集 $\xcoch(T)^+$ 与 $\tilde T$ 关于 $\tilde B$ 的支配权集 $\xch(\tilde T)^+$ 一致。

回忆由 $\tilde G$ 的正根生成的半群 $\tilde Q^+\subset \xch(\tilde T)$ 的下述描述：$\la\in \tilde Q^+$ 当且仅当存在最高权表示 $L_\mu$ 使 $\mu-\la$ 作为 $L_\mu$ 的权出现。由定理 \ref{MV theory}，这等价于存在某个 $\mu$ 使 $t^{\mu-\la}\in \Gr_{\leq \mu}$，又等价于 $\la$ 是 $G$ 的正余根之和。因此由 $G$ 的正余根生成的半群 $(Q^\vee)^+\subset\xcoch(T)$ 与半群 $\tilde Q^+$ 一致，进而 $G$ 的单余根集与 $\tilde G$ 的单根集一致。为完成定理 \ref{birth of dual group} 的证明，只剩下

\begin{lem}
根数据 $(\xch,\xcoch,\Phi,\Phi^\vee)$ 由支配权半群 ${\xch}^+$ 与单根集 $\Delta\subset\Phi$ 唯一确定。
\end{lem}

\begin{proof}
事实上，${\xch}^+$ 确定 $\xch\otimes \bR$ 的支配 Weyl 房，从而确定 Weyl 群 $W$；于是 $\Phi=W(\Delta)$。细节留作练习。
\end{proof}

上述论证还给出

\begin{cor}\label{canonical TB}
超上同调函子 $\on{H}^*$ 上的分次对应于单参数子群
\[2\rho:\bG_m\to \tilde T\subset \tilde G,\]
其中余特征 $2\rho\in \xcoch(\tilde T)=\xch(T)$ 是 $\tilde G$ 的正余根之和。

特别地，$\tilde T$ 是 $2\rho(\bG_m)$ 在 $\tilde G$ 中的中心化子，而 $\tilde B$ 是包含 $\tilde T$ 且使 $2\rho$ 关于之支配的唯一 Borel 子群。因此 $(\tilde T,\tilde B)$ 不依赖 $T\subset B$ 与 $B\subset G$ 的选取。
\end{cor}

本小节最后给出几何 Satake 等价的若干应用。首先它给出 $\hat G$ 的表示的某些「典范基」。

\begin{thm}\label{MV basis}~
\begin{enumerate}
\item 经由定理 \ref{MV theory} (2)，$S_\la\cap \Gr_{\leq \mu}$ 的不可约分支给出 $L_\mu$ 的 $\la$-权空间的一组自然基。

\item 经由推论 \eqref{Satake fiber}，$\Gr_{\leq \mmu}\cap m^{-1}(t_\la)$ 中维数为 $(\rho,|\mmu|-\la)$ 的不可约分支给出 $\Hom_G(V_\la, V_{\mu_1}\otimes\cdots\otimes V_{\mu_n})$ 的一组自然基。
\end{enumerate}
\end{thm}

(1) 中的基通常称为 Mirkovi\'c--Vilonen 基，$S_\la \cap \Gr_{\leq \mu}$ 的不可约分支通常称为 MV 圈；(2) 中的基有时称为 Satake 基。

下面给出几何 Satake 等价对拓扑的几个漂亮应用。

\begin{ex}
(i) 有典范同构
\[\on{H}^*(\Gr(i,n))\cong \wedge^i \Big(\bigoplus_{j=0}^{n-1}\Ql[-2j](-j)\Big).\]

(ii) 设 $X_n\subset \bP^{n+1}$ 是光滑 $n$ 维二次超曲面，则
\[\on{H}^*(X_n)\cong\left\{\begin{array}{ll}
\bigoplus_{j=0}^n \Ql[-2j](-j) &  n \mbox{ 为奇数,} \\ 
\bigoplus_{j=0}^n \Ql[-2j](-j)\bigoplus \Ql[-n](-\frac{n}{2}) &  n \mbox{ 为偶数.}
\end{array}\right.\]
证明是把 $\Gr(i,n)$ 与 $X_n$ 分别等同于 $\on{PGL}_n$、$\on{SO}_{n+2}$ 的极小 Schubert 簇。
\end{ex}

回忆若 $k=\bC$，则 $\Gr\cong \Omega K$，从而是 $H$-空间。因此上同调环 $\on{H}^*(\Gr)$ 有自然的 $\bZ$-分次交换且余交换的 Hopf 代数结构。$k$ 一般时这仍成立：卷积映射 $m:\Gr\tilde\times \Gr\to \Gr$ 诱导等变上同调的余乘法
\[\on{H}^*_{L^+G}(\Gr)\to \on{H}^*_{L^+G}(\Gr)\otimes_{R_G}\on{H}^*_{L^+G}(\Gr).\]
引理 \ref{left=right} 中的类似陈述对 $\on{H}^*_{L^+G}(\Gr)$ 成立。因此可沿 $R_G\to\Ql$ 特殊化，得到通常上同调的余乘法。

由这类 Hopf 代数的分类，$\on{H}^*(\Gr)=U(\fraka)$（或当 $k=\bC$ 时 $\on{H}^*(\Omega K)=U(\fraka)$），其中 $\fraka$ 是 $\on{H}^*(\Gr_G)$ 的本原部分。因 $\on{H}^*(\Gr)$ 通过杯积作用于 $\on{H}^*(\mA)$，一般 Tannaka 形式化诱导映射 $\fraka\to \hat \frakg=\Lie \hat G$。

\begin{thm}\label{canonical pinning}
固定同构 $\Ql(1)\simeq \Ql$。
\begin{enumerate}
\item Chern 类 $c_1(\mO(1))\in \fraka$ 及其在 $\hat \frakg$ 中的像，是 $\hat \frakg$ 的主幂幺元 $N$。

\item 映射 $\fraka\to\hat\frakg$ 单射，并把 $\fraka$ 与 $N$ 在 $\hat\frakg$ 中的中心化子等同。

\item $(\tilde G,\tilde B,\tilde T, N)$ 是 $\tilde G$ 的典范带钉。
\end{enumerate}
\end{thm}

(1) 与 (2) 在 \cite[\S 4]{Gi1} 与 \cite{YZ} 中证明。这里只指出 $N$ 是主幂幺元是相交上同调的硬 Lefschetz 定理的推论（参见 \cite[Lemma 4.4.2]{Gi1}）。(3) 在 \cite{YZ} 中证明。该定理特别给出当 $k=\bC$ 时紧 Lie 群的基于环路群的描述。

\begin{rmk}
(i) 本定理有整系数版本，用 $\Gr$ 的同调表述（\cite[Theorem 1.1]{YZ}），证明依赖几何 Satake 等价的整系数版本 \cite{MV}。

(ii) 类似地可用 $\hat\frakg$ 的正则中心化子描述等变上同调 $\on{H}^*_{L^+G}(\Gr)$，见 \cite[Theorem 2.12]{BFM} 与 \cite[Proposition 6.6]{YZ}。这是研究 $L^+G$-等变导出范畴 $D_{L^+G}(\Gr)$（又称导出 Satake 范畴）的第一步，见 \cite{BF}。
\end{rmk}

% ============================================================
%  ch14  第五讲 §5.4 融合积
%  原文行 2930–3074
% ============================================================

\subsection{融合积}\label{fusion}

Beilinson--Drinfeld Grassmann 簇的分解结构使我们能把卷积积 \eqref{Lus conv} 解释为融合积，由此可一并推出命题 \ref{fiber functor} 与命题 \ref{existence of comm}。

首先，用定义 $\on{P}_{L^+G}(\Gr_G)$ 的同一套手续（见 \S\ref{A:ind}），得到 $\Gr_{X^I}$ 上 $(L^+G)_{X^I}$-等变反常层的 Abel 范畴 $\on{P}_{(L^+G)_{X^I}}(\Gr_{X^I})$。仿照 Reich（见 \cite{Rei}），定义 $\Sat_{X^I}\subset \on{P}_{(L^+G)_{X^I}}(\Gr_{X^I})$ 为那些关于态射 $q_I:\Gr_{X^I}\to X^I$ 满足 ULA 条件的层构成的全子范畴（ULA 性质的综述见 \S\ref{ULA sheaf}）。由定理 \ref{III:factorization} 与定理 \ref{ULA perverse}，对每个 $\phi: J\to I$ 有函子
\begin{equation}\label{Sat restriction}
\Delta(\phi)^\bullet:=(\nu^\phi)^*[|I|-|J|]\cong(\nu^\phi)^![|J|-|I|]:\Sat_{X^J}\to\Sat_{X^I}.
\end{equation}

给定满射 $\phi:J\twoheadrightarrow I=\{1,2,\ldots,n\}$，有从卷积 Grassmann 簇到 Beilinson--Drinfeld Grassmann 簇的卷积映射
\[m_\phi:\Gr_{X^{J_1}}\tilde\times\cdots\tilde\times\Gr_{X^{J_n}}\to \Gr_{X^J}.\]
对 $\mA_i\in\Sat_{X^{J_i}}$，可用 \eqref{torsor E} 在 $\Gr_{X^{J_1}}\tilde\times\cdots\tilde\times\Gr_{X^{J_n}}$ 上构造「外扭曲积」$\mA_1\tilde\boxtimes\cdots\tilde\boxtimes\mA_n$。然后可定义「外卷积积」
\[\mA_1\boxstar \mA_2\boxstar \cdots\boxstar \mA_n:=(m_\phi)_!(\mA_1\tilde\boxtimes\cdots\tilde\boxtimes\mA_n).\]
关键陈述如下。

\begin{prop}\label{conv=fusion}~
\begin{enumerate}
\item 复形 $\mA_1\boxstar\mA_2\boxstar\cdots\boxstar\mA_n$ 属于 $\Sat_{X^{J}}$，因此有定义良好的外卷积积函子
\[\boxstar_{i\in I}:\Sat_{X^{J_1}}\times\cdots\times\Sat_{X^{J_n}}\to \Sat_{X^J},\quad (\mA_1,\ldots,\mA_n)\mapsto \mA_1\boxstar\cdots\boxstar\mA_n.\]

\item 有典范同构
\begin{equation}\label{eqfusion}
\mA_1\boxstar\cdots\boxstar\mA_n\cong j_{!*}((\mA_1\boxtimes\cdots\boxtimes\mA_n)|_{X^{(\phi)}}),
\end{equation}
其中 $j: X^{(\phi)}\to X^J$ 记开嵌入，而 $(\mA_1\boxtimes\cdots\boxtimes\mA_n)|_{X^{(\phi)}}$ 经同构 \eqref{BD factorization} 看作 $\Gr_{X^J}\times_{X^J}X^{(\phi)}$ 上的反常层。
\end{enumerate}
\end{prop}

\begin{proof}
首先注意若所有 $\mA_i$ 都 ULA，则 $\mA_1\tilde\boxtimes\cdots\tilde\boxtimes\mA_n$ 也 ULA。这是因为（在光滑拓扑下）局部上支集 $\on{Supp}(\mA_1\tilde\times\cdots\tilde\times\mA_n)$ 同构于乘积 $\on{Supp}(\mA_1)\times\cdots\times\on{Supp}(\mA_n)$。于是由定理 \ref{ULA prop} (3) 与 (4)' 即得。

再用定理 \ref{ULA prop} (2) 可知 $\mA_1\boxstar\cdots\boxstar\mA_n$ 关于 $\Gr_{X^J}\to X^J$ 是 ULA。此外它在 $X^{(\phi)}$ 上的限制反常，因为由引理 \ref{generic conv}，在同构 \eqref{BD factorization} 下
\[\mA_1\boxstar \mA_2\boxstar \cdots\boxstar \mA_n|_{X^{(\phi)}}=\mA_1\boxtimes\cdots\boxtimes\mA_n|_{X^{(\phi)}}.\]
故由定理 \ref{ULA perverse}，它反常，且 $\mA_1\boxstar\cdots\boxstar\mA_n\cong j_{!*}((\mA_1\boxtimes\cdots\boxtimes\mA_n)|_{X^{(\phi)}})$。最后，$(L^+G)_{X^J}$-等变性显然。
\end{proof}

下面推出若干推论（即下文 (i)--(iii)），特别是给出命题 \ref{fiber functor} 的另一证明与命题 \ref{existence of comm} 的证明。

(i) 结合 \eqref{Sat restriction}，可在每个 $\Sat_{X^J}$ 上定义内卷积积函子
\begin{equation}\label{internal convolution}
\ostar_{i\in I}: \times_{i\in I}\Sat_{X^J}\to \Sat_{X^J},\quad \mA_1\ostar\mA_2\ostar\cdots\ostar\mA_n:= \Delta^\bullet(\mA_1\boxstar\mA_2\boxstar\cdots\boxstar\mA_n),
\end{equation}
其中 $\Delta: X^I\to (X^I)^n$ 是对角嵌入，$\Delta^\bullet=\Delta^*[-(n-1)|I|]$。与 $\Sat_G$ 同样的理由表明每个 $\Sat_{X^J}$ 都是幺半范畴。

\begin{rmk}\label{triv group}
可以试着用同一公式 \eqref{internal convolution} 在整 $\on{P}_{(L^+G)_{X^J}}(\Gr_{X^J})$ 上定义内卷积积 $\ostar$。但下例说明 $\ostar$ 在这个更大的范畴上并不反常正合，这也是我们要求 ULA 条件的原因之一。

设 $G=\{e\}$ 平凡。则 $\on{P}_{L^+G}(\Gr_X)$ 就是 $X$ 上反常层的范畴，而 $\Sat_X$ 是 $X$ 上局部系（平移 $-1$ 度）构成的范畴。按定义
\[\mA_1\boxstar\mA_2=\mA_1\boxtimes\mA_2,\quad \mA_1\ostar\mA_2=\mA_1\otimes\mA_2[-1].\]
特别地，若 $\mA_1$ 或 $\mA_2$ 是摩天楼层，则 $\mA_1\ostar\mA_2$ 不反常。
\end{rmk}

注意有自然函子
\[\Sat_G\to \Sat_X,\quad \mA\mapsto \mA_X:=\Ql\tilde\boxtimes\mA[1],\]
其中 $\mA_X=\Ql\tilde\boxtimes\mA[1]$ 是 $\Gr_X=X\tilde\times\Gr$ 上的外扭曲积（见 \eqref{III: twisted constr}）。

\begin{lem}~
\begin{enumerate}
\item 函子 $\mA\mapsto \mA_X$ 完全忠实。
\item 对每个 $n$ 有典范同构
\[(\mA_1\star\cdots\star\mA_n)_X\cong (\mA_1)_X\ostar\cdots\ostar(\mA_n)_X.\]
特别地，函子 $\Sat_G\to \Sat_X,\ \mA\mapsto \mA_X$ 上有自然幺半结构。
\end{enumerate}
\end{lem}

\begin{proof}
注意图
\[X\tilde\times \on{Supp}(\mA)\leftarrow \hat X^+\times \on{Supp}(\mA)\to \on{Supp}(\mA)\]
中的两个映射都可写成纤维几何连通的平滑满射的逆向极限，故 (1) 由 \cite[Proposition 4.2.5]{BBD} 推出。(2) 由典范同构
\[X\tilde\times (\Gr\tilde\times\cdots\tilde\times\Gr)\cong \Gr_X\tilde\times\cdots\tilde\times\Gr_X|_\Delta\]
推出。
\end{proof}

\begin{rmk}\label{conv perverse}
本引理与命题 \ref{conv=fusion} 合起来给出命题 \ref{conv prod} 的证明。事实上它证明了更强的陈述（如 \cite{Ga} 中命题 6 之后的注所述），因为无需假设 $\mA_1$ 是 $L^+G$-等变的。
\end{rmk}

(ii) 注意 $(q_I)_*\mA$ 在 $X^I$ 上 ULA。令
\begin{equation}\label{relative coh}
\mH^*(\mA):=\bigoplus_i \mH^i(q_I)_*\mA[i].
\end{equation}
由定理 \ref{ULA prop} (1)，这是 $X^I$ 上的分次局部系。于是有典范同构
\begin{equation}\label{monoidal functor}
\mH^*(\mA_1\boxstar\cdots\boxstar\mA_n)\cong \mH^*(\mA_1)\boxtimes\cdots\boxtimes\mH^*(\mA_n),
\end{equation}
因为由 K\"unneth 公式它在 $X^{(\phi)}$ 上成立。于是得到到 $X^I$ 上分次局部系范畴的函子 $\mH^*: \Sat_{X^I}\to \on{Loc}^{\on{gr}}(X^I)$，它与外卷积积、张量积相容。因此 $\mH^*$ 关于内卷积积与张量积有自然幺半结构
\[\mH^*: (\Sat_{X^I},\ostar)\to (\on{Loc}^{\on{gr}}(X^I),\otimes).\]
对每个点 $i_x:\{x\}\to X$，取茎 $i_x^*$ 是从 $\on{Loc}^{\on{gr}}(X)$ 到分次向量空间范畴 $\on{Vect}^{\on{gr}}_{\Ql}$ 的对称幺半函子（两者都配以分次交换约束 $c_{\on{gr}}$）。虽然 $i_x^*$ 依赖点 $x$ 的选取，复合
\begin{equation}\label{decomp coh}
\Sat_G\to \Sat_X\stackrel{\mH}{\to} \on{Loc}^{\on{gr}}(X)\stackrel{i_x^*}{\to}\on{Vect}^{\on{gr}}_{\Ql}
\end{equation}
却不依赖，且由 proper 基变换典范同构于超上同调函子 $\on{H}^*$。于是也得到 $\on{H}^*$ 上的幺半结构。留给读者验证：它与命题 \ref{fiber functor} 中构造的幺半结构一致。

(iii) \eqref{eqfusion} 右端的表达式通常称为 $\mA_1,\cdots,\mA_n$ 的\textbf{融合积}。这一解释显出卷积积更多来自 $X^I$ 对称性的对称性。下面的简单陈述已足以证明命题 \ref{existence of comm}；更一般的讨论见注 \ref{symmetry on fusion}。

设 $\mA_1,\mA_2\in\Sat_X$。回忆有对合 $\Delta(\sigma):\Gr_{X^2}\cong\Gr_{X^2}$，见命题 \ref{III:Gr2}。由该命题与 \eqref{eqfusion} 可知存在典范同构
\[\Delta(\sigma)^*(\mA_1\boxstar\mA_2)\cong \mA_2\boxstar\mA_1\]
使下图交换
\[\begin{CD}
\sigma^*\mH^*(\mA_1\boxstar\mA_2)@>\cong>> \mH^*\Delta(\sigma)^*(\mA_1\boxstar\mA_2)@>\cong>>\mH^*(\mA_2\boxstar\mA_1)\\
@V\cong VV & & @VV\cong V\\
\sigma^*(\mH^*(\mA_1)\boxtimes\mH^*(\mA_2))@>c_{\on{gr}}>> \mH^*(\mA_2)\boxtimes\mH^*(\mA_1),
\end{CD}\]
其中 $c_{\on{gr}}$ 是（例如在 $(x,y)\in X^2$ 处）由
\[\begin{split}\sigma^*(\mH^*(\mA_1)\boxtimes\mH^*(\mA_2))_{(x,y)}=\mH^*(\mA_1)_y\otimes\mH^*(\mA_2)_x \\ \stackrel{c_{\on{gr}}}{\cong}\mH^*(\mA_2)_x\otimes \mH^*(\mA_1)_y=(\mH^*(\mA_2)\boxtimes\mH^*(\mA_1))_{(x,y)}\end{split}\]
给出的同构。

因此

\begin{cor}\label{existence of comm2}
存在典范同构 $c'_{\mA_1,\mA_2}:\mA_1\ostar\mA_2\cong\mA_2\ostar\mA_1$ 使下图交换
\[\begin{CD}
\mH^*(\mA_1\ostar\mA_2)@>c'_{\mA_1,\mA_2}>\cong>\mH^*(\mA_2\ostar\mA_1)\\
@V\cong VV@VV\cong V\\
\mH^*(\mA_1)\otimes\mH^*(\mA_2)@>c_{\on{gr}}>>\mH^*(\mA_2)\otimes\mH^*(\mA_1).
\end{CD}\]
\end{cor}

最后，因 $\Sat_G$ 是 $\Sat_X$ 的全子范畴，且超上同调函子 $\on{H}^*:\Sat_G\to \on{Vect}_{\Ql}$ 同构于 \eqref{decomp coh}，命题 \ref{existence of comm} 得证。

\begin{rmk}\label{symmetry on fusion}
更一般地，对每个 $\phi: J\twoheadrightarrow I$（$I$ 无次序）有函子
\[\underset{\phi}{\fstar}: \times_{i\in I}\Sat_{X^{J_i}}\to \Sat_{X^J},\quad ( \mA_{i}, i\in I)\mapsto \underset{\phi}{\fstar}\mA_i:=j_{!*}((\boxtimes_i \mF_i)|_{X^{(\phi)}}).\]
（右端关于 $q_I$ 是 ULA，这是命题 \ref{conv=fusion} 的推论。）于是由定理 \ref{III:factorization}，对每个 $K\stackrel{\phi}{\twoheadrightarrow} J\stackrel{\psi}{\twoheadrightarrow} I$ 有函子的典范同构
\begin{equation}\label{fusion prod}
\underset{\phi}{\fstar}= \underset{\psi\phi}{\fstar}(\times_{i\in I} \underset{\phi_i}{\fstar}): \times_{j} \Sat_{X^{K_j}}\to \Sat_{X^K},
\end{equation}
以及
\begin{equation}
\Delta(\phi)^\bullet \underset{\psi\phi}{\fstar}= \underset{\psi}{\fstar}(\times_{i\in I}\Delta(\phi_i)^\bullet): \times_{i}\Sat_{X^{K_i}}\to \Sat_{X^J}.
\end{equation}
由这些公式与 \cite[Proposition 1.5]{DM} 直接推出：$(\Sat_{X^I},\ostar)$ 有张量范畴结构，且 $\mH^*: (\Sat_{X^I},\ostar)\to (\on{Loc}^{\on{gr}}(X^I),\otimes)$ 有张量结构。
\end{rmk}

\begin{ex}
设 $G=\GL_n$。上述用融合积构造交换约束的方式与例 \ref{momega1BD} 合起来蕴含：在例 \ref{conv v.s. springer} 中把 $\Gr_{\omega_1,\ldots,\omega_1}\cap m^{-1}(t^\la)$ 等同于 Springer 纤维的该等同下，对称群 $S_m$ 在 $\Hom(V_\la, V_{\omega_1}^{\otimes m})$ 上的作用与 Springer 作用相差一个符号。特别地，由此得到 Schur--Weyl 对偶。注意该等同也是 Laumon 从 Springer 层构造自守层的出发点（参见 \cite{FGKV}）。
\end{ex}

% ============================================================
%  ch15  第五讲 §5.5 自举（bootstrap）
%  原文行 3075–3210
% ============================================================

\subsection{自举}

要把几何 Satake 对应用于（几何）朗兰兹纲领，重要的是不假设 $k$ 代数闭，并把它推广为 $\Sat_{X^I}$（对各个 $I$）与涉及 Langlands 对偶群 $\hat G$ 的某些范畴之间的等价。本小节说明如何自举先前的结果以得到这样的推广。

先讲一些一般性的形式框架。

\begin{dfn}
设 $K$ 是域 $E$ 上的线性代数群，$\mC$ 是 $E$-线性 Abel 范畴且容许无穷直和（于是可以与 $E$-向量空间张量）。$\mC$ 中 $K$ 的\textbf{表示}是一个对象 $M$ 连同态射 $\rho: M\to M\otimes_{E}\mO_K$，满足余作用通常的交换图。$K$-表示的态射 $f: (M_1,\rho_1)\to (M_2,\rho_2)$ 是 $\mC$ 中与 $\rho_i$ 相容的态射 $f:M_1\to M_2$。$\mC$ 中 $K$-表示的范畴记作 $\on{Rep}(K,\mC)$。

一般地，若 $\mC$ 未必容许无穷直和，先过渡到其归纳完备化 $\on{Ind}\mC$，再把 $\on{Rep}(K,\mC)$ 定义为 $\on{Rep}(K,{\on{Ind}\mC})$ 中那些底对象 $M$ 属于 $\mC$ 的 $(M,\rho)$ 构成的全子范畴。容易验证：当 $\mC$ 容许无穷直和时，新定义与旧定义一致。
\end{dfn}

\begin{ex}
若 $\mC=\on{Vect}_E$ 是有限维 $E$-向量空间范畴，则 $\on{Rep}(K,\mC)$ 就是 $K$ 的有限维代数表示范畴，为简洁记作 $\Rep_E(K)$。
\end{ex}

\begin{ex}\label{KLoc}
设 $\Ga$ 是投射有限群，$\on{Rep}^c_\Ql(\Ga)$ 记 $\Ga$ 的连续 $\ell$-进表示范畴。设 $\on{Rep}(K,{\on{Rep}^c_\Ql(\Ga)})$ 是上述 $K$ 在 $\on{Rep}^c_\Ql(\Ga)$ 中表示的范畴。则容易看出，$\on{Rep}(K,{\on{Rep}^c_\Ql(\Ga)})$ 的对象恰是 $\Ga\times K$ 的表示，它对第一个因子连续、对第二个因子代数。
\end{ex}

注意若 $F:\mC\to \mD$ 是加性函子，它典范地提升为函子 $F:\on{Rep}(K,{\mC})\to \on{Rep}(K,{\mD})$。

表示论中若干经典构造可推广到此框架，择要列举。给定 $K$-线性表示 $V$ 与 $M\in \on{Rep}(K,{\mC})$，张量积 $M\otimes V$ 作为 $\on{Rep}(K,{\mC})$ 的对象有意义。若 $\mC$ 是幺半 Abel 范畴且 $M,N\in \on{Rep}(K,{\mC})$，则 $M\otimes N$ 带有自然的对角 $K$-作用，因此 $\mC$ 的幺半结构可升级为 $\on{Rep}(K,{\mC})$ 上的幺半结构。还可定义不变量 $M^K$ 为映射 $\rho:M\to M\otimes \mO_K$ 与映射 $M=M\otimes_EE\stackrel{e_K}{\to} M\otimes_E\mO_K$ 的等化子。则 $M\mapsto M^K$ 是从 $\on{Rep}(K,{\mC})$ 到 $\mC$ 的左正合函子。

现在设 $k$ 是未必代数闭的域，$\bar k$ 是其可分闭包，$\Ga_k=\Gal(\bar k/k)$。设 $G$ 是 $k$ 上连通约化群，$\underline G=G\otimes_k\mO$，$\Gr$ 记 $\underline G$ 的仿射 Grassmann 簇。设 $\hat G=\Aut^{\otimes}(\on{H}^*)$ 是 $\Sat_{G_{\bar k}}$ 的 Tannaka 群，配以命题 \ref{canonical pinning} 的典范带钉\footnote{要把它与通常组合构造给出的群等同，其实并不需要定理 \ref{birth of dual group}。}。下面先说明如何把完整 Langlands 对偶群 ${^L}G$ 定义为 $\hat G$ 与 $\Ga_k$ 的半直积。事实上将看到 ${^L}G$ 有两个版本。

如 \cite{RZ} 所解释，有 $\Ga_k$ 在 $\hat G$ 上的典范作用 $\on{act}^{\on{geo}}$，定义如下。每个 $\ga\in \Ga_k$ 诱导自同构 $\Gr\otimes \bar k\cong \Gr\otimes \bar k$，故拉回 $\ga^*$ 诱导 $\Sat_{G_{\bar k}}$ 的张量自同构。此外有张量函子的典范同构 $\al_\ga: \on{H}^*\cong \on{H}^*\circ \ga^*$。于是对 $g\in \hat G$（看作张量函子 $\on{H}^*$ 的自同构），$\ga(g)$ 定义为 $\on{H}^*$ 的张量自同构
\begin{equation}\label{gaaction}
\on{H}^*\stackrel{\al_\ga}{\longrightarrow}\on{H}^*\circ\ga^*\stackrel{g\circ \id}{\to}\on{H}^*\circ\ga^*\stackrel{\al^{-1}_\ga}{\to}\on{H}^*.
\end{equation}
由于 $\al_\ga$ 保持上同调次数，$\on{act}^{\on{geo}}$ 保持 $2\rho:\bG_m\to \hat G$，从而由推论 \ref{canonical TB} 保持 $\tilde T\subset \tilde B\subset \hat G$；但它不保持命题 \ref{canonical pinning} 中的主幂幺元 $N$。

设 $\hat G_\ad$ 是 $\hat G$ 的伴随商。则余特征 $\bG_m\stackrel{2\rho}{\to} \hat G\to \hat G_\ad$ 有唯一的平方根，记作 $\rho_\ad$。设 $\on{cycl}:\Ga_k\to\bZ_\ell^\times$ 是由 $\Ga_k$ 在 $\bar k$ 的 $\ell^\infty$-单位根上的作用定义的分圆特征。定义同态
\[\chi:\Ga_k\stackrel{\on{cycl}}{\to} \bZ_\ell^\times\stackrel{\rho_\ad}{\to} \hat G_{\ad}(\Ql),\]
它通过内自同构诱导 $\Ga_k$ 在 $\hat G$ 上的作用 $\Ad_{\chi}:\Ga_k\to\Aut(\hat G)$。

\begin{lem}\label{comparison}
作用 $\on{act}^{\on{alg}}:=\on{act}^{\on{geo}}\circ \Ad^{-1}_{\chi}$ 保持命题 \ref{canonical pinning} 中定义的 $\hat G$ 的带钉。
\end{lem}

\begin{proof}
这由第一 Chern 类 $c_1(\mO(1))$ 属于 $\on{H}^2(\Gr_G\otimes\bar k, \Ql(1))^{\Ga_k}$ 推出。细节见 \cite[Proposition A.6]{RZ}。
\end{proof}

\begin{rmk}
本注后文不需要。回忆总有 $\Ga_k$ 通过外自同构在 $G\otimes \bar k$ 上的作用 $\psi: \Ga_k\to \on{Out}(G\otimes \bar k)$。另一方面，由定理 \ref{birth of dual group} 有典范同构 $\on{Out}(G\otimes\bar k)\cong \on{Out}(\hat G)=\Aut(\hat G, \tilde B, \tilde T, N)$。如 \cite{RZ} 所解释，在此同构下 $\on{act}^{\on{alg}}$ 与 $\psi$ 一致。
\end{rmk}

把 $G$ 的完整 Langlands 对偶群定义为
\[{^L}G^{\on{alg}}:= \hat G\rtimes_{\on{act}^{\on{alg}}}\Ga_k.\]
因 $\Ga_k$ 投射有限，且其在 $\hat G$ 上的作用 $\on{act}^{\on{alg}}$ 通过有限商分解，可把 ${^L}G^{\on{alg}}$ 看作 pro-代数群。

另一方面，按上述讨论，考虑
\[{^L}G^{\on{geo}}:= \hat G\rtimes_{\on{act}^{\on{geo}}}\Ga_k\]
也很自然，它与 $\Sat_G$ 联系更紧密。把 ${^L}G^{\on{geo}}$ 看作拓扑群。

作为一条勘误，我们指出 \cite[Corollary A.8]{RZ} 有误，正确陈述如下。

\begin{lem}\label{geom Langlands dual}
设 $\chi: \Ga_k\to \hat G_\ad(\Ql)$ 可提升为连续同态 $\tilde \chi:\Ga_k\to \hat G(\Ql)$。则有（拓扑群的）自然同构
\[{^L}G^{\on{geo}}\cong{^L}G^{\on{alg}},\quad (g,\ga)\mapsto (g\tilde\chi(\ga),\ga).\]
\end{lem}

\begin{rmk}\label{lifting condition}
提升 $\tilde \chi$ 在下列情形存在。
\begin{enumerate}
\item[(i)] 余特征 $\rho_\ad$ 提升为余特征 $\tilde \rho: \bG_m\to \hat G$。例如 $\hat G=\GL_n$ 时成立。注意并不要求 $2\tilde \rho=2\rho$。
\item[(ii)] 分圆特征 $\on{cycl}: \Ga_k\to \bZ_\ell^\times$ 有平方根，即存在半 Tate 扭 $\Ql(\frac{1}{2})$。例如 $\on{char} k=p>0$，或 $k$ 是 $p$-进域（$p\neq \ell$）时成立。
\end{enumerate}
\end{rmk}

\begin{ex}\label{geom v.s. alg dual}
这是两个群不同构的例子。取 $k=\bR$、$G=\on{PGL}_2(\bR)$。则 ${^L}G^{\on{alg}}=\SL_2\times \Gal(\bC/\bR)$，而 ${^L}G^{\on{geo}}=\SL_2\rtimes \Gal(\bC/\bR)$，其中复共轭通过
\[\begin{pmatrix} a & b\\ c & d \end{pmatrix}\mapsto \begin{pmatrix} a & -b\\ -c & d \end{pmatrix}\]
作用在 $\SL_2$ 上。注意 ${^L}G^{\on{alg}}$ 的中心同构于 $\bZ/2\times \bZ/2$，而 ${^L}G^{\on{geo}}$ 的中心同构于 $\bZ/4$，故两者不同构。
\end{ex}

虽然 ${^L}G^{\on{geo}}$ 不能看作 pro-代数群，但总可把它嵌入更大的 pro-代数群，如下。令
\begin{equation}\label{Deligne modification}
\hat G^T:=\hat G\rtimes\bG_m,
\end{equation}
其中 $\bG_m$ 通过 $\rho_\ad:\bG_m\to \hat G_\ad$ 以内自同构作用在 $\hat G$ 上。例如若 $\hat G=\SL_2$，则 $\hat G^T=\GL_2$。作用 $\on{act}^{\on{alg}}$ 自然延拓为 $\Ga_k$ 在 $\hat G^T$ 上的作用，仍记作同一记号。令
\[{^L} G^T=\hat G^T\rtimes_{\on{act}^{\on{alg}}}\Ga_k,\]
它可看作 pro-代数群。以 $d: \hat G^T\to \bG_m$ 记自然投影，则有短正合列
\[1\to {^L}G^{\on{geo}}\to {^L}G^T\stackrel{d\times \on{cycl}^{-1}}{\longrightarrow}\bG_m\to 1.\]

\begin{rmk}
注意 $\hat G^T$ 其实是 $G$ 被 $\bG_m$ 的中心扩张的对偶群。看来 ${^L}G^T$ 正是 \cite[Definition 5.3.2]{BuGe} 中出现的 $C$-群。如上面所提示的，也如 \cite{BuGe} 所建议，为了给自守表示配 Galois 表示，这似乎是更自然的对象。
\end{rmk}

以 $\on{Rep}^c_{\Ql}({^L}G^{\on{geo}})$ 记 ${^L}G^{\on{geo}}$ 的有限维连续表示中限制到 $\hat G$ 为代数者的范畴。类似地定义 $\on{Rep}^c_{\Ql}({^L}G^T)$。设 $B$ 是 $\bG_m$ 的正则表示，经同态 $d\times \on{cycl}^{-1}$ 看作 $\on{Rep}_{\Ql}({^L}G^T)$ 的归纳完备化中的代数对象。回忆对幺半范畴 $\mC$ 与代数对象 $B\in \mC$，可谈 $\mC$ 中的 $B$-模，其范畴记作 $B\on{-mod}(\mC)$。则有等价
\[B\on{-mod}(\on{Rep}^c_{\Ql}({^L}G^T))\simeq \on{Rep}^c_{\Ql}({^L}G^{\on{geo}}),\quad V\mapsto V\otimes_B\Ql,\]
其中 $B\to \Ql$ 是余单位映射（对偶于 $1\in \bG_m$）。

现在以 $\Sat_G$ 记 $\Gr$ 上 $L^+G$-等变反常层范畴，配以卷积积（注意 \eqref{Lus conv} 定义在 $k$ 上）。例如若 $G=\{e\}$ 平凡，这就是 $\Ga_k$ 的连续 $\ell$-进表示范畴 $\on{Rep}^c_{\Ql}(\Ga_k)$，配以通常张量积。下面是几何 Satake 等价的完整版本。

\begin{thm}\label{full Sat equiv}
超上同调函子 $\on{H}^*:\Sat_G\to \on{Rep}^c_{\Ql}(\Ga_k)$ 提升为典范等价
\begin{equation}\label{full Sat}
F:\Sat_G\simeq \on{Rep}^c_{\Ql}({^L}G^{\on{geo}}),
\end{equation}
且下图交换（差一个典范同构）
\[\begin{CD}
\Sat_G@>\simeq>> \on{Rep}^c_{\Ql}({^L}G^{\on{geo}})\\
@V{\text{拉回}}VV@VV{\text{限制}}V\\
\Sat_{G_{\bar k}}@>\simeq>> \on{Rep}_{\Ql}(\hat G).
\end{CD}\]
\end{thm}

\begin{proof}
勾勒证明（另一途径见 \cite[\S~5]{Ri}）。因 $\Gr$ 上的层拉回到 $\Gr\otimes \bar k$ 是 $\Ga_k$-等变的，由一般形式框架（\cite[Lemma A.3]{RZ}）可知超上同调按定量所述分解，且图交换。剩下要证它是等价，将构造拟逆函子来证。

为说明想法，假设 $G$ 分裂，从而每个 $L^+G$-轨道定义在 $k$ 上。若 $G$ 非分裂，可先过渡到使 $G_{k'}$ 分裂的有限扩张 $k'/k$，再做进一步的 Galois 下降。

对 Schubert 簇 $\Gr_{\leq \mu}$，以 $\IC_\mu$ 记其相交上同调复形（故 $\IC_\mu|_{\Gr_\mu}=\Ql[(2\rho,\mu)]$）。先考虑 $\Sat_G$ 中由 $\mu\in \xcoch(T)^+$、$i\in \bZ$ 的 $\IC_\mu(i)$ 的直和张成的全子范畴 $\Sat_G^T$。

\begin{lem}~
\begin{enumerate}
\item 范畴 $\Sat_G^T$ 是 Tannaka 范畴，纤维函子由 $\on{H}^*$ 给出。
\item Tannaka 群 $\Aut^{\otimes}(\on{H}^*)$ 同构于 \eqref{Deligne modification} 中定义的 $\hat G^T$。
\end{enumerate}
\end{lem}

\begin{proof}
(1) 事实上由推论 \ref{Satake fiber}，$\Sat_G^T$ 对卷积积封闭。又因命题 \ref{existence of comm} 中的交换约束 $c$ 是 $\Ga_k$-等变的（因 $c_{\on{vect}}$ 是 $\Ga_k$-等变的），$\Sat_G^T$ 对称幺半。于是如定理 \ref{Tannakian structure} 的论证，它是 Tannaka 范畴。

(2) 记该 Tannaka 群为 $\tilde G^T$。有正合列
\[1\to \hat G\to \tilde G^T\stackrel{d}{\to} \bG_m\to 1,\]
其中第一个映射由拉回 $\Sat_G^T\to \Sat_{G_{\bar k}}$ 诱导，第二个由 $\IC_e(i)$（$i\in\bZ$）的直和张成的全子范畴给出。投影 $d:\tilde G^T\to \bG_m$ 有分裂 $\bG_m\to \tilde G^T$，对应于 $\on{H}^*(\mA)$ 上的下述分次
\[\on{H}^*(\mA)=\bigoplus_i \Ql(i)\otimes \Hom(\Ql(i), \on{H}^*(\mA)).\]
于是如引理 \ref{geom Langlands dual}，用 $c_1(\mO(1))\in \on{H}^2(\Gr_G\otimes\bar k,\bQ_\ell(1))^{\Ga_k}$ 即得 $\tilde G^T=\hat G^T$。
\end{proof}

固定 $\on{H}^*$ 的拟逆 $S^T: \Rep_\Ql(\hat G^T)\simeq \Sat_G^T$。

现在考虑整个范畴 $\Sat_G$，将用到本节开头的一般形式框架。设 $R^T$ 是 $\hat G^T$ 上正则函数环，经左正则表示看作 $\on{Ind}\on{Rep}_{\Ql}(\hat G^T)$ 的对象。令 $\mR^T=S^T(R^T)\in \on{Ind}\Sat^T_G$ 为对应对象。因 $R^T$ 还有右正则表示，$\mR^T$ 其实是 $\on{Rep}(\hat G^T,{\on{Ind}\Sat_G})$ 中的对象。因 $G$ 分裂，$\Rep^c_\Ql({^L}G^T)=\on{Rep}(\hat G^T,{\on{Rep}^c_{\Ql}(\Ga_k)})$。$\on{Rep}(\hat G^T,{\on{Rep}^c_{\Ql}(\Ga_k)})$ 中的对象 $V$ 可看作带 $\hat G^T$ 作用的 $\Spec k$ 上层。于是由一般形式框架，$e_*V\star\mR^T$ 作为 $\Rep(\hat G^T, \on{Ind}\Sat_G)$ 的对象有意义，故 $(e_*V\star \mR^T)^{\hat G^T}$ 是 $\on{Ind}\Sat_G$ 的对象。由此得到张量函子
\[\on{Rep}(\hat G^T,{\on{Rep}^c_{\Ql}(\Ga_k)})\to \on{Ind}\Sat_G,\quad V\mapsto (e_*V\star \mR^T)^{\hat G^T},\]
仍记作 $S^T$，因为它在 $\Rep_\Ql(\hat G^T)\subset \on{Rep}(\hat G^T,{\on{Rep}^c_{\Ql}(\Ga_k)})$ 上的限制恰是开头固定的函子 $S^T$。

回忆把 $B=\mO_{\bG_m}$ 看作 $\on{Rep}(\hat G^T,{\on{Rep}^c_{\Ql}(\Ga_k)})$ 的对象。则
\[\mB:=S^T(B)= \IC_e\otimes B.\]
因此 $S^T$ 把 $\on{Rep}(\hat G^T,{\on{Rep}^c_{\Ql}(\Ga_k)})$ 中的 $B$-模对象送到 $\on{Ind}\Sat_G$ 中的 $\mB$-模对象。于是可定义
\[S: \on{Rep}^c_{\Ql}({^L}G^{\on{geo}})\simeq B\on{-mod}(\on{Rep}^c_{\Ql}({^L}G^T)) \to \on{Ind}\Sat_G:\quad V\mapsto S^T(V)\otimes_{B}\Ql,\]
其中如前 $B\to\Ql$ 是余单位映射。

注意对 $V\in \on{Rep}(\hat G^T,{\on{Rep}^c_{\Ql}(\Ga_k)})$，
\[\on{H}^*\circ S^T(V)=\on{H}^*((e_*V\star \mR^T)^{\hat G^T})=(V\otimes R^T)^{\hat G^T}=V,\]
看作 ${^L}G^{\on{geo}}$ 的表示。由此可知 $S$ 正是 $\on{H}^*:\Sat_G\to \on{Rep}^c_{\Ql}({^L}G^{\on{geo}})$ 所需的拟逆。细节留作练习。
\end{proof}

\begin{rmk}\label{simplify}
若注 \ref{lifting condition} 中某一条件成立，定理的证明可简化。例如若 $\rho\in\xch(T)$ 或存在半 Tate 扭，可定义标准化的相交上同调复形
\[\IC_\mu^{\on{norm}}=\IC_\mu((2\rho,\mu)),\]
并定义由这些标准化 IC 层的直和张成的子范畴 $\Sat_G^N\subset \Sat_G$。由推论 \ref{Satake fiber}，它是 $\Sat_G$ 的幺半子范畴，且拉回诱导幺半范畴的等价
\[\Sat_G^N\simeq \Sat_{G_{\bar k}}.\]
因此 $\hat G$ 上正则函数环 $R$ 给出对象 $\mR\in \on{Rep}(\hat G,{\on{Ind}\Sat_G})$，可用它代替 $\mR^T$ 直接定义函子 $S$。
\end{rmk}

因 ${^L}G^{\on{alg}}$ 可看作 pro-代数群，有代数表示范畴 $\on{Rep}_\Ql({^L}G^{\on{alg}})$。若提升 $\tilde \chi:\Ga_k\to \hat G(\Ql)$ 存在，则经引理 \ref{geom Langlands dual} 中的自然同构，可把 $\on{Rep}_\Ql({^L}G^{\on{alg}})$ 看作 $\on{Rep}^c_{\Ql}({^L}G^{\on{geo}})$ 的全子范畴。在 $\rho\in \xch(T)$ 或存在半 Tate 扭的假设下，\cite{RZ} 说明了如何在等价 \eqref{full Sat} 下选出 $\Sat_G$ 中对应于 $\on{Rep}_\Ql({^L}G^{\on{alg}})$ 的全子范畴（\emph{同上}记作 $\on{P}_{L^+G}^f(\Gr_G)$）。当 $k$ 有限时，陈述尤为简洁。

\begin{prop}
设 $k=\bF_q$ 有限。以 $\Sat_G^0\subset \Sat_G$ 记权为零的半单纯反常层构成的全子范畴。则等价 \eqref{full Sat} 限制为等价
\[F:\Sat_G^0\simeq \on{Rep}_\Ql({^L}G^{\on{alg}}).\]
\end{prop}

\begin{proof}
见 \cite[Remark A.9, Proposition A.10]{RZ}。注意对 $\mA\in\Sat_G^0$，$\Ga_k$ 在 $F(\mA)$ 上的作用与 $\Ga_k$ 在 $\oplus_i \on{H}^i(\mA(\frac{i}{2}))$ 上的自然作用一致。
\end{proof}

下面描述对所有 $\Sat_{X^I}$ 的几何 Satake 等价的一个版本。为简单起见，假设 $G$ 分裂，且 $\rho\in \xch(T)$ 或存在半 Tate 扭。则如上所述，有完全忠实函子
\[S: \Rep(\hat G)\simeq \Sat_G^N\to \Sat_G.\]
设 $X$ 是 $k$ 上的曲线。对有限非空集 $I$，以 $\Loc(X^I)$ 记 $X^I$ 上 $\Ql$-局部系范畴，另有 $\Rep(\hat G^I,\Loc(X^I))$。联系这些范畴的有下述函子。设 $\phi:J\twoheadrightarrow I$。
\begin{enumerate}
\item[(i)] 限制 $\Delta(\phi)^*: \on{Rep}(\hat G^J,{\on{Loc}(X^J)})\to \on{Rep}(\hat G^I, {\on{Loc}(X^I)})$，其中 $\hat G^I$ 经嵌入 $\hat G^I\to \hat G^J$ 作用在 $\Delta(\phi)^*\mF$ 上。
\item[(ii)] 外积 $\boxtimes_{i\in I}:\prod_{i\in I} \on{Rep}(\hat G^{J_i},{\on{Loc}(X^{J_i})})\to \on{Rep}(\hat G^J,\on{Loc}(X^J))$。
\end{enumerate}

\begin{thm}\label{multiple point Sat}
存在典范等价
\[S_I: \on{Rep}(\hat G^I, \on{Loc}(X^I))\simeq \Sat_{X^I},\]
它与限制函子及外积函子在显然的意义下相容。
\end{thm}

\begin{rmk}
本定理的一个版本最早由 Gaitsgory 证明（\cite{Ga2}）。事实上他在两边都考虑了更大的范畴：左端为 $\on{Rep}(\hat G^I, \on{P}(X^I))$，右端为 $\on{P}_{(L^+G)_{X^I}}(\Gr_{X^I})$。然而如 (1) 中那样的限制函子在这些更大范畴上没有定义（见注 \ref{triv group}）。我们的途径与\emph{同上}不同。
\end{rmk}

下述推论在 \cite{Laf} 中起关键作用。

\begin{cor}
对每个 $I$ 有完全忠实函子 $S_I: \on{Rep}(\hat G^I)\to \Sat_{X^I}$，且对每个满射 $\phi: J\twoheadrightarrow I$ 有函子的典范同构 $\Delta(\phi)^*S_J\cong S_I\circ \Res_{\hat G^I}^{\hat G^J}$，其中 $\Res_{\hat G^I}^{\hat G^J}$ 是嵌入 $\hat G^I\to \hat G^J$ 对应的通常限制函子。
\end{cor}

说明该定理的证明思路。首先断言函子 $\mH^* :\Sat_{X^I}\to \on{Loc}(X^I)$ 可典范地提升为函子
\[F_I: \Sat_{X^I}\to \on{Rep}(\hat G^I,{\on{Loc}(X^I)}).\]
即：设 $\eta^I$ 是 $X^I$ 的通用点，$\overline{\eta^I}$ 是 $\eta^I$ 上的几何点。由于限制函子
\[i_{\overline{\eta^I}}^*: \on{Loc}(X^I)\to \on{Rep}^c_{\Ql}(\Gal(\overline{\eta^I}/\eta^I))\]
完全忠实，只需说明每个 $\mH^*(\mA)_{\overline{\eta^I}}$ 带有 $\hat G^I$ 的典范作用。然而 $(\Gr_{X^I})_{\eta^I}\cong (\prod_i \Gr_X)_{\eta^I}$，故所需提升由定理 \ref{full Sat equiv} 推出。

为反方向定义函子，需要 $\on{Rep}(\hat G^I, {\on{Ind}\Sat_{X^I}})$ 中类似于 $\Gr$ 上层 $\mR^T$（或在我们的假设下为注 \ref{simplify} 中的层 $\mR$）的对象。定义
\[\mR_X:= S(R)_X,\]
它（由定理 \ref{full Sat equiv} 中的论证）是 $\on{Rep}(\hat G, {\on{Ind}\Sat_X})$ 的对象。对一般 $I$，定义
\[\mR_{X^I}:=\boxstar_{i\in I} \mR_X\]
为 $\mR_X$ 的自外卷积积，自然看作 $\on{Rep}(\hat G^I,{\on{Ind}\Sat_{X^I}})$ 的对象。

于是可定义函子 $S_{I}: \on{Rep}(\hat G^I,{\on{Ind}\on{Loc}(X^I)})\to \on{Ind}\Sat_{X^I}$ 为
\[S_{I}(\mF)= ((e_I)_*\mF[|I|]\ostar\mR_{X^I})^{\hat G^I}.\]
断言：对每个 $\mA\in\Sat_{X^I}$ 有典范同构 $S_IF_I(\mA)\cong \mA$；对每个 $\mF\in \on{Rep}(\hat G^I, {\on{Loc}(X^I)})$ 有典范同构 $F_IS_I(\mF)\cong \mF$。注意由 ULA 性质，限制函子
\[\Sat_{X^I}\to \on{P}_{((L^+G)_{X^I})_{\eta^I}}((\Gr_{X^I})_{\eta^I})\]
完全忠实。因此只需在 $\eta^I$ 上构造这些同构，这就化归到定理 \ref{full Sat equiv}。

由 $\mR_{X^I}$ 的分解性质容易验证 $S_I$ 与外卷积积相容；由卷积积作为融合积的解释，$F_I$ 与到对角线的限制相容。定理得证。

% ============================================================
%  ch16  第五讲 §5.6 从几何 Satake 到经典 Satake
%  原文行 3300–3400
% ============================================================

\subsection{从几何 Satake 到经典 Satake}\label{applications of sat}

本小节讨论几何 Satake 等价与经典 Satake 同构的关系。为简单起见，假设 $G$ 在 $k=\bF_q$ 上分裂；一般情形见 \cite{RZ}。以 $\sigma$ 记 $\bF_q$ 的几何 Frobenius 自同构。固定半 Tate 扭 $\Ql(\frac{1}{2})$，等价地固定 $q$ 在 $\Ql$ 中的一个平方根 $q^{1/2}$（$=\tr(\sigma,\Ql(\frac{1}{2}))^{-1}$）。为联系几何 Satake 与经典 Satake，还需选取同构 $\iota:\Ql\cong \bC$。可安排 $\Ql(\frac{1}{2})$ 与 $\iota$ 的选取使
\[\iota(q^{1/2})\in \bR^{<0}.\]
这样选取的好处是：在 $\iota$ 下，几何 Frobenius 在 $\Ql[i](\frac{i}{2})$ 上的迹对每个 $i\in \bZ$ 都\emph{为正}。

以 $H_G$ 记经典球 Hecke 代数，即 $G(F)$ 上 $G(\mO)$-双不变的 $\bZ$-值函数全体，配以卷积积
\[(f*g)(x)=\int_{G(F)}f(y)g(y^{-1}x)dy,\]
其中测度取使 $G(\mO)$ 的体积为 $1$ 者。对 $\mu\in \xcoch(T)^+$，以 $c_\mu$ 记 $G(\mO)t^\mu G(\mO)$ 上的特征函数。则 $\{c_\mu,\ \mu\in\xcoch(T)^+\}$ 构成 $H_G$ 的一组 $\bZ$-基。

另一方面，以 $R(\hat G)$ 记 $\hat G$ 的表示环，它是以不可约表示的特征为基的 $\bZ$-代数。对 $\mu\in \xcoch(T)^+$，以 $L_\mu=\on{H}^*(\IC_\mu)$ 记相应的不可约表示，$\chi_\mu$ 记相应的特征。

经典 Satake 同构（或者更确切地说，Langlands 的重新解释）是典范同构
\[\mS: H_G\otimes\bC\cong R(\hat G)\otimes\bC.\]
细节见 \cite{Gr}。下面解释几何途径。

回忆 Grothendieck 的层-函数字典给每个 $\mA\in \Sat_G$（更一般地，$\Gr_G$ 上每个 $\ell$-进复形）配上函数
\[f_\mA: \Gr(k)\to \Ql,\quad f_\mA(x)=\sum_i(-1)^i\tr(\sigma_x, \on{H}^i_{\bar x}(\mA)),\]
其中 $\bar x$ 记 $x$ 上的几何点，$\on{H}^i_{\bar x}$ 记 $\bar x$ 处的茎上同调，$\sigma_x\in \Gal(\bar x/x)$ 记几何 Frobenius 元。因 $\mA$ 是 $L^+G$-等变的，可把 $f_\mA$ 看作 $H_G\otimes_\bZ\Ql$ 的元素。此外由定义有

\begin{lem}\label{conv vs classical conv}
对 $\mA_1,\mA_2\in\Sat_G$，有 $f_{\mA_1\star\mA_2}=f_{\mA_1}*f_{\mA_2}$。
\end{lem}

\begin{proof}
设 $m:\Gr\tilde\times\Gr\to \Gr$ 是卷积映射。给定 $x\in \Gr(k)$，纤维 $m^{-1}(x)(k)$ 可等同于 $\{(y, y^{-1}x)\mid y\in G(F)/G(\mO)\}$，且由构造有茎的典范同构 $(\mA_1\tilde\boxtimes\mA_2)_{(y,y^{-1}x)}\cong (\mA_1)_y\otimes(A_2)_{y^{-1}x}$。于是按测度的选取，
\[\begin{split}
f_{\mA_1}*f_{\mA_2}(x)&=\int_{G(F)}f_{\mA_1}(y)f_{\mA_2}(y^{-1}x)dy=\sum_{y\in G(F)/G(\mO)}f_{\mA_1}(y)f_{\mA_2}(y^{-1}x)\\
&=\sum_{y\in m^{-1}(x)(k)} \tr(\sigma_{\bar y}, (\mA_1\tilde\boxtimes\mA_2)_{\bar y})=f_{\mA_1\star\mA_2}(x).
\end{split}\]
\end{proof}

现在令 $f_\mu=f_{\IC_\mu[-(2\rho,\mu)]}$，回忆 $\IC_\mu|_{\Gr_\mu}=\Ql[(2\rho,\mu)]$。由（仿射）Schubert 簇相交上同调复形茎上同调的纯性（参见 \cite{KL}），可写
\begin{equation}\label{stalk IC}
f_\mu=c_\mu+\sum_{\la<\mu}a_{\mu\la}(q)c_\la,
\end{equation}
其中
\[a_{\mu\la}(x)=\sum_i a_{\mu\la,i}x^i\in \bN[x]\]
是 Kazhdan--Lusztig 多项式。回忆 $a_{\mu\la,i}=\dim \on{H}^{2i-(2\rho,\mu)}_{t^\la}\IC_\mu$ 是 $\IC_\mu$ 在 $t^\la$ 处次数 $2i-(2\rho,\mu)$ 的茎上同调。

回忆幺半 Abel 范畴的 $K$-群是自然的 $\bZ$-代数（即所谓的 $K$-环）。则几何 Satake 等价诱导 $\bZ$-代数同构
\[K(\Sat_G^N)\cong K(\on{Rep}_\Ql(\hat G)),\]
其中回忆 $\Sat_G^N$ 是由标准化 IC 层张成的 $\Sat_G$ 的幺半子范畴（见注 \ref{simplify}）。以 $\sigma_q$ 记 $k=\bF_q$ 的几何 Frobenius。则由 \eqref{stalk IC} 与引理 \ref{conv vs classical conv}，层-函数字典诱导代数同构
\[f: K(\Sat_G^N)\otimes\Ql\to H_G\otimes\Ql.\]
另一方面，把 $[V]$ 送到其特征 $\chi_V$ 总给出典范同构 $\on{Ch}:K(\on{Rep}_\Ql(\hat G))\cong R(\hat G)$。合起来得到同构
\[\mS':H_G\otimes\Ql\stackrel{f^{-1}}{\cong} K(\Sat_G^N)\otimes\Ql \cong K(\on{Rep}_\Ql(\hat G))\otimes\Ql\stackrel{\on{Ch}}{\cong}R(\hat G)\otimes\Ql.\]

\begin{lem}
在所选同构 $\iota:\Ql\cong\bC$ 下，$\mS'=\mS$。特别地
\[\mS^{-1}(\chi_\mu)=(-\iota(q^{1/2}))^{-(2\rho,\mu)}f_\mu.\]
\end{lem}

证明的梗概见 \cite{RZ} 末尾。首先，函子 $\on{CT}:\Sat_G\to \Sat_T$ 自然提升为函子
\[\on{CT}^N: \Sat_G^N\to \Sat_T^N.\]
（事实上由定理 \ref{MV theory} (iii) 可知 $\on{H}_c^*(\mA)$ 是 Tate 型的。）在层-函数字典与所选的 $\iota$ 下，函子 $\on{CT}^N$ 恰好对应经典 Satake 变换（例如显式公式见 \cite[\S\ 3]{Gr}）。

\begin{rmk}
(i) 若不愿引入平方根 $q^{1/2}$，可用等价 $\Sat_G^T\cong \Rep_\Ql(\hat G^T)$ 表述 Satake 同构的一个版本。

(ii) 由于对每个 $\mA\in \Sat_G^N$ 有 $\iota f_\mA\in H_G\otimes\bZ[q^{\pm 1/2}]$，同构 $\mS'$ 限制为同构 $H_G\otimes \bZ[q^{\pm 1/2}]\simeq R(\hat G)\otimes \bZ[q^{\pm 1/2}]$。若还有 $\rho\in \xch(T)$，它进一步限制为同构 $H_G\otimes\bZ[q^{-1}]\simeq R(\hat G)\otimes\bZ[q^{-1}]$。这与经典理论一致（见 \cite[\S 3]{Gr}）。
\end{rmk}

还要指出，Kazhdan--Lusztig 多项式 $a_{\mu\la}$ 可经 $\hat G$ 作如下解释。回忆命题 \ref{canonical pinning} 给出的 $\hat G$ 上的带钉，它用
\begin{equation}\label{BK fil}
N_iV=(\ker N)^{i+1}
\end{equation}
在 $\hat G$ 的任意表示 $V$ 上定义递增滤过（Brylinski--Kostant 滤过）。对 $\la\in\xcoch(T)$，以 $V(\la)$ 记 $V$ 在 $\tilde T$ 作用下的相应权子空间。则滤过 \eqref{BK fil} 诱导
\begin{equation}\label{BKfilmu}
N_iV(\la)=V(\la)\cap N_iV.
\end{equation}

下述定理在 \cite{Lu} 与 \cite{Bry} 中证明，\cite[\S~5]{Gi1} 给出另一方法（另见 \cite[\S~5]{Z11}）。

\begin{thm}
设 $\mu\in\xcoch(T)^+$，$L_\mu=\on{H}^*(\IC_\mu)$ 是 $\hat G$ 的相应表示。则
\[a_{\mu\la}(x)=\sum_i \dim \on{gr}^N_iL_\mu(\la)x^i.\]
\end{thm}

% ============================================================
%  ch17  附录 A §A.1 层论补充：等变反常层
%  原文行 3405–3537
% ============================================================

\appendix

\section{层论补充}\label{App}

\subsection{等变反常层范畴}

\subsubsection{基本性质}
回忆对 $k$ 上的概形 $X$，有可构造 $\ell$-进复形的有界导出范畴 $D_c^b(X)$，其中反常层范畴 $\on{P}(X)$ 作为全 Abel 子范畴。若 $X$ 带有线性代数群 $K$ 的作用（作用映射记作 $\on{act}$），则可以定义 $X$ 上 $K$-等变反常层的 Abel 范畴 $\on{P}_K(X)$，即：$\on{P}_K(X)$ 的对象是 $X$ 上的反常层 $\mF$ 连同一个沿两个映射
\[\on{act},\ \pr_2:K\times X\longrightarrow X\]
的同构 $\theta$，满足 (i) $e^*\theta=\id$，其中 $e:X\to K\times X$ 由 $K$ 的单位元给出；(ii) $K\times K\times X$ 上自然的余圈条件。

\begin{rmk}
还有更复杂的概念：$X$ 上（可构造）$\ell$-进层的（有界）等变导出范畴 $\on{D}_{K,c}^b(X)$，它不能像上面那样简单地定义。或者如 \cite{BeLu} 定义，或者把它看作代数叠 $[X/K]$ 上 $\ell$-进层的导出范畴（例如参见 \cite{LZ}）。本书只在非常形式的层面偶尔用到这个概念，即这些三角范畴之间（在不同概形上）存在良好的 Grothendieck 六函子形式体系。\zy{该六函子形式体系的一般框架见 nLab 词条 six operations（\url{https://ncatlab.org/nlab/show/six+operations}，2026-09-30 修订）。[来源：nLab six operations，2026-09-30。]}
\end{rmk}

按定义有遗忘函子 $\on{P}_K(X)\to \on{P}(X)$。

\begin{lem}\label{A:forgetful}
若 $K$ 连通，则该遗忘函子完全忠实，其本质像由那些满足 $\on{act}^*\mF\cong\pr_2^*\mF$ 的反常层 $\mF$ 构成。
\end{lem}

本引理的意义在于：$K$-等变性是 $X$ 上反常层的一个\emph{性质}，而非附加结构。注意该陈述对不连通群不成立：若 $K$ 有限，则 $\on{P}_K(\on{pt})$ 等价于 $K$ 的有限维表示范畴。该陈述对等变导出范畴也不成立。

\begin{proof}
简述第一个断言的证明，第二个留作练习。该函子显然忠实。设 $(\mF_1,\theta_1),(\mF_2,\theta_2)$ 是两个 $K$-等变反常层，$\varphi: \mF_1\to \mF_2$ 是态射，要证它与 $\theta_i$ 自动相容，即
\begin{equation}\label{A:comp}
\theta_2^{-1}\circ\on{act}^*(\varphi)\circ\theta_1=\pr_2^*(\varphi).
\end{equation}
注意因 $K$ 连通，由 \cite[Proposition 4.2.5]{BBD}，$\pr^*[\dim K]: \on{P}(X)\to \on{P}(K\times X)$ 完全忠实。于是要验证 \eqref{A:comp}，只需证 $e^*(\theta_2^{-1}\circ\on{act}^*(\varphi)\circ\theta_1)=e^*(\pr_2^*(\varphi))$，而这按定义成立。
\end{proof}

我们将用到等变范畴的下述两条性质。

\begin{lem}\label{A:unip}
设 $K_1\subset K$ 是闭正规子群。
\begin{enumerate}
\item 若 $K_1$ 在 $X$ 上的作用自由，且 $[X/K_1]$ 可由代数空间 $\bar X$ 表示，则沿 $q:X\to \bar X$ 的拉回诱导范畴等价
\begin{equation}\label{equivariance free action}
q^*[\dim K_1]:\on{P}_{K/K_1}(\bar X)\simeq \on{P}_K(X).
\end{equation}

\item 设 $K_1$ 连通。若 $K_1$ 在 $X$ 上的作用平凡，则遗忘函子
\begin{equation}\label{unipotence}
\on{P}_K(X)\to \on{P}_{K/K_1}(X)
\end{equation}
是范畴等价。
\end{enumerate}
\end{lem}

\begin{proof}
由反常层的 \'etale 下降，可假设 $X=\bar X\times K_1$，于是 (1) 容易推出。(2) 由引理 \ref{A:forgetful} 推出。
\end{proof}

\begin{rmk}\label{equivariant descent}
若把 $\on{P}_{K/K_1}$ 换成等变导出范畴，等价 \eqref{equivariance free action} 仍成立，但证明更不平凡（见 \cite[\S\ 2.6.2]{BeLu}）。
\end{rmk}

\subsubsection{扭曲外积}\label{A:twist product}
设 $p:E\to X$ 是 $K$-主丛，$Y$ 是带 $K$ 作用的概形。设 $\mF\in\on{P}(X)$、$\mG\in \on{P}_K(Y)$，则可如下构造 $X\tilde\times Y$ 上的反常层，即扭曲外积 $\mF\tilde\boxtimes\mG$。$E\times Y$ 上的反常层 $p^*\mF[\dim K]\boxtimes \mG$ 按构造是 $K$-等变的。于是由引理 \ref{A:unip} (i)，它下降为 $X\tilde\times Y$ 上唯一（在唯一同构意义下）的反常层，即 $\mF\tilde\boxtimes\mG$。注意由注 \ref{equivariant descent}，外扭曲积其实可对 $\ell$-进复形定义。

\subsubsection{等变上同调}\label{A:char class}
设 $K$ 是 $k$ 上的线性代数群，$R_{K}=\on{H}^*(\bB K)$ 是 $K$ 的分类叠的上同调。它是有限生成的分次增广交换 $\Ql$-代数，同构于多项式代数。具体地，它可实现如下。设 $\{E_n\to B_n\}$ 是一列 $K$-主丛，逼近 $K$ 的分类空间。例如把 $K$ 嵌入某个 $\GL_r$ 使 $\GL_r/K$ 拟仿射；对大的 $n$，令 $E_n:=S_{n,r}$ 为 Stiefel 簇，即 $\Gr(r,n)$ 上的通用 $\GL_r$-主丛。则 $B_n:=E_n/K$ 可由概形表示，$B_n\subset B_{n+1}$ 是闭嵌入，且 $\on{H}^*(\underrightarrow\lim_n B_n)= \on{H}^*(\bB K)$。例如若 $K=T$ 是环面，则有典范同构 $\Spec R_T\cong \Lie T$。若 $K$ 幂幺，则 $R_K=\Ql$。

设 $\mF\in \on{P}_K(X)$（或更一般地 $\mF\in\on{D}^b_{K,c}(X)$），则可谈 $K$-等变上同调 $\on{H}_K^*(X,\mF)$。它或者定义为 $\mF$ 作为商叠 $[X/K]$ 上层的上同调，或者更具体地
\[\on{H}_K^*(X,\mF):=\on{H}^*(\underrightarrow\lim_n (B_n\tilde\times X), \Ql\tilde\boxtimes \mF).\]
由构造，$\on{H}^*_K(X,\mF)$ 是 $\on{H}^*(\underrightarrow\lim_n B_n)=R_{K}$ 上的模。

由定义可知：若 $K_1\subset K$ 自由作用于 $X$，$\bar X$ 为如引理 \ref{A:unip} 中的商，则
\begin{equation}\label{coh free action}
\on{H}^*_K(X,q^*\mF)=\on{H}^*_{K/K_1}(\bar X,\mF).
\end{equation}
另一方面，若 $K_1$ 幂幺且平凡作用于 $X$，则
\begin{equation}\label{coh unip}
\on{H}^*_K(X,\mF)=\on{H}^*_{K/K_1}(X,\mF).
\end{equation}

现在设 $K$ 连通，$\mF\in\on{P}_K(X)$。与投影 $B_n\tilde\times X\to B_n$ 相伴的 Leray 谱序列诱导一个谱序列，其 $E_2$-项为 $\on{H}^*(\bB K)\otimes \on{H}^*(X,\mF)$，收敛到 $\on{H}^*_K(X,\mF)$。仿照 \cite{GKM}，称 $\mF$ \textbf{等变形式}（equivariantly formal），若该谱序列在 $E_2$-项退化。

\begin{thm}
若 $X$ proper，则 $X$ 的相交上同调层（看作 $X$ 上 $K$-等变反常层）等变形式。
\end{thm}

证明依赖「权」的理论，见 \cite[Theorem 14.1]{GKM}。事实上同一论证适用于几何来源的半单纯反常层。

对等变形式的复形 $\mF$，有非典范的同构 $\on{H}^*_K(X,\mF)\cong R_{K}\otimes\on{H}^*(X,\mF)$。特别地此时 $\on{H}^*_K(X,\mF)$ 是秩等于 $\dim \on{H}^*(X,\mF)$ 的自由 $R_{K}$-模。典范地有同构
\begin{equation}\label{de-equiv}
\on{H}^*(X,\mF)\cong \on{H}^*_K(X,\mF)\otimes_{R_{K}}\Ql,
\end{equation}
其中 $R_{K}\to\Ql$ 是增广映射。此外有等变 K\"unneth 公式：设 $X$、$Y$ 都 proper，$\mF\in\on{P}_K(X)$、$\mG\in\on{P}_K(Y)$ 是半单纯反常层，则有典范同构
\begin{equation}\label{equiv Kunneth}
\on{H}_K^*(X,\mF)\otimes_{R_{K}}\on{H}_K^*(Y,\mG)\cong \on{H}^*_K(X\times Y,\mF\boxtimes\mG).
\end{equation}
这由通常的 K\"unneth 公式与谱序列的退化推出。

再叙述等变上同调的局部化定理（拓扑版本见 \cite[(6.2)]{GKM}）。首先对 $\mF\in \on{D}_{T,c}^b(X)$ 可定义紧支撑 $T$-等变上同调为
\[\on{H}_{T,c}^*(X,\mF):=\on{H}_c^*(\underrightarrow\lim_n (B_n\tilde\times X), \Ql\tilde\boxtimes \mF).\]

\begin{thm}\label{equiv localization}
设 $K=T$ 是作用于 $X$ 的环面，$i:X^0\subset X$ 是 $T$-不动点子概形。则对 $\mF\in\on{D}_{T,c}^b(X)$，自然映射
\[\on{H}^*_T(X^0,i^!\mF)\to \on{H}^*_T(X,\mF),\quad \on{H}^*_{T,c}(X,\mF)\to \on{H}^*_{T,c}(X^0,i^*\mF)\]
在局部化到 $R_{T}$ 的通用点后成为同构。
\end{thm}

\begin{proof}
简述第一个同构的证明，第二个由 Verdier 对偶推出。设 $j: X-X^0\to X$ 是开补集。由 distinguished triangle $i_!i^!\mF\to \mF\to j_*j^*\mF\to$，化归到证明：若 $X^0=\emptyset$，则 $\on{H}^*_T(X,\mF)$ 是挠 $R_T$-模。可取 $X$ 的有限覆盖 $X=\cup X_i$，使每个 $X_i$ 是 $T$-不变的，且存在非平凡子环面 $T_i\subset T$ 自由作用于 $X_i$。于是由 \eqref{coh free action}，$\on{H}^*_T(X_i,\mF|_{X_i})$ 是挠 $R_T$-模。从而由谱序列论证，$\on{H}^*_T(X,\mF)$ 也是挠 $R_T$-模。
\end{proof}

\subsubsection{作用于归纳概形的 pro-代数群}\label{A:ind}
设 $X$ 是归纳概形。可把 $X$ 上的反常层范畴定义为 $\on{P}(X)=\underrightarrow\lim\on{P}(X_i)$，其中 $X=\underrightarrow\lim X_i$ 是把 $X$ 表为 $k$ 上有限型闭子概形的表示。这与表示的选取无关。

现设 $X$ 带有 pro-代数群 $K$ 的作用。假设
\[(S)\quad\quad \left.\begin{array}{c}\mbox{每个几何点的稳定化子只有有限多个连通分支。}\end{array}\right.\]
则可以如下定义 $X$ 上 $K$-等变反常层的范畴。

首先存在表示 $X=\underrightarrow\lim X_i$，使每个 $X_i$ 是 $k$ 上有限型的 $K$-稳定闭子概形，且 $K$ 在 $X_i$ 上的作用通过某个代数商 $K_i$ 分解。可取 $K_i$ 为 $K$ 模去连通正规子群所得的代数商。此外，若 $X_i\subset X_j$ 是闭子概形，可安排 $K_i$ 是 $K_j$ 的商。

然后定义 $\on{P}_K(X_i):=\on{P}_{K_i}(X_i)$。由引理 \ref{A:unip} (ii) 可知 $\on{P}_K(X)$ 在典范等价意义下与 $K_i$ 的选取无关。因此有完全忠实函子
\[\on{P}_K(X_i)=\on{P}_{K_i}(X_i)\simeq \on{P}_{K_j}(X_i)\to \on{P}_{K_j}(X_j)=\on{P}_K(X_j).\]
最后定义
\[\on{P}_K(X)=\underrightarrow\lim \on{P}_K(X_i).\]

现设 $E\to X$ 是有限型归纳概形上的 $K$-主丛，$Y$ 是满足 (S) 的 $K$-归纳概形。设 $\mF\in\on{P}(X)$、$\mG\in\on{P}_K(Y)$，可类似地把扭曲积 $\mF\tilde\boxtimes\mG$ 定义为 $\on{P}(X\tilde\times Y)$ 的对象。即：$\mF$ 支在某个 $X_i$ 上，$\mG$ 支在某个 $Y_j$ 上，$K$ 在 $Y_j$ 上的作用通过 $K_j$ 分解。则 $E_{ij}:=E|_{X_i}\times^KK_i\to X_i$ 是 $K_i$-主丛，且
\[X_i\tilde\times Y_j= E|_{X_i}\times^KY_j= E_{ij}\times^{K_j}Y_j\]
是 $X\tilde\times Y$ 的闭子概形。于是 \S\ref{A:twist product} 的构造给出 $\on{P}(X_i\tilde\times Y_j)\subset \on{P}(X\tilde\times Y)$ 中的对象 $\mF\tilde\boxtimes\mG$。

最后，设 $K^u$ 是 $K$ 的 pro-幂幺根，且 $K/K^u$ 代数。设 $X$ 是满足 (S) 的 $K$-归纳概形。则对 $\mF\in \on{P}_K(X)$ 有定义良好的 $K$-等变上同调 $H_K^*(X,\mF)$，它是 $R_{K}=\on{H}^*(\bB (K/K^u))$ 上的模。\S\ref{A:char class} 中关于等变上同调的全部性质在此情形下仍成立。

% ============================================================
%  ch18  附录 A §A.2 普遍局部无环性（ULA）
%  原文行 3538–3700
% ============================================================

\subsection{普遍局部无环性}\label{ULA sheaf}

设 $S$ 是概形，$s$ 是 $S$ 的几何点，以 $S_{(s)}$ 记 $S$ 在 $s$ 处的严格 Hensel 化。若 $t$ 是 $S_{(s)}$ 的几何点，则形式地写 $t\to s$，并称其为特化映射。

回忆 \cite{SGA4.5} 中的下述定义。

\begin{dfn}
设 $f:X\to S$ 是 $k$ 上有限型概形的态射。$X$ 上的 $\ell$-进复形 $\mF$ 称为关于 $f$ \textbf{局部无环}，若对每个几何点 $x\in X$ 与每个几何点 $t\in S_{(f(x))}$，自然映射 $R\Gamma(X_{(x)},\mF)\to R\Gamma( (X_{(x)})_t,\mF)$ 是同构，其中 $(X_{(x)})_t=X_{(x)}\times_{S_{(f(x))}}t$。若它在任意基变换 $S'\to S$ 之后仍局部无环，则称\textbf{普遍局部无环}（ULA）。
\end{dfn}

局部无环性可改写如下。设 $f:X\to S$ 如上，$\mF$ 是 $X$ 上的 $\ell$-进复形，$s$ 是 $S$ 的几何点，$t$ 是 $S_{(s)}$ 的几何点，$j_t: X\times_{S} t\to X$ 是自然映射。记
\[\Psi_{t\to s}(\mF):= ((j_t)_*j_t^*\mF)|_{X_s},\]
并把它看作沿特化映射 $t\to s$ 的邻近圈（回忆按约定 $(j_t)_*$ 是导出的）。事实上，若 $(S,s,\eta)$ 是严格局部 Hensel 链（即 $S$ 是局部严格 Hensel 离散赋值环的谱，$s$ 与 $\eta$ 是其闭点与一般点），$t$ 是 $\eta$ 上的几何点，则 $\Psi_{t\to s}(\mF)$ 按定义就是 $\mF$ 的一维邻近圈 $\Psi(\mF_\eta)$。

\begin{lem}\label{LA reform}
复形 $\mF$ 关于 $f$ 局部无环，当且仅当对如上每个 $s$ 与 $t$，自然映射
\begin{equation}\label{sp}
\mF|_{X_s}\to \Psi_{t\to s}(\mF)
\end{equation}
是同构。
\end{lem}

\begin{proof}
事实上，局部无环性说的正是上述映射在每个几何点 $x\in X_s$ 处诱导茎的同构。
\end{proof}

\begin{rmk}\label{LA explanation}
注意若 $(S,s,\eta)$ 是链，则引理 \ref{LA reform} 说 $\mF$ 关于 $f$ 局部无环当且仅当 $\mF_s\to \Psi(\mF_\eta)$ 是同构，等价地消没圈 $\Phi(\mF)$ 消没。因此对高维底 $S$，粗略地说，层 $\mF$ 关于 $f: X\to S$ 是 ULA，意味着沿 $S$ 的每个方向都没有消没圈。引入高维底上的邻近圈与消没圈函子即可使这一说法精确（参见 \cite[Example 1.7(b)]{Il3}）。
\end{rmk}

下面是所需的 ULA 层的基本性质。

\begin{thm}\label{ULA prop}~
\begin{enumerate}
\item 设 $S$ 是 $k$ 上的光滑簇。则 $S$ 上的复形 $\mF$ 关于恒等映射 ULA，当且仅当每个上同调层 $\mH^i(\mF)$ 是 $S$ 上的局部系。

\item 若 $f:X\to Y$ 是 $S$ 上的 proper 态射，且 $\mF$ 关于 $X\to S$ 是 ULA，则 $f_*\mF$ 关于 $S$ 是 ULA。

\item 若 $f:X\to Y$ 是 $S$ 上的光滑态射，则 $f^*\mF$ 关于 $X\to S$ 是 ULA，当且仅当 $\mF$ 关于 $Y\to S$ 是 ULA。

\item 设 $f_i:X_i\to S, i=1,2$ 是两个态射，$\mF_i$ 是 $X_i$ 上关于 $f_i$ 的 ULA 复形。则 $\mF_1\boxtimes_S\mF_2$ 关于 $f_1\times_Sf_2:X_1\times_SX_2\to S$ 是 ULA。

\item[(4)'] 设 $f_i:X_i\to S_i, i=1,2$，$\mF_i$ 是 $X_i$ 上关于 $f_i$ 的 ULA 复形。则 $\mF_1\boxtimes\mF_2$ 关于 $f_1\times f_2$ 是 ULA。
\end{enumerate}
\end{thm}

\begin{proof}
(1) 由定义可知 $S$ 上真正的层关于 $\id:S\to S$ 是 ULA 当且仅当它是局部系；一般情形对 $\mF$ 的上同调幅度归纳。

(2) 与 (3) 由函子 $\Psi_{t\to s}$ 与 proper 推前、光滑拉回交换推出。(4) 见 \cite[Corollary 2.5]{Il3}。(4)' 由 (4) 推出，因为对每个态射 $S\to S_1\times S_2$ 有
\[(X_1\times X_2)\times_{S_1\times S_2}S=(X_1\times_{S_1}S)\times_S(X_2\times_{S_2}S).\]
\end{proof}

下面讨论一个与 ULA 反常层有关的结果。

\begin{thm}\label{ULA perverse}
设 $f: X\to S$ 是 $k$-簇的态射，$S$ 光滑，$D\subset S$ 是光滑有效除子。设 $i: Z=f^{-1}(D)\to X$ 是闭嵌入，$j:U=X-Z\to X$ 是开补集。若 $X$ 上的复形 $\mF$ 关于 $f$ 是 ULA，且 $\mF|_U$ 反常，则
\[\mF\cong j_{!*}(\mF|_U), \quad i^!\mF[1]\cong i^*\mF[-1].\]
特别地 $\mF$ 也反常；此外 $i^*\mF[-1]$ 反常，且关于 $fi$ 是 ULA。
\end{thm}

\begin{proof}
问题在 \'etale 拓扑下是局部的，故可假设存在\emph{光滑}映射 $g: S\to \bA^1$ 使 $D=g^{-1}(0)$。由 \cite[Lemma 2.14]{SGA4.5}，$\mF$ 关于 $gf: X\to \bA^1$ 是 ULA。于是化归到 $f:X\to \bA^1$、$Z=f^{-1}(0)$ 的情形。由注 \ref{LA explanation}，$i^*\mF\to \Psi(j^*\mF)$ 是同构。由此
\[i^*\mF[-1]\cong \Psi(j^*\mF)[-1]=:\mC\]
反常（\cite[Corollary 4.5]{Il2}）。由 distinguished triangle
\[ j_!j^*\mF\to \mF\to i_*i^*\mF\to\]
可知 $\mF$ 只处于反常上同调次数 $-1$ 与 $0$。此外，由 distinguished triangle（见 \cite[Lemma 3.11]{SGA4.5} 的证明）
\[i^*j_*j^*\mF\to \Psi(j^*\mF)^P\stackrel{1-t}{\to} \Psi(j^*\mF)^P\to\]
可知反常上同调层 ${^p}\mH^r i^*j_*j^*\mF\cong \mC$（$r=-1,0$）且对 $r\neq -1,0$ 消没，其中 $P$ 记野惯性群，$t$ 记驯惯性群的生成元。于是由 distinguished triangle
\[i_*i^!\mF\to \mF\to j_*j^*\mF\to\]
得 $i^!\mF[1]\cong \mC$、$\mF$ 反常，且有反常层短正合列
\[0\to \mF\to j_*j^*\mF\to {^p}\mH^0j_*j^*\mF\to 0.\]
这又推出 $\mF\cong j_{!*}j^*\mF$。
\end{proof}

\begin{rmk}
当底 $S$ 光滑时，局部无环性还有 \cite[\S~5.1,~Appendix B]{BGa} 中引入的另一等价定义。按该定义，定理 \ref{ULA prop} (iv) 由 Gaitsgory 证明（私人通信），定理 \ref{ULA prop} 其余部分与定理 \ref{ULA perverse} 由 Reich 在 \cite{Rei} 中证明。
\end{rmk}

% ch19 空（附录 A.2 已并入 ch18；参考文献在 ch21）

% ch20 空（附录 A.2 已并入 ch18；参考文献在 ch21）


\clearpage
\phantomsection
\addcontentsline{toc}{section}{参考文献}
\begin{thebibliography}{99}

\bibitem{Ap} Jarod Alper, \emph{Adequate moduli spaces and geometrically reductive group schemes}, Algebr. Geom. \textbf{1} (2014), no. 4, 489--531.

\bibitem{BG} Vladimir Baranovsky and Victor Ginzburg, \emph{Conjugacy classes in loop groups and $G$-bundles on elliptic curves}, Internat. Math. Res. Notices (1996), no. 15, 733--751.

\bibitem{BL} Arnaud Beauville and Yves Laszlo, \emph{Conformal blocks and generalized theta functions}, Comm. Math. Phys. \textbf{164} (1994), no. 2, 385--419.

\bibitem{BL2} Arnaud Beauville and Yves Laszlo, \emph{Un lemme de descente}, C. R. Acad. Sci. Paris S\'er. I Math. \textbf{320} (1995), no. 3, 335--340.

\bibitem{BLS} Arnaud Beauville, Yves Laszlo and Christoph Sorger, \emph{The Picard group of the moduli of $G$-bundles on a curve}, Compositio Math. \textbf{112} (1998), no. 2, 183--216.

\bibitem{Bea} Arnaud Beauville, \emph{Symplectic singularities}, Invent. Math. \textbf{139} (2000), no. 3, 541--549.

\bibitem{BBD} A. A. Be{\u\i}linson, J. Bernstein and P. Deligne, \emph{Faisceaux pervers},  \textbf{100} (1981), 5--171.

\bibitem{Be} A. A. Be{\u\i}linson, \emph{How to glue perverse sheaves},  \textbf{1289} (1984--1986), 42--51.

\bibitem{BBE} Alexander Beilinson, Spencer Bloch and H{\'e}l{\`e}ne Esnault, \emph{$\epsilon$-factors for Gauss-Manin determinants}, Mosc. Math. J. \textbf{2} (2002), no. 3, 477--532.

\bibitem{BD0} Alexander Beilinson and Vladimir Drinfeld, \emph{Chiral algebras}, American Mathematical Society, Providence, RI, 2004.

\bibitem{BD} Alexander Beilinson and Vladimir Drinfeld, \emph{Quantization of Hitchin's integrable system and Hecke eigensheaves}, available at www.math.uchicago.edu/~mitya/langlands.

\bibitem{BeLu} Joseph Bernstein and Valery Lunts, \emph{Equivariant sheaves and functors}, Springer-Verlag, Berlin, 1994.

\bibitem{BFM} Roman Bezrukavnikov, Michael Finkelberg and Ivan Mirkovi{\'c}, \emph{Equivariant homology and $K$-theory of affine Grassmannians and Toda lattices}, Compos. Math. \textbf{141} (2005), no. 3, 746--768.

\bibitem{BF} Roman Bezrukavnikov and Michael Finkelberg, \emph{Equivariant Satake category and Kostant-Whittaker reduction}, Mosc. Math. J. \textbf{8} (2008), no. 1, 39--72, 183.

\bibitem{BS} Bhargav Bhatt and Peter Scholze, \emph{Projectivity of Witt vector affine Grassmannians}, available at http://arxiv.org/abs/1507.06490.

\bibitem{Br} Tom Braden, \emph{Hyperbolic localization of intersection cohomology}, Transform. Groups \textbf{8} (2003), no. 3, 209--216.

\bibitem{BGa} A. Braverman and D. Gaitsgory, \emph{Geometric Eisenstein series}, Invent. Math. \textbf{150} (2002), no. 2, 287--384.

\bibitem{Bry} Ranee Kathryn Brylinski, \emph{Limits of weight spaces, Lusztig's $q$-analogs, and fiberings of adjoint orbits}, J. Amer. Math. Soc. \textbf{2} (1989), no. 3, 517--533.

\bibitem{BuGe} K. Buzzard and T. Gee, \emph{The conjectural connections between automorphic representations and Galois representations}, available at http://arxiv.org/abs/1009.0785.

\bibitem{SGA4.5} Pierre Deligne, \emph{Th\'eor\`emes de finitude en cohomologie $\ell$-adique},  (1977), 233--261.

\bibitem{DM} Pierre Deligne and James S. Milne, \emph{Tannakian categories},  \textbf{900} (1982), 101--228.

\bibitem{DS} V. G. Drinfel'd and Carlos Simpson, \emph{$B$-structures on $G$-bundles and local triviality}, Math. Res. Lett. \textbf{2} (1995), no. 6, 823--829.

\bibitem{Dr} Vladimir Drinfeld, \emph{Infinite-dimensional vector bundles in algebraic geometry: an introduction},  \textbf{244} (2006), 263--304.

\bibitem{Fa} Gerd Faltings, \emph{Algebraic loop groups and moduli spaces of bundles}, J. Eur. Math. Soc. (JEMS) \textbf{5} (2003), no. 1, 41--68.

\bibitem{Fa2} Gerd Faltings, \emph{Theta functions on moduli spaces of $G$-bundles}, J. Algebraic Geom. \textbf{18} (2009), no. 2, 309--369.

\bibitem{FGKV} E. Frenkel, D. Gaitsgory, D. Kazhdan and K. Vilonen, \emph{Geometric realization of Whittaker functions and the Langlands conjecture}, J. Amer. Math. Soc. \textbf{11} (1998), no. 2, 451--484.

\bibitem{FZ} Edward Frenkel and Xinwen Zhu, \emph{Any flat bundle on a punctured disc has an oper structure}, Math. Res. Lett. \textbf{17} (2010), no. 1, 27--37.

\bibitem{Ga} D. Gaitsgory, \emph{Construction of central elements in the affine Hecke algebra via nearby cycles}, Invent. Math. \textbf{144} (2001), no. 2, 253--280.

\bibitem{Ga2} D. Gaitsgory, \emph{On de Jong's conjecture}, Israel J. Math. \textbf{157} (2007), 155--191.

\bibitem{Ga3} Dennis Gaitsgory, \emph{Contractibility of the space of rational maps}, Invent. Math. \textbf{191} (2013), no. 1, 91--196.

\bibitem{GL} Dennis Gaitsgory and Jacob Lurie, \emph{Weil's conjecture for function fields}, available at http://math.harvard.edu/~lurie/papers/tamagawa.pdf.

\bibitem{Gi1} Victor Ginzburg, \emph{Perverse sheaves on a loop group and Langlands' duality}, available at http://arxiv.org/abs/alg-geom/9511007.

\bibitem{GKM} Mark Goresky, Robert Kottwitz and Robert MacPherson, \emph{Equivariant cohomology, Koszul duality, and the localization theorem}, Invent. Math. \textbf{131} (1998), no. 1, 25--83.

\bibitem{Gortz} Ulrich G{\"o}rtz, \emph{Affine Springer fibers and affine Deligne-Lusztig varieties},  (2010), 1--50.

\bibitem{Gr} Benedict H. Gross, \emph{On the Satake isomorphism},  \textbf{254} (1998), 223--237.

\bibitem{Ha} Thomas J. Haines, \emph{Equidimensionality of convolution morphisms and applications to saturation problems}, Adv. Math. \textbf{207} (2006), no. 1, 297--327.

\bibitem{He} Jochen Heinloth, \emph{Uniformization of $\scr G$-bundles}, Math. Ann. \textbf{347} (2010), no. 3, 499--528.

\bibitem{Il2} Luc Illusie, \emph{Autour du th\'eor\`eme de monodromie locale}, Ast\'erisque (1994), no. 223, 9--57.

\bibitem{Il3} Luc Illusie, \emph{Around the Thom-Sebastiani theorem}, available at http://www.math.u-psud.fr/~illusie/Illusie-nearby-cycles27d.pdf.

\bibitem{Kac} Victor G. Kac, \emph{Infinite-dimensional Lie algebras}, 3 ed., Cambridge University Press, Cambridge, 1990.

\bibitem{KWWY} Joel Kamnitzer, Ben Webster, Alex Weekes and Oded Yacobi, \emph{Yangians and quantizations of slices in the affine Grassmannian}, Algebra Number Theory \textbf{8} (2014), no. 4, 857--893.

\bibitem{KL} David Kazhdan and George Lusztig, \emph{Schubert varieties and Poincar\'e duality},  (1979), 185--203.

\bibitem{Ku} Shrawan Kumar, \emph{Demazure character formula in arbitrary Kac-Moody setting}, Invent. Math. \textbf{89} (1987), no. 2, 395--423.

\bibitem{KNR} Shrawan Kumar, M. S. Narasimhan and A. Ramanathan, \emph{Infinite Grassmannians and moduli spaces of $G$-bundles}, Math. Ann. \textbf{300} (1994), no. 1, 41--75.

\bibitem{KN} Shrawan Kumar and M. S. Narasimhan, \emph{Picard group of the moduli spaces of $G$-bundles}, Math. Ann. \textbf{308} (1997), no. 1, 155--173.

\bibitem{Laf} Vincent Lafforgue, \emph{Chtoucas pour les groupes r\'{e}ductifs et param\'{e}trisation de Langlands globale}, available at http://arxiv.org/abs/1209.5352.

\bibitem{LS} Yves Laszlo and Christoph Sorger, \emph{The line bundles on the moduli of parabolic $G$-bundles over curves and their sections}, Ann. Sci. \'Ecole Norm. Sup. (4) \textbf{30} (1997), no. 4, 499--525.

\bibitem{LZ} Yifeng Liu and Weizhe Zheng, \emph{Enhanced six operations and base change theorem for Artin stacks}, available at http://arxiv.org/abs/1211.5948.

\bibitem{Lu0} G. Lusztig, \emph{Green polynomials and singularities of unipotent classes}, Adv. in Math. \textbf{42} (1981), no. 2, 169--178.

\bibitem{Lu} George Lusztig, \emph{Singularities, character formulas, and a $q$-analog of weight multiplicities},  \textbf{101} (1981), 208--229.

\bibitem{Ma} Olivier Mathieu, \emph{Formules de caract\`eres pour les alg\`ebres de Kac-Moody g\'en\'erales}, Ast\'erisque (1988), no. 159-160, 267.

\bibitem{MS} James S. Milne and Kuang-yen Shih, \emph{Conjugates of Shimura varieties},  \textbf{900} (1982), 280--356.

\bibitem{MV} I. Mirkovi{\'c} and K. Vilonen, \emph{Geometric Langlands duality and representations of algebraic groups over commutative rings}, Ann. of Math. (2) \textbf{166} (2007), no. 1, 95--143.

\bibitem{MVy} I. Mirkovi\'c and M. Vybornov, \emph{Quiver varieties and Beilinson-Drinfeld Grassmannians of type $A$},

\bibitem{Mum} David Mumford, \emph{Abelian varieties}, Published for the Tata Institute of Fundamental Research, Bombay; Oxford University Press, London, 1970.

\bibitem{Na} David Nadler, \emph{Matsuki correspondence for the affine Grassmannian}, Duke Math. J. \textbf{124} (2004), no. 3, 421--457.

\bibitem{NP} B. C. Ng{\^o} and P. Polo, \emph{R\'esolutions de Demazure affines et formule de Casselman-Shalika g\'eom\'etrique}, J. Algebraic Geom. \textbf{10} (2001), no. 3, 515--547.

\bibitem{OZ1} Denis Osipov and Xinwen Zhu, \emph{A categorical proof of the Parshin reciprocity laws on algebraic surfaces}, Algebra Number Theory \textbf{5} (2011), no. 3, 289--337.

\bibitem{OZ2} Denis Osipov and Xinwen Zhu, \emph{The two-dimensional Contou-Carr\`ere symbol and reciprocity laws}, available at http://arxiv.org/abs/1305.6032.

\bibitem{PR} G. Pappas and M. Rapoport, \emph{Twisted loop groups and their affine flag varieties}, Adv. Math. \textbf{219} (2008), no. 1, 118--198.

\bibitem{PR2} G. Pappas and M. Rapoport, \emph{$\Phi$-modules and coefficient spaces}, Mosc. Math. J. \textbf{9} (2009), no. 3, 625--663, back matter.

\bibitem{PZ} G. Pappas and X. Zhu, \emph{Local models of Shimura varieties and a conjecture of Kottwitz}, Invent. Math. \textbf{194} (2013), no. 1, 147--254.

\bibitem{PS} Andrew Pressley and Graeme Segal, \emph{Loop groups}, The Clarendon Press, Oxford University Press, New York, 1986.

\bibitem{R3} Michael Rapoport, \emph{A positivity property of the Satake isomorphism}, Manuscripta Math. \textbf{101} (2000), no. 2, 153--166.

\bibitem{Ra} Michael Rapoport, \emph{A guide to the reduction modulo $p$ of Shimura varieties}, Ast\'erisque (2005), no. 298, 271--318.

\bibitem{Rei} Ryan Cohen Reich, \emph{Twisted geometric Satake equivalence via gerbes on the factorizable Grassmannian}, Represent. Theory \textbf{16} (2012), 345--449.

\bibitem{RZ} Timo Richarz and Xinwen Zhu, \emph{Construction of the full Langlands dual group via the geometric Satake correspondence},

\bibitem{Ri} Timo Richarz, \emph{A new approach to the geometric Satake equivalence}, Doc. Math. \textbf{19} (2014), 209--246.

\bibitem{Sch} Peter Scholze and Jared Weinstein, \emph{$p$-adic geometry},

\bibitem{Se2} Jean-Pierre Serre, \emph{Local fields}, Springer-Verlag, New York-Berlin, 1979.

\bibitem{So0} Christoph Sorger, \emph{On moduli of $G$-bundles of a curve for exceptional $G$}, Ann. Sci. \'Ecole Norm. Sup. (4) \textbf{32} (1999), no. 1, 127--133.

\bibitem{So} Christoph Sorger, \emph{Lectures on moduli of principal $G$-bundles over algebraic curves},  \textbf{1} (1999), 1--57.

\bibitem{Sp} T. A. Springer, \emph{Reductive groups},  (1977), 3--27.

\bibitem{St} Stack project authors, \emph{Stack project},  (2016).

\bibitem{Ti} J. Tits, \emph{Reductive groups over local fields},  (1977), 29--69.

\bibitem{Ti3} Jacques Tits, \emph{Groupes associ\'es aux alg\`ebres de Kac-Moody}, Ast\'erisque (1989), no. 177-178, Exp.\ No.\ 700, 7--31.

\bibitem{YZ} Zhiwei Yun and Xinwen Zhu, \emph{Integral homology of loop groups via Langlands dual groups}, Represent. Theory \textbf{15} (2011), 347--369.

\bibitem{Yu} Zhiwei Yun, \emph{Lectures on Springer theories and orbital integrals},

\bibitem{Z07} Xinwen Zhu, \emph{Affine Demazure modules and $T$-fixed point subschemes in the affine Grassmannian}, Adv. Math. \textbf{221} (2009), no. 2, 570--600.

\bibitem{Z10} Xinwen Zhu, \emph{On the coherence conjecture of Pappas and Rapoport}, Ann. of Math. (2) \textbf{180} (2014), no. 1, 1--85.

\bibitem{Z11} Xinwen Zhu, \emph{The geometric Satake correspondence for ramified groups}, Ann. Sci. \'Ec. Norm. Sup\'er. (4) \textbf{48} (2015), no. 2, 409--451.

\bibitem{Z14} Xinwen Zhu, \emph{Affine Grassmannians and the geometric Satake in mixed characteristic}, available at http://arxiv.org/abs/1407.8519.

\end{thebibliography}


\end{document}
